%==========================%
% My article template for research paper & lecture note
% For lecture note, remove H
% For research article, remove C,D and I
% stored in Evernote, May 2018
%==========================%
%\documentclass[11pt]{elsarticle}
\documentclass[11pt,authoryear]{elsarticle}
\usepackage[margin=2.3cm, includehead]{geometry}
\makeatletter
\def\ps@pprintTitle{%
  \let\@oddhead\@empty
  \let\@evenhead\@empty
  \let\@oddfoot\@empty
  \let\@evenfoot\@oddfoot
}
\makeatother
%==========================%
% B. Common pkg for all articles
%==========================%
\usepackage{currfile}
\usepackage{mathrsfs}
\usepackage[mathscr]{eucal}
\DeclareMathAlphabet{\mathpzc}{OT1}{pzc}{m}{it}
\usepackage{natbib}
\usepackage[notcite, notref]{showkeys} % show labels when printed
\usepackage{makeidx}
\usepackage{multirow}
\usepackage{multicol}
\usepackage{booktabs}
\usepackage{enumerate}
\usepackage{setspace}
\usepackage{float}
\usepackage{hhline}
\usepackage{bbm}
\usepackage[dvipsnames,svgnames,table]{xcolor}
\usepackage{graphicx}
\usepackage{epstopdf}
\usepackage[normalem]{ulem}
\usepackage{listings}
\usepackage{kotex}
%\usepackage[bookmarks]{hyperref} % hyper link for cross referencing
\usepackage{amsmath}
\usepackage{amsfonts}
\usepackage{url}
\usepackage{amssymb}
%\usepackage{graphics}
\usepackage{dsfont}
\usepackage{algorithm}
\usepackage{algpseudocode}
\usepackage{epsfig}
\usepackage{verbatim}
\usepackage{amsthm}
% new LaTeX commands
\geometry{left=1in,right=1in,top=1in,bottom=1in}
\usepackage{framed}    % to make framed boxes
\usepackage{scalefnt}
\usepackage{alltt}
\usepackage{tikz}
\usetikzlibrary{arrows.meta}
\usepackage{caption}
\usepackage{hyperref}
\hypersetup{
    bookmarks=true,         % show bookmarks bar?
    unicode=false,          % non-Latin characters in Acrobat's bookmarks
    pdftoolbar=true,        % show Acrobat's toolbar?
    pdfmenubar=true,        % show Acrobat's menu?
    pdffitwindow=true,      % page fit to window when opened
    pdftitle={My title},    % title
    pdfauthor={Author},     % author
    pdfsubject={Subject},   % subject of the document
    pdfnewwindow=true,      % links in new window
    pdfkeywords={keywords}, % list of keywords
    colorlinks=true,        % false: boxed links; true: colored links
    linkcolor=blue,         % color of internal links
    citecolor=blue,        % color of links to bibliography
    filecolor=magenta,      % color of file links
    urlcolor=cyan           % color of external links
}

\onehalfspacing
\def\iid{\buildrel {\rm i.i.d.} \over \sim}
\def\wc{\buildrel {\rm (d)} \over \longrightarrow}
\def\eqd{\buildrel {\rm (d)} \over =}
\def\eqas{\buildrel {\rm a.s.} \over =}


\def\cp{\buildrel {\rm p} \over \longrightarrow}
\def\cw#1 { \overset{\mathbb{P}}{\underset{#1}{\longrightarrow}} }
\def\Real{\mathbb{R}}
\def\Natu0{\mathbb{N}_0}
\def\P#1{{\mathbb{P}}\left(#1\right)}
\def\Pn#1#2{{\mathrm{P}_#1}\left(#2\right)}

\def\R{{\mathbb R}}  %%


\def\bTheta{\boldsymbol{\Theta}}
\def\br{\boldsymbol{r}}
\def\bY{\boldsymbol{Y}}
\def\bZ{\boldsymbol{Z}}
\def\bmu{\boldsymbol{\mu}}


\def\E#1{{\mathbb E}\left[#1\right]}
\def\Ew#1#2{{\mathbb E}_{#1}\left[#2\right]}


\def\Var#1{{\rm Var}\left(#1\right)}
\def \I#1{\mathbb{I}\left\{#1\right\}}
\def \hix#1{ {\rm HIX}\left(#1\right) }
\def \six#1{ {\rm SIX}\left(#1\right) }
\def \kix#1{ {\rm KIX}\left(#1\right) }

\def \rhix#1{ {{\rm RHIX}}\left(#1\right) }

\def \rcov#1#2 {{\rm cov}_{#1}\left( #2\right)}
\def \corr#1{ {{\rm corr}}\left(#1\right) }
\def\cov#1{{\rm  Cov}\left(#1\right)}
\def\cor#1{{\rm  corr}\left[#1\right]}


\newcommand{\Es}[1]{\operatorname{\mathbb E}_{s}\!\left[#1\right]}
\newcommand{\kappaT}{\kappa_{k}}          % cumulant sub‑k
\newcommand{\muST}{\mu_{s,\tau}}          % tilted mean shorthand

\DeclareMathOperator{\esssup}{ess\,sup}
\DeclareMathOperator{\essinf}{ess\,inf}
\DeclareMathOperator*{\argmin}{arg\,min}
\DeclareMathOperator*{\argmax}{arg\,max}
\newtheorem{notation}{Notation}
\newtheorem{example}{Example}
\renewcommand{\emph}[1]{\textbf{#1}}

\DeclareMathOperator{\sgn}{{\rm sgn}}
\oddsidemargin0in
\textwidth6.5in
\addtolength{\topmargin}{-.75in}
\textheight 8.5in
\newtheorem{lemma}{Lemma}
\newtheorem{theorem}{Theorem}
\newtheorem{definition}{Definition}
\newtheorem{corollary}{Corollary}
\newtheorem{remark}{Remark}
\newtheorem{proposition}{Proposition}
\newtheorem{assumption}{Assumption}
\newtheorem{example0}{Example}
%\newtheorem*{toy*}{Toy Model}
\newtheorem{claim}{Claim}
\newtheorem{model}{Model}
\newtheorem{aux}{Auxiliary Model}
\newtheorem{question}{Question}

%\newtheorem{structure}{Structure}
%\newtheorem{regular}{Regularity Conditions}
\makeatletter
\newtheoremstyle{nopunct}% name
  {}{}              % Space above/below
  {\itshape}        % Body font
  {}                % Indent amount
  {\bfseries}       % Theorem head font
  {}                % Punctuation after theorem head
  { }               % Space after theorem head
  {}                % Theorem head spec
\makeatother

\theoremstyle{nopunct}
\newtheorem*{intermedmodel}{}
\numberwithin{equation}{section}




\newcommand{\bl}[1]{\textcolor{blue}{{#1}}}
\newcommand{\red}[1]{\textcolor{red}{{#1}}}
\newcommand{\mg}[1]{\textcolor{magenta}{{#1}}}


\begin{document}
	\begin{titlepage}	
		\thispagestyle{empty}
%\begin{frontmatter}
\title{Hurdle-Lognormal state-space model}

\author{Jae Youn Ahn\fnref{thirdfoot}}%\corref{bbb}}
%\author{\qquad Qihe Tang \fnref{fifthfoot}}




%\cortext[bbb]{Corresponding authors/equal contribution}
\fntext[thirdfoot]{Department of Statistics, Ewha Womans University, Seoul, Republic of Korea. Email: \url{jaeyahn@ewha.ac.kr}}
%\fntext[fifthfoot]{School of Risk and Actuarial Studies, UNSW, Sydney, Australia. Email: \url{qihe.tang@unsw.edu.au}}




\begin{abstract}

To be written.

\end{abstract}


\maketitle


\textbf{Keywords:} 


\vspace{.2cm}

JEL Classification: C3, C5


\vspace{.2cm}


Version of \today
\end{titlepage}

\newpage



\section{Introduction}

Many insurance payment variables have a mixed distribution consisting
of a point mass at zero and a continuous distribution over the positive
real line. A zero observation may represent the absence of an insured
event, a claim, or a payment during a given period, whereas a positive
observation records the corresponding payment amount. Such data can be
modelled through several approaches, including collective risk models,
Tweedie distributions, and two-part models. In a two-part model, the
loss at time \(t\) is represented as
\[
Y_t=Z_tS_t,
\]
where the binary variable \(Z_t\) indicates whether a positive loss
occurs and the positive continuous variable \(S_t\) represents its
amount. This decomposition makes it possible to model the occurrence
and positive amount separately while retaining a direct interpretation
of their combined effect on the total loss.

For longitudinal ratemaking, it is important to account for dependence
among repeated losses from the same policyholder or risk. Such
dependence is commonly introduced through random effects representing
unobserved heterogeneity. In a two-part model, a shared random effect may
simultaneously influence the occurrence and positive amount, or separate
but correlated random effects may be assigned to the two components.
The latter formulation is more flexible because it allows the
unobserved factors affecting occurrence and severity to differ, while
their association is captured through the correlation between the
random effects. This flexibility, however, generally makes likelihood,
posterior, and predictive calculations more difficult.

The distinction between tractability for estimation and tractability
for prediction is particularly important in insurance. For parameter
estimation, one seeks an analytically manageable likelihood function and
posterior distribution of the random effects. For experience ratemaking,
the principal quantity is instead the conditional expectation
\[
\E{Y_{t+1}\mid Y_{1:t}=y_{1:t}},
\]
which provides the one-step-ahead a posteriori premium under the
equivalence principle. For excess-of-loss coverage, the related
conditional stop-loss expectation
\[
\E{(Y_{t+1}-d)_+\mid Y_{1:t}=y_{1:t}},
\qquad d>0,
\]
is also of direct actuarial interest. A model may admit a tractable
likelihood without providing corresponding analytical expressions for
these predictive quantities. In many random-effects and copula-based
two-part models, they must therefore be evaluated through additional
numerical integration or simulation.

A particularly useful analytical structure arises from combining a
Gaussian random effect with a probit link. For dependent binary
observations, the resulting Bernoulli-probit random effect model
expresses joint probabilities, posterior densities, and predictive
probabilities through multivariate Gaussian cumulative distribution
functions
\citep{durante2019conjugate,fasano2021closed,cao2022scalable}.
Two-part random-effects and stochastic-process models for longitudinal
semicontinuous data have also been studied extensively in the
statistical and biostatistical literature
\citep{farewell2017two,yiu2018two}. In the actuarial literature, joint
modelling of dependent binary insurance outcomes is likewise well
established
\citep{young2009multivariate,frees2010dependent,guo2025entity,
oh2026matrix}. Nevertheless, the analytical advantages of the
Gaussian-probit construction appear to have received limited attention
for insurance two-part models, particularly for deriving conditional
premium formulas under correlated occurrence and severity random
effects.

Motivated by this gap, this paper studies a Bernoulli-lognormal probit
two-part model with a bivariate Gaussian random effect. The first random
effect enters the probit model for the occurrence indicator, whereas the
second enters the lognormal model for the positive loss amount. Their
correlation links the occurrence and severity components and allows
positive-loss observations to update the latent risk associated with
future claim occurrence. Thus, the model retains the interpretability
and flexibility of the standard bivariate random-effects formulation
while preserving a tractable Gaussian structure.

The main contribution of the paper is to show that this structure
provides a fully analytical environment for both ratemaking and
statistical inference. First, we derive a closed-form expression for
\[
\E{Y_{t+1}\mid Y_{1:t}=y_{1:t}}
\]
in terms of a ratio of multivariate Gaussian cumulative distribution
functions. The derivation combines Gaussian updating based on the
observed positive losses, Gaussian exponential tilting, and a
Gaussian product-probit identity. Second, we derive a closed-form
expression for the conditional stop-loss expectation
\[
\E{(Y_{t+1}-d)_+\mid Y_{1:t}=y_{1:t}},
\]
which provides an analytical one-step-ahead premium for the layer above
a deductible \(d\). These formulas avoid direct numerical integration
over the bivariate random effect.

For statistical inference, we further derive the mixed
discrete-continuous likelihood of the observed loss history and the
posterior density of the bivariate random effect. Consequently, the same
Gaussian representation that produces the predictive premiums also
provides tractable likelihood-based and Bayesian inference. To establish
the main analytical mechanism, we first review the Bernoulli-probit
random effect model and its closed-form joint probabilities, predictive
probabilities, and posterior density. We then introduce the
Bernoulli-lognormal bivariate random effect model and derive its
one-step-ahead mean and stop-loss predictions, followed by the
corresponding likelihood and posterior results.

\section{Survival representation of counting distributions}
\label{sec.3}
Hurdle representations provides the flexibility in modelling the count data by decomposing a
count distribution according to whether the count crosses a specified
threshold. The standard hurdle representation uses the threshold,
separating \(Y=0\) from \(Y\ge1\). It can be extended to a double-hurdle representation, separating \(Y=1\) from \(Y\ge2\) conditional on
having already crossed the first threshold. Motivated by these models, this section presents a more
general threshold-by-threshold view of count distributions, in which the count
is described by its successive passages from one level to the next. 

We begin with the notation used throughout the paper. Let \(\mathbb N\) denote
the set of positive integers and let
\[
\mathbb N_0:=\mathbb N\cup\{0\}.
\]
We write \(\Real\), \(\Real_{\ge0}\), and \(\Real_{>0}\) for the real line,
the nonnegative real line, and the positive real line, respectively. For \(k\in\mathbb N\) and for any set \(A\), we write
\(
A^k
\)
for the \(k\)-fold Cartesian product of \(A\). In particular,
\(\mathbb N^k\), \(\mathbb N_0^k\), \(\Real^k\),
\(\Real_{\ge0}^k\), and \(\Real_{>0}^k\) denote the corresponding
\(k\)-dimensional product spaces. 

For \(k\in\mathbb N\), and for real numbers \(x_1,\ldots,x_r\), we use 
\[
x_{1:k}:=(x_1,\ldots,x_k)^\top\in\Real^k
\]
and all the vector is regarded as column vector. 
Especially, we use $0_{1:k}=(0, \cdots, 0)^\top\in\Real^k$, 
and similarly, 
in a matrix environment, we write
\(
0_{i\times j}
\) to be $i\times j$ matrix of zero. 
We also write
\[
e_{1:k} = (0, \cdots, 0, 1)^\top\in\Real^k.
\]
If \(x_{1:k},y_{1:k}\in\Real^k\), their Euclidean inner product is denoted by
\[
\langle x_{1:k},y_{1:k}\rangle
:=
\sum_{j=1}^{k}x_jy_j.
\]
Throughout, we adopt the standard conventions that a product over an
empty index set equals \(1\), whereas a sum over an empty index set
equals \(0\). That is, if \(\mathcal I=\varnothing\), then
\[
\prod_{j\in\mathcal I} a_j:=1,
\qquad
\sum_{j\in\mathcal I} a_j:=0.
\]


We then introduce some distributions used in this paper.
\begin{itemize}
  \setlength\itemsep{0em}
	\item ${\rm Gamma}(a, b)$: gamma distribution with shape parameter $a\in\Real_{>0}$  and
	rate  parameter$b\in\Real_{>0}$. 
	\item ${\rm Beta}(a, b)$: beta distribution on $(0,1)$ with parameters $a\in\Real_{>0}$ and $b\in\Real_{>0}$.
	\item ${\rm Pois}( \lambda)$: Poisson distribution
 with mean $\lambda\in\Real_{>0}$.
\item ${\rm Geo}(p)$: geometric distribution on \(\mathbb N_0\) with parameter
\(p\in(0,1)\) so that its mean is \(p/(1-p)\).
	 mean $\lambda\in\Real_{>0}$ and variance  $\lambda + \lambda^2/\Gamma$ 
\item ${\rm N}(\mu, b^2)$: normal distribution with mean $\mu\in\Real_{>0}$ and variance $\sigma^2\in\Real_{>0}$.
\item ${\rm MVN}_k(\mu_{1:k}, \boldsymbol\Sigma)$: multivariate normal distribution with \(\mu_{1:k}\in\Real^k\) and
a positive definite covariance matrix \(\boldsymbol\Sigma\in\Real^{k\times k}\). 


\end{itemize}

For a positive semidefinite covariance matrix
\(\boldsymbol\Sigma\in\Real^{k\times k}\), let
\[
\Phi_k(x_{1:k}\,;\,\boldsymbol\Sigma),
\qquad
x_{1:k}\in\Real^k,
\]
denote the cumulative distribution function at \(x_{1:k}\) of
\[
{\rm MVN}_k(0_{1:k},\boldsymbol\Sigma).
\]
That is, if
\[
X_{1:k}
\sim
{\rm MVN}_k(0_{1:k},\boldsymbol\Sigma),
\]
then
\[
\Phi_k(x_{1:k}\,;\,\boldsymbol\Sigma)
:=
\P{
X_1\le x_1,\ldots,X_k\le x_k
}.
\]
If \(\boldsymbol\Sigma\) is positive definite, let
\[
\phi_k(x_{1:k}\,;\,\boldsymbol\Sigma),
\qquad
x_{1:k}\in\Real^k,
\]
denote the corresponding probability density function with respect to
the \(k\)-dimensional Lebesgue measure. 

When the covariance argument is omitted, we use the standard
multivariate normal distribution:
\[
\phi_k(x_{1:k})
:=
\phi_k(x_{1:k}\,;\,\boldsymbol I_k),
\qquad
\Phi_k(x_{1:k})
:=
\Phi_k(x_{1:k}\,;\,\boldsymbol I_k),
\]
where \(\boldsymbol I_k\) is the \(k\times k\) identity matrix. Thus,
\(\phi_k(x_{1:k})\) and \(\Phi_k(x_{1:k})\) are, respectively, the
probability density function and cumulative distribution function of
\[
{\rm MVN}_k(0_{1:k},\boldsymbol I_k).
\]
When \(k=1\), for \(x\in\Real\), we simply write
\(\phi(x)\) and \(\Phi(x)\), omitting the subscript \(k=1\).

















\section{Review}
Here we provide the useful auxiliary result for this paper.

The following is the well-known result of Gaussian exponential tilting.
\begin{lemma}
\label{lem.4}
Let
\[
X\sim N(\mu,v),
\]
where \(\mu\in\Real\) and \(v>0\). For \(h\in\Real\), let
\[
\widetilde X
\sim
N(\mu+vh,v).
\]
Then, for every Borel-measurable function
\(g:\Real\to\Real\) satisfying
\[
\E{
\left|
\exp(hX)g(X)
\right|
}
<\infty,
\]
we have
\[
\E{
\exp(hX)g(X)
}
=
\exp\left(
h\mu+\frac{h^2v}{2}
\right)
\E{
g(\widetilde X)
}.
\]
\end{lemma}

\begin{proof}

The integrability assumption gives
\[
\begin{aligned}
\infty
&>
\E{
\left|
\exp(hX)g(X)
\right|
}\\
&=
\int_{\Real}
\exp(hx)|g(x)|
\frac{1}{\sqrt{2\pi v}}
\exp\left(
-\frac{(x-\mu)^2}{2v}
\right)
\,{\rm d}x\\
&=
\exp\left(
h\mu+\frac{h^2v}{2}
\right)
\int_{\Real}
|g(x)|
\frac{1}{\sqrt{2\pi v}}
\exp\left(
-\frac{(x-(\mu+vh))^2}{2v}
\right)
\,{\rm d}x\\
&=
\exp\left(
h\mu+\frac{h^2v}{2}
\right)
\E{
\left\vert g(\widetilde X)\right\vert
}
\end{aligned}
\]
which ensures
\[
\E{\left\vert g(\widetilde X) \right\vert }<\infty.
\]
Consequently, the same calculation with \(|g|\) replaced by \(g\) is
justified and gives
\[
\E{
\exp(hX)g(X)
}
=
\exp\left(
h\mu+\frac{h^2v}{2}
\right)
\E{
g(\widetilde X)
},
\]
which proves the result.
\end{proof}


The following is the multivariate version of the Gaussian exponential
tilting identity.

\begin{lemma}[Multivariate Gaussian exponential tilting]
\label{lem.3}
Let
\[
X_{1:t}
\sim
\operatorname{MVN}_t(\mu_{1:t},\boldsymbol\Sigma),
\]
where \(\mu_{1:t}\in\Real^t\) and
\(\boldsymbol\Sigma\in\Real^{t\times t}\) is positive definite covariance matrix. For
\(h_{1:t}\in\Real^t\), let
\[
\widetilde X_{1:t}
\sim
\operatorname{MVN}_t
\left(
\mu_{1:t}+\boldsymbol\Sigma h_{1:t},
\boldsymbol\Sigma
\right).
\]
Then, for every Borel-measurable function
\(g:\Real^t\to\Real\) satisfying
\[
\E{
\left|
\exp\left(h_{1:t}^{\top}X_{1:t}\right)
g(X_{1:t})
\right|
}
<\infty,
\]
we have
\[
\begin{aligned}
&\E{
\exp\left(h_{1:t}^{\top}X_{1:t}\right)
g(X_{1:t})
}\\
&\quad =
\exp\left(
h_{1:t}^{\top}\mu_{1:t}
+
\frac12
h_{1:t}^{\top}\boldsymbol\Sigma h_{1:t}
\right)
\E{
g(\widetilde X_{1:t})
}.
\end{aligned}
\]
\end{lemma}

\begin{proof}
The integrability assumption gives
\[
\begin{aligned}
&\infty
>
\E{
\left|
\exp\left(h_{1:t}^{\top}X_{1:t}\right)
g(X_{1:t})
\right|
}\\
&=
\int_{\Real^t}
\exp\left(h_{1:t}^{\top}x_{1:t}\right)
|g(x_{1:t})|
\frac{1}
{(2\pi)^{t/2}|\boldsymbol\Sigma|^{1/2}}
\exp\left(
-\frac12
(x_{1:t}-\mu_{1:t})^\top
\boldsymbol\Sigma^{-1}
(x_{1:t}-\mu_{1:t})
\right)
\,{\rm d}x_{1:t}\\
&=
\int_{\Real^t}
|g(x_{1:t})|
\frac{1}
{(2\pi)^{t/2}|\boldsymbol\Sigma|^{1/2}}\\
&\qquad\times
\exp\left(
h_{1:t}^{\top}\mu_{1:t}
+
\frac12
h_{1:t}^{\top}\boldsymbol\Sigma h_{1:t}
-\frac12
\left(
x_{1:t}-\mu_{1:t}-\boldsymbol\Sigma h_{1:t}
\right)^\top
\boldsymbol\Sigma^{-1}
\left(
x_{1:t}-\mu_{1:t}-\boldsymbol\Sigma h_{1:t}
\right)
\right)
\,{\rm d}x_{1:t}\\
&=
\exp\left(
h_{1:t}^{\top}\mu_{1:t}
+
\frac12
h_{1:t}^{\top}\boldsymbol\Sigma h_{1:t}
\right)\\
&\quad\times
\int_{\Real^t}
|g(x_{1:t})|
\frac{1}
{(2\pi)^{t/2}|\boldsymbol\Sigma|^{1/2}}
\exp\left(
-\frac12
\left(
x_{1:t}-\mu_{1:t}-\boldsymbol\Sigma h_{1:t}
\right)^\top
\boldsymbol\Sigma^{-1}
\left(
x_{1:t}-\mu_{1:t}-\boldsymbol\Sigma h_{1:t}
\right)
\right)
\,{\rm d}x_{1:t}\\
&=
\exp\left(
h_{1:t}^{\top}\mu_{1:t}
+
\frac12
h_{1:t}^{\top}\boldsymbol\Sigma h_{1:t}
\right)
\E{
\left|g(\widetilde X_{1:t})\right|
}
\end{aligned}
\]
which ensures
\[
\E{
\left|g(\widetilde X_{1:t})\right|
}
<\infty.
\]
Consequently, the same calculation with \(|g|\) replaced by \(g\) is
justified and gives
\[
\begin{aligned}
&\E{
\exp\left(h_{1:t}^{\top}X_{1:t}\right)
g(X_{1:t})
}\\
&\quad =
\exp\left(
h_{1:t}^{\top}\mu_{1:t}
+
\frac12
h_{1:t}^{\top}\boldsymbol\Sigma h_{1:t}
\right)
\E{
g(\widetilde X_{1:t})
},
\end{aligned}
\]
which proves the result.
\end{proof}


\begin{lemma}[Gaussian product-probit identity]
\label{lem.5}
Let
\[
X_{1:t}=(X_1,\ldots,X_t)^\top
\]
be a Gaussian vector with mean
\(\mu_{1:t}\in\Real^t\) and covariance matrix
\(\boldsymbol\Sigma\in\Real^{t\times t}\). Then
\[
\E{\Phi_{t}(X_{1:t})}
=
\E{
\prod_{j=1}^t
\Phi(X_j)
}
=
\Phi_t
\left(
\mu_{1:t};
\boldsymbol I_t+\boldsymbol\Sigma
\right).
\]
\end{lemma}

\begin{proof}
Let
\[
W_{1:t}:=(W_1,\ldots,W_t)^\top
\sim
\operatorname{MVN}(0_{1:t},\boldsymbol I_t)
\]
be independent of \(X_{1:t}\). Then, we have
\[
\begin{aligned}
%\E{\Phi_{t}(X_{1:t})}
%&=
%\E{
\Phi_{t}(X_{1:t})&=
\P{W_{1:t}\le X_{1:t}\mid X_{1:t}}\\
%}\\
%&=
%\E{\P{W_{1:t}\le X_{1:t}\mid X_{1:t}}}\\
&=\prod_{j=1}^{t}\P{W_j\le X_j\mid X_j}\\
&=\prod_{j=1}^{t}\Phi(X_j)
\end{aligned}
\]
which further implies
\[
\begin{aligned}
\E{\Phi_{t}(X_{1:t})}=\E{\E{\prod_{j=1}^{t}\Phi(X_j)}}.
\end{aligned}
\]

%Since \(W_1,\ldots,W_t\) are independent,
Furthermore, from the following observations
\[
\begin{aligned}
\E{\Phi_{t}(X_{1:t})}&=
%\E{
%\prod_{j=1}^t
%\Phi(X_j)
%}\\
%&=
%\E{
%\P{
%W_{1:t}\le X_{1:t}
%\mid
%X_{1:t}
%}
%}\\
%&=
%\P{
%W_{1:t}\le X_{1:t}
%}\\
&=
\P{
W_{1:t}-(X_{1:t}-\mu_{1:t})
\le
\mu_{1:t}
}
\end{aligned}
\]
and
\[
W_{1:t}-(X_{1:t}-\mu_{1:t})
\sim
\operatorname{MVN}
\left(
0_{1:t},
\boldsymbol I_t+\boldsymbol\Sigma
\right),
\]
we conclude
\[
\P{
W_{1:t}-(X_{1:t}-\mu_{1:t})
\le
\mu_{1:t}
}
=
\Phi_t
\left(
\mu_{1:t};
\boldsymbol I_t+\boldsymbol\Sigma
\right),
\]
which proves the result.
\end{proof}


\subsection{Review: Bernoulli-probit random effect model}

Joint modelling of dependent binary observations is an important topic
in the statistical literature. Among the available approaches, models
combining a Gaussian random effect with a probit link have recently attracted
considerable attention because they flexibly accommodate dependence
among repeated binary observations while retaining analytically
tractable Gaussian representations
\citep{farewell2017two,durante2019conjugate,fasano2021closed,
cao2022scalable}.

Although joint modelling of binary observations is also well established
in the actuarial and insurance literature
\citep{young2009multivariate,frees2010dependent,guo2025entity,
oh2026matrix}, the analytical tractability offered specifically by the
Bernoulli-probit random effect model appears to have received limited
attention in actuarial applications. We therefore briefly review its
 random-effect formulation and derive closed-form expressions for
the joint probability, one-step prediction, and posterior distribution
of the random effect. Although these results are known in the statistical
literature; see, for example, \citet{farewell2017two}, we include them for
completeness and as a foundation for the Bernoulli-lognormal models
developed in the subsequent sections.

\begin{model}[Bernoulli-probit AR(1) state-space model]
\label{model.1}
Let
\[
\sigma\in\Real_{>0}\quad\hbox{and}\quad \left(c_t\right)_{t\in\mathbb N}\subseteq\Real
\]
be predetermined exogenous constants. We define the  sequence
\((Z_t)_{t\in\mathbb N}\) through the following random effect assumption and observation
equation.

\begin{itemize}
  \item[i.] Random effect: Assume
  \[
  \Theta\sim \operatorname{N}(0, \sigma^2).
  \]

  \item[ii.] Observation equation:
  Conditional on the random effect $\Theta$, assume that
   \(Z_1,Z_2,\ldots\) are mutually independent and
  \[
  Z_t\mid \Theta
  \sim
  \operatorname{Ber}\left(\Phi(c_t+\Theta)\right),
  \qquad t=1,2,\ldots.
  \]
\end{itemize}
\end{model}


The following theorem provides closed-form expressions for the joint
probability of the observed binary history and the corresponding
one-step-ahead prediction under Model~\ref{model.1}. In particular, the
joint probability, and hence the likelihood, is expressed in terms of a
multivariate normal cumulative distribution function, whereas the
one-step-ahead prediction is obtained as the ratio of the joint
probabilities for the augmented and observed histories.


\begin{theorem}[Closed-form one-step prediction]
\label{thm.1}
Consider Model~\ref{model.1}. For a realized history
\[
Z_{1:t}=z_{1:t},
\qquad\hbox{for}\quad
z_1, \cdots, z_t\in\{0,1\}, 
\]
define
\begin{equation}\label{eq.10}
x_j:=\begin{cases}
2z_j-1, &j=1,\ldots,t;\\
1, &j=t+1 .
\end{cases}
\qquad 
\end{equation}
For \(k\in\{t,t+1\}\), define the \(k\times k\) matrix
\begin{equation}\label{eq.8}
\boldsymbol X_k
:=
x_{1:k}x_{1:k}^\top
=
\left[x_ix_j\right]_{i,j=1}^k.
\end{equation}
Then, we have the following results.
\begin{itemize}
  \item[i.]   The joint probability is 
obtained by
\[
  \P{Z_{1:t}=z_{1:t}}=\Phi_t
\left(
\left(
x_1c_1,\ldots,
x_tc_t
\right)^\top;
\boldsymbol I_t
+
\sigma^2\boldsymbol X_t
\right).
  \]
  \item[ii.] The one-step-ahead prediction 
  is obtained by
  \[
\begin{aligned}
&\E{
Z_{t+1}\mid Z_{1:t}=z_{1:t}
} =
\frac{
\Phi_{t+1}
\left(
\left(
x_1c_1,\ldots,
x_tc_t,
x_{t+1}c_{t+1}
\right)^\top;
\boldsymbol I_{t+1}
+
\sigma^2\boldsymbol X_{t+1}
\right)
}{
\Phi_t
\left(
\left(
x_1c_1,\ldots,
x_tc_t
\right)^\top;
\boldsymbol I_t
+
\sigma^2\boldsymbol X_t
\right)
}.
\end{aligned}
\]
\end{itemize}

\end{theorem}

\begin{proof}




We first show part i.
\begin{equation}\label{eq.thm1.history}
\begin{aligned}
\P{Z_{1:t}=z_{1:t}}
&=
\Phi_t
\left(
\left(
x_1c_1,\ldots,
x_tc_t
\right)^\top;
\boldsymbol I_t
+
\sigma^2\boldsymbol X_t
\right).
\end{aligned}
\end{equation}
By the conditional independence of
\(Z_1,\ldots,Z_t\) and the observation equation in Model~\ref{model.1},
\[
\begin{aligned}
\P{Z_{1:t}=z_{1:t}}
&=
\E{
\prod_{j=1}^{t}
\P{Z_j=z_j\mid\Theta}
}\\
&=
\E{
\prod_{j=1}^{t}
\Phi\left(x_j(c_j+\Theta)\right)
}\\
&=
\Phi_t
\left(
\left(
x_1c_1,\ldots,
x_tc_t
\right)^\top;
\boldsymbol I_t
+
\sigma^2\boldsymbol X_t
\right)
\end{aligned}
\]
where
the last equality follows from Lemma \ref{lem.5} and
\[
\left(x_1(c_1+\Theta), \cdots, x_t(c_t+\Theta)\right)^\top \sim \operatorname{MVN}_t\left(
\left(x_1c_1,\cdots,x_tc_t\right)^\top, \sigma^2\boldsymbol X_t
 \right).
\]
This proves
\eqref{eq.thm1.history}.



Now, we prove part ii, by first observing
\begin{equation}\label{eq.1}
\begin{aligned}
\E{Z_{t+1}\mid Z_{1:t}=z_{1:t}}
&=
\P{Z_{t+1}=1\mid Z_{1:t}=z_{1:t}}\\
&=
\frac{
\P{Z_{1:t+1}=(z_1,\ldots,z_t,1)^\top}
}{
\P{Z_{1:t}=z_{1:t}}
}.
\end{aligned}
\end{equation}
where the denominator is in \eqref{eq.thm1.history}. 
Applying the result in part i to the numerator in \eqref{eq.1}, we obtain
\[
\P{Z_{1:t+1}=(z_1,\ldots,z_t,1)^\top}=
\Phi_{t+1}
\left(
\left(
x_1c_1,\ldots,
x_tc_t,
x_{t+1}c_{t+1}
\right)^\top;
\boldsymbol I_{t+1}
+
\sigma^2\boldsymbol X_{t+1}
\right)
\]
which proves the result.
\end{proof}


The following corollary provides a closed-form expression for the
posterior density of the random effect conditional on the observed
binary history under Model~\ref{model.1}. The result follows directly from
Bayes' formula and the joint probability obtained in
part i of Theorem~\ref{thm.1}.



\begin{corollary}[Posterior distribution of the random effect]
\label{cor.2}
Under the setting of Theorem~\ref{thm.1}, the conditional distribution
\[
\Theta\mid z_{1:t}=z_{1:t}
\]
has density \(\pi(\theta\mid Z_{1:t}=z_{1:t})\) at
\(\theta\in\Real\) given by
\begin{equation}\label{eq.posterior.theta.m0.density}
\begin{aligned}
\pi(\theta\mid Z_{1:t}=z_{1:t})
&=
\frac{
\displaystyle
\frac{1}{\sigma}
\phi\left(\frac{\theta}{\sigma}\right)
\prod_{j=1}^t
\Phi\left(x_j(c_j+\theta)\right)
}{
\Phi_t
\left(
(x_1c_1,\ldots,x_tc_t)^\top;
\boldsymbol I_t+\sigma^2\boldsymbol X_t
\right)
}.
\end{aligned}
\end{equation}
\end{corollary}

\begin{proof}
Let \(\pi(\theta)\) denote the prior density of \(\Theta\) at
\(\theta\in\Real\). By Bayes' formula, we have
\[
\begin{aligned}
\pi(\theta\mid Z_{1:t}=z_{1:t})
&=
\frac{
\pi(\theta)
\P{Z_{1:t}=z_{1:t}\mid\Theta=\theta}
}{
\P{Z_{1:t}=z_{1:t}}
}\\
&=
\frac{
\frac{1}{\sigma}
\phi\left(\frac{\theta}{\sigma}\right)
\prod_{j=1}^t
\P{Z_j=z_j\mid\Theta=\theta}
}{
\P{Z_{1:t}=z_{1:t}}
}\\
&=
\frac{
\frac{1}{\sigma}
\phi\left(\frac{\theta}{\sigma}\right)
\prod_{j=1}^t
\Phi\{x_j(c_j+\theta)\}
}{
\P{Z_{1:t}=z_{1:t}}
}.
\end{aligned}
\]
Finally, substituting the expression for
\(\P{Z_{1:t}=z_{1:t}}\) from part i of Theorem \ref{thm.1} gives
\eqref{eq.posterior.theta.m0.density}, which proves the result.
\end{proof}











\section{Bernoulli-lognormal probit bivariate random effect model}\label{sec.6}


In the following we extend Model \ref{model.1} additionally introduce positive part.

\begin{model}[Bernoulli-lognormal probit shared random effect model]
\label{m.2.uni}
Let
\[
\eta,\sigma,\tau>0, \quad 
(c_t)_{t\in\mathbb N},(d_t)_{t\in\mathbb N}\subseteq\Real
\]
be predetermined exogenous constants.
We define the sequence \((Y_t)_{t\in\mathbb N}\)
through the following hierarchical specification.

\begin{itemize}
  \item[i.] Random effect:
  Assume
  \[
  \Theta\sim\operatorname{N}(0,\sigma^2).
  \]

  \item[ii.] Observation equations:
  Conditional on \(\Theta\), assume that the pairs
  \((Z_t,S_t)\), \(t\in\mathbb N\), are mutually independent
  and that, for each \(t\), their components are independent,
  with
  \[
  Z_t\mid\Theta
  \sim
  \operatorname{Ber}\left(\Phi(c_t+\Theta)\right),
  \qquad
  S_t\mid\Theta
  \sim
  \operatorname{LN}(d_t+\eta\Theta,\tau^2).
  \]
  The observed loss is defined by
  \[
  Y_t=Z_tS_t,
  \qquad t\in\mathbb N.
  \]
\end{itemize}
\end{model}

Model~\ref{m.2.uni} is equivalent to Model~\ref{model.2}
with \(\rho=1\), under the identification
\[
\Theta_1=\Theta,
\qquad
\Theta_2=\eta\Theta,
\qquad
\sigma_1=\sigma,
\qquad
\sigma_2=\eta\sigma.
\]
The parameter \(\eta\) allows the shared random effect
to enter the occurrence and log-severity components
with different magnitudes.



The following remark collects the notation used in the results for
Model~\ref{m.2.uni}.

\begin{remark}\label{rem.1}
Under Model~\ref{m.2.uni}, an observed history
\(Y_{1:t}=y_{1:t}\in\Real_{\ge0}^t\) reveals the information on all occurrence
indicators, but only the severities corresponding to positive losses.
For \(j=1,\ldots,t\), define
\begin{equation}\label{eq.12}
z_j:=1_{[y_j>0]},
\qquad
s_j:=
\begin{cases}
y_j, & y_j>0,\\
\mathrm{NA}, & y_j=0,
\end{cases}
\end{equation}
where \(\mathrm{NA}\) indicates that the realization of \(S_j\)
is unobserved.

Let
\begin{equation}\label{eq.14}
\mathcal I_t
:=
\{j\in\{1,\ldots,t\}:y_j>0\},
\qquad
n_t:=|\mathcal I_t|
\end{equation}
denote the set of periods with positive losses and its cardinality.
When \(n_t\ge1\), write
\[
\mathcal I_t=\{i_1,\ldots,i_{n_t}\},
\qquad
i_1<\cdots<i_{n_t},
\]
and define the corresponding severity vector and its observed value by
\[
S_{\mathcal I_t}
:=
(S_{i_1},\ldots,S_{i_{n_t}})^\top,
\qquad
s_{\mathcal I_t}
:=
(y_{i_1},\ldots,y_{i_{n_t}})^\top.
\]
The observed history \(Y_{1:t}=y_{1:t}\) is therefore
equivalently represented by
\[
Z_{1:t}=z_{1:t}
\qquad\text{and}\qquad
S_{\mathcal I_t}=s_{\mathcal I_t}.
\]
When \(n_t=0\), the severity vectors are understood to be empty,
and conditioning on the observed history reduces to conditioning
on \(Z_{1:t}=z_{1:t}\) alone.
\end{remark}


\begin{theorem}[Closed-form one-step mean prediction]
\label{thm.2.uni}
Consider Model~\ref{m.2.uni} with the realized observations
\[
Y_{1:t}=y_{1:t}\in\Real_{\ge 0}^t,
\]
and let
\[
Z_{1:t}=z_{1:t},
\qquad
S_{\mathcal I_t}=s_{\mathcal I_t},
\qquad
\mathcal I_t,
\qquad
n_t
\]
be defined in Remark~\ref{rem.1}.
For the given \(Z_{1:t}=z_{1:t}\), define
\[
x_1,\ldots,x_t,x_{t+1}\in\{-1,1\}
\]
as in \eqref{eq.10}, and let the \(k\times k\) matrix
\(\boldsymbol X_k\) be defined in \eqref{eq.8} for
\(k\in\{t,t+1\}\).
Then, we have the following results.

\begin{itemize}
\item[i.] We have
\begin{equation}\label{eq.posterior.shared}
\Theta
\mid
S_{\mathcal I_t}=s_{\mathcal I_t}
\sim
N(m_t^{[1]},C_t^{[1]}),
\end{equation}
where
\begin{equation}\label{eq.222}
C_t^{[1]}
:=
\left(
\frac{1}{\sigma^2}
+
\frac{n_t\eta^2}{\tau^2}
\right)^{-1}
\quad\hbox{and}\quad
m_t^{[1]}
:=
\frac{\eta C_t^{[1]}}{\tau^2}
\sum_{j\in\mathcal I_t}
(\log s_j-d_j).
\end{equation}


\item[ii.] Furthermore, we have
\[
\begin{aligned}
&\E{
Y_{t+1}
\mid
Y_{1:t}=y_{1:t}
}\\
&\quad=
\exp\left(
d_{t+1}
+
\frac{\tau^2}{2}
+
\eta m_t^{[1]}
+
\frac{\eta^2C_t^{[1]}}{2}
\right)
\frac{
\Phi_{t+1}
\left(
m_{1:t+1}^{[1**]};
\boldsymbol I_{t+1}
+
C_t^{[1]}\boldsymbol X_{t+1}
\right)
}{
\Phi_t
\left(
m_{1:t}^{[1*]};
\boldsymbol I_t
+
C_t^{[1]}\boldsymbol X_t
\right)
},
\end{aligned}
\]
where
\[
m_{1:t}^{[1*]}
:=
\left(
x_1(c_1+m_t^{[1]}),
\ldots,
x_t(c_t+m_t^{[1]})
\right)^\top
\]
and
\[
m_{1:t+1}^{[1**]}
:=
\left(
x_1(c_1+m_t^{[1]}+\eta C_t^{[1]}),
\ldots,
x_{t+1}(c_{t+1}+m_t^{[1]}+\eta C_t^{[1]})
\right)^\top.
\]
\end{itemize}
\end{theorem}

\begin{proof}
We start with the proof of part i.
Using
\[
\log S_j
\mid
\Theta
\sim
N(d_j+\eta\Theta,\tau^2),
\qquad j=1,\ldots,t,
\]
and Bayes' formula, the posterior density of \(\Theta\) at
\(\theta\in\Real\) conditional on
\(S_{\mathcal I_t}=s_{\mathcal I_t}\), denoted by
\(\pi(\theta\mid S_{\mathcal I_t}=s_{\mathcal I_t})\),
is given by
\[
\begin{aligned}
&\pi(
\theta
\mid
S_{\mathcal I_t}=s_{\mathcal I_t}
)\\
&\quad\propto
\pi(\theta)
f(
s_{\mathcal I_t}
\mid
\Theta=\theta
)\\
&\quad\propto
\exp\left(
-\frac{\theta^2}{2\sigma^2}
-
\frac{1}{2\tau^2}
\sum_{j\in\mathcal I_t}
(\log s_j-d_j-\eta\theta)^2
\right)\\
&\quad\propto
\exp\left(
-\frac{1}{2C_t^{[1]}}(\theta-m_t^{[1]})^2
\right),
\end{aligned}
\]
where all proportionalities are with respect to \(\theta\),
and the last proportionality follows from \eqref{eq.222}.
This proves part i.

Now, we prove part ii.
By the definition of conditional expectation, we have
\begin{equation}\label{eq.6}
\begin{aligned}
\E{
Y_{t+1}
\mid
Y_{1:t}=y_{1:t}
}
&=
\E{
Y_{t+1}
\mid
Z_{1:t}=z_{1:t},
S_{\mathcal I_t}=s_{\mathcal I_t}
}\\
&=
\frac{
\E{
Y_{t+1}1_{\left[Z_{1:t}=z_{1:t}\right]}
\mid
S_{\mathcal I_t}=s_{\mathcal I_t}
}
}{
\P{
Z_{1:t}=z_{1:t}
\mid
S_{\mathcal I_t}=s_{\mathcal I_t}
}
}.
\end{aligned}
\end{equation}
We first calculate the denominator
\[
\P{
Z_{1:t}=z_{1:t}
\mid
S_{\mathcal I_t}=s_{\mathcal I_t}
}.
\]
By the conditional independence in Model~\ref{m.2.uni},
\begin{equation}\label{eq.9}
\begin{aligned}
\P{
Z_{1:t}=z_{1:t}
\mid
S_{\mathcal I_t}=s_{\mathcal I_t}
}&=
\E{
\P{
Z_{1:t}=z_{1:t}
\mid
\Theta,
S_{\mathcal I_t}=s_{\mathcal I_t}
}
\mid
S_{\mathcal I_t}=s_{\mathcal I_t}
}\\
&=
\E{
\prod_{j=1}^t
\Phi\left(
x_j(c_j+\Theta)
\right)
\mid
S_{\mathcal I_t}=s_{\mathcal I_t}
}\\
&=
\Phi_t
\left(
m_{1:t}^{[1*]};
\boldsymbol I_t
+
C_t^{[1]}\boldsymbol X_t
\right),
\end{aligned}
\end{equation}
where the last equality follows from Lemma~\ref{lem.5} combined
with the following observation obtained from part i:
\[
\begin{aligned}
&\left[
\left(
x_1(c_1+\Theta),
\ldots,
x_t(c_t+\Theta)
\right)^\top
\ \middle|\
S_{\mathcal I_t}=s_{\mathcal I_t}
\right]\sim
\operatorname{MVN}
\left(
m_{1:t}^{[1*]},
C_t^{[1]}\boldsymbol X_t
\right).
\end{aligned}
\]

Now, we calculate the numerator
\[
\E{
Y_{t+1}1_{\left[Z_{1:t}=z_{1:t}\right]}
\mid
S_{\mathcal I_t}=s_{\mathcal I_t}
}
\]
in \eqref{eq.6}.
First, observe that
\[
\begin{aligned}
\E{Y_{t+1}\mid\Theta}
&=
\E{Z_{t+1}\mid\Theta}
\E{S_{t+1}\mid\Theta}\\
&=
\Phi(c_{t+1}+\Theta)
\exp\left(
d_{t+1}
+
\eta\Theta
+
\frac{\tau^2}{2}
\right).
\end{aligned}
\]
Then, using this and the tower property, we have
\begin{equation}\label{eq.numerator.shared}
\begin{aligned}
&\E{
Y_{t+1}1_{\left[Z_{1:t}=z_{1:t}\right]}
\mid
S_{\mathcal I_t}=s_{\mathcal I_t}
}\\
&\quad=
\E{
\E{Y_{t+1}\mid\Theta}
\P{Z_{1:t}=z_{1:t}\mid\Theta}
\mid
S_{\mathcal I_t}=s_{\mathcal I_t}
}\\
&\quad=
\exp\left(
d_{t+1}
+
\frac{\tau^2}{2}
\right)\times
\E{
\exp(\eta\Theta)
\Phi(c_{t+1}+\Theta)
\prod_{j=1}^t
\Phi\left(
x_j(c_j+\Theta)
\right)
\mid
S_{\mathcal I_t}=s_{\mathcal I_t}
}\\
&\quad=
\exp\left(
d_{t+1}
+
\frac{\tau^2}{2}
\right)
\E{
\exp(\eta\Theta)
\prod_{j=1}^{t+1}
\Phi\left(
x_j(c_j+\Theta)
\right)
\mid
S_{\mathcal I_t}=s_{\mathcal I_t}
},
\end{aligned}
\end{equation}
where the last equality follows from the definition \(x_{t+1}=1\).

%Define
%\[
%g(\theta)
%:=
%\prod_{j=1}^{t+1}
%\Phi\left(
%x_j(c_j+\theta)
%\right),
%\qquad \theta\in\Real.
%\]
%Since \(0\le g(\theta)\le1\), the Gaussian distribution in part i
%ensures that the integrability condition in Lemma~\ref{lem.4}
%is satisfied.
Since the probit function is finite, applying Lemma~\ref{lem.4} with \(h=\eta\) to
\eqref{eq.numerator.shared}, we have
\begin{equation}\label{eq.7}
\begin{aligned}
&\E{
Y_{t+1}1_{\left[Z_{1:t}=z_{1:t}\right]}
\mid
S_{\mathcal I_t}=s_{\mathcal I_t}
}\\
&\quad=
\exp\left(
d_{t+1}
+
\frac{\tau^2}{2}
+
\eta m_t^{[1]}
+
\frac{\eta^2C_t^{[1]}}{2}
\right)\\
&\qquad\times
\E{
\prod_{j=1}^{t+1}
\Phi\left(
x_j(c_j+\widetilde{\Theta})
\right)
\mid
S_{\mathcal I_t}=s_{\mathcal I_t}
},
\end{aligned}
\end{equation}
where
\[
\widetilde{\Theta}
\mid
S_{\mathcal I_t}=s_{\mathcal I_t}
\sim
N(m_t^{[1]}+\eta C_t^{[1]},C_t^{[1]}).
\]

Since
\[
\begin{aligned}
&\left[
\left(
x_1(c_1+\widetilde{\Theta}),
\ldots,
x_{t+1}(c_{t+1}+\widetilde{\Theta})
\right)^\top
\ \middle|\
S_{\mathcal I_t}=s_{\mathcal I_t}
\right]\\
&\quad\sim
\operatorname{MVN}
\left(
m_{1:t+1}^{[1**]},
C_t^{[1]}\boldsymbol X_{t+1}
\right),
\end{aligned}
\]
applying Lemma~\ref{lem.5} to \eqref{eq.7}, we have
\[
\begin{aligned}
&\E{
Y_{t+1}1_{\left[Z_{1:t}=z_{1:t}\right]}
\mid
S_{\mathcal I_t}=s_{\mathcal I_t}
}\\
&\quad=
\exp\left(
d_{t+1}
+
\frac{\tau^2}{2}
+
\eta m_t^{[1]}
+
\frac{\eta^2C_t^{[1]}}{2}
\right)\\
&\qquad\times
\Phi_{t+1}
\left(
m_{1:t+1}^{[1**]};
\boldsymbol I_{t+1}
+
C_t^{[1]}\boldsymbol X_{t+1}
\right).
\end{aligned}
\]
Combining this expression with \eqref{eq.9} in \eqref{eq.6}
concludes the proof of part ii.
\end{proof}



\begin{theorem}[Closed-form one-step stop-loss prediction]
\label{thm.stop.loss.uni}
Under the same setting as in Theorem~\ref{thm.2.uni},
let \(d>0\). Then, we have
\begin{equation}\label{eq.stop.loss.prediction.uni}
\begin{aligned}
&\E{
(Y_{t+1}-d)_+
\mid
Y_{1:t}=y_{1:t}
}\\
&\quad=
\frac{
\begin{aligned}
&\exp\left(
d_{t+1}
+
\frac{\tau^2}{2}
+
\eta m_t^{[1]}
+
\frac{\eta^2C_t^{[1]}}{2}
\right)
\Phi_{t+2}
\left(
\mu_{1:t+2}^{[1*]};
\boldsymbol\Omega_{t+2}^{[1]}
\right)-
d\,
\Phi_{t+2}
\left(
\mu_{1:t+2}^{[1**]};
\boldsymbol\Omega_{t+2}^{[1]}
\right)
\end{aligned}
}{
\Phi_t
\left(
m_{1:t}^{[1*]};
\boldsymbol I_t
+
C_t^{[1]}\boldsymbol X_t
\right)
},
\end{aligned}
\end{equation}
where
\[
\mu_{1:t+2}^{[1*]}
:=
\begin{pmatrix}
m_{1:t+1}^{[1**]}\\[0.15cm]
\dfrac{
d_{t+1}+\eta m_t^{[1]}+\eta^2C_t^{[1]}+\tau^2-\log d
}{\tau}
\end{pmatrix}
\]
and
\[
\mu_{1:t+2}^{[1**]}
:=
\begin{pmatrix}
m_{1:t}^{[1*]}\\[0.1cm]
x_{t+1}(c_{t+1}+m_t^{[1]})\\[0.15cm]
\dfrac{
d_{t+1}+\eta m_t^{[1]}-\log d
}{\tau}
\end{pmatrix},
\]
with
\[
\boldsymbol\Omega_{t+2}^{[1]}
:=
\begin{pmatrix}
\boldsymbol I_{t+1}
+
C_t^{[1]}\boldsymbol X_{t+1}
&
\dfrac{\eta C_t^{[1]}}{\tau}x_{1:t+1}
\\[0.4cm]
\dfrac{\eta C_t^{[1]}}{\tau}x_{1:t+1}^{\top}
&
1+\dfrac{\eta^2C_t^{[1]}}{\tau^2}
\end{pmatrix}.
\]
\end{theorem}

\begin{proof}
%First, observe that, from
%\[
%\left[
%\log S_{t+1}\mid\Theta
%\right]
%\sim
%N(d_{t+1}+\eta\Theta,\tau^2),
%\]
%we have the stop-loss expectation given by
%\[
%\begin{aligned}
%&\E{
%(S_{t+1}-d)_+
%\mid
%\Theta
%}\\
%&\quad=
%\exp\left(
%d_{t+1}
%+
%\eta\Theta
%+
%\frac{\tau^2}{2}
%\right)
%\Phi\left(
%\frac{
%d_{t+1}+\eta\Theta+\tau^2-\log d
%}{\tau}
%\right)-
%d\,
%\Phi\left(
%\frac{
%d_{t+1}+\eta\Theta-\log d
%}{\tau}
%\right).
%\end{aligned}
%\]

Since \(d>0\), observe that we have
\[
(Y_{t+1}-d)_+
=
Z_{t+1}(S_{t+1}-d)_+.
\]
As a result, the conditional independence in
Model~\ref{m.2.uni} and the stop-loss expectation of the normal distribution give
\[
\begin{aligned}
&\E{
(Y_{t+1}-d)_+
\mid
\Theta
}\\
&=
\Phi(c_{t+1}+\Theta)
\Bigg[
\exp\left(
d_{t+1}
+
\eta\Theta
+
\frac{\tau^2}{2}
\right)
\Phi\left(
\frac{
d_{t+1}+\eta\Theta+\tau^2-\log d
}{\tau}
\right)-
d\,
\Phi\left(
\frac{
d_{t+1}+\eta\Theta-\log d
}{\tau}
\right)
\Bigg].
\end{aligned}
\]

Similar to \eqref{eq.6}, we have
\begin{equation}\label{eq.stop.loss.ratio.uni}
\begin{aligned}
&\E{
(Y_{t+1}-d)_+
\mid
Y_{1:t}=y_{1:t}
}\\
&\quad=
\frac{
\E{
(Y_{t+1}-d)_+
1_{\left[Z_{1:t}=z_{1:t}\right]}
\mid
S_{\mathcal I_t}=s_{\mathcal I_t}
}
}{
\P{
Z_{1:t}=z_{1:t}
\mid
S_{\mathcal I_t}=s_{\mathcal I_t}
}
}.
\end{aligned}
\end{equation}
The denominator is obtained from \eqref{eq.9}.

We now calculate the numerator. Using the preceding conditional
stop-loss expectation, the tower property, the conditional
independence in Model~\ref{m.2.uni}, and \(x_{t+1}=1\), we obtain
\begin{equation}\label{eq.stop.loss.numerator.uni}
\begin{aligned}
&\E{
(Y_{t+1}-d)_+
1_{\left[Z_{1:t}=z_{1:t}\right]}
\mid
S_{\mathcal I_t}=s_{\mathcal I_t}
}\\
&\quad=
\exp\left(
d_{t+1}
+
\frac{\tau^2}{2}
\right)
\E{
\exp(\eta\Theta)
\prod_{j=1}^{t+1}
\Phi\left(
x_j(c_j+\Theta)
\right)
\Phi\left(
\frac{
d_{t+1}+\eta\Theta+\tau^2-\log d
}{\tau}
\right)
\ \bigg|\
S_{\mathcal I_t}=s_{\mathcal I_t}
}\\
&\qquad\qquad\qquad-
d\,
\E{
\prod_{j=1}^{t+1}
\Phi\left(
x_j(c_j+\Theta)
\right)
\Phi\left(
\frac{
d_{t+1}+\eta\Theta-\log d
}{\tau}
\right)
\ \bigg|\
S_{\mathcal I_t}=s_{\mathcal I_t}
}.
\end{aligned}
\end{equation}

For the expectation in the first term in
\eqref{eq.stop.loss.numerator.uni}, since the product of the Gaussian
CDFs is bounded by \(1\) and observing the Gaussian distribution in
part i of Theorem~\ref{thm.2.uni}, Lemma~\ref{lem.4} with \(h=\eta\) gives.
\[
\begin{aligned}
&\E{
\exp(\eta\Theta)
\prod_{j=1}^{t+1}
\Phi\left(
x_j(c_j+\Theta)
\right)
\Phi\left(
\frac{
d_{t+1}+\eta\Theta+\tau^2-\log d
}{\tau}
\right)
\ \bigg|\
S_{\mathcal I_t}=s_{\mathcal I_t}
}\\
&\quad=
\exp\left(
\eta m_t^{[1]}+\frac{\eta^2C_t^{[1]}}{2}
\right)
\E{
\prod_{j=1}^{t+1}
\Phi\left(
x_j(c_j+\widetilde{\Theta})
\right)
\Phi\left(
\frac{
d_{t+1}+\eta\widetilde{\Theta}+\tau^2-\log d
}{\tau}
\right)
\ \bigg|\
S_{\mathcal I_t}=s_{\mathcal I_t}
},
\end{aligned}
\]
where
\[
\widetilde{\Theta}
\mid
S_{\mathcal I_t}=s_{\mathcal I_t}
\sim
N(m_t^{[1]}+\eta C_t^{[1]},C_t^{[1]}).
\]

Here,
\[
\left[
\begin{pmatrix}
x_1(c_1+\widetilde{\Theta})\\
\vdots\\
x_{t+1}(c_{t+1}+\widetilde{\Theta})\\[0.15cm]
\dfrac{
d_{t+1}+\eta\widetilde{\Theta}+\tau^2-\log d
}{\tau}
\end{pmatrix}
\ \middle|\
S_{\mathcal I_t}=s_{\mathcal I_t}
\right]
\]
follows the Gaussian distribution with mean
\(\mu_{1:t+2}^{[1*]}\) and covariance matrix
\[
\begin{pmatrix}
C_t^{[1]}\boldsymbol X_{t+1}
&
\dfrac{\eta C_t^{[1]}}{\tau}x_{1:t+1}
\\[0.4cm]
\dfrac{\eta C_t^{[1]}}{\tau}x_{1:t+1}^{\top}
&
\dfrac{\eta^2C_t^{[1]}}{\tau^2}
\end{pmatrix}
=
\boldsymbol\Omega_{t+2}^{[1]}-\boldsymbol I_{t+2},
\]
which follows from the conditional distribution of
\(\widetilde{\Theta}\).

Therefore, Lemma~\ref{lem.5} gives
\begin{equation}\label{eq.stop.loss.first.uni}
\begin{aligned}
&\E{
\exp(\eta\Theta)
\prod_{j=1}^{t+1}
\Phi\left(
x_j(c_j+\Theta)
\right)
\Phi\left(
\frac{
d_{t+1}+\eta\Theta+\tau^2-\log d
}{\tau}
\right)
\ \bigg|\
S_{\mathcal I_t}=s_{\mathcal I_t}
}\\
&\quad=
\exp\left(
\eta m_t^{[1]}+\frac{\eta^2C_t^{[1]}}{2}
\right)
\Phi_{t+2}
\left(
\mu_{1:t+2}^{[1*]};
\boldsymbol\Omega_{t+2}^{[1]}
\right).
\end{aligned}
\end{equation}

Now, we calculate the expectation in the second term in
\eqref{eq.stop.loss.numerator.uni}.
Similar to the distribution in the first term, part i of
Theorem~\ref{thm.2.uni} shows that
\[
\begin{aligned}
&\left[
\begin{pmatrix}
x_1(c_1+\Theta)\\
\vdots\\
x_{t+1}(c_{t+1}+\Theta)\\[0.15cm]
\dfrac{
d_{t+1}+\eta\Theta-\log d
}{\tau}
\end{pmatrix}
\ \middle|\
S_{\mathcal I_t}=s_{\mathcal I_t}
\right]\sim
\operatorname{MVN}_{t+2}
\left(
\mu_{1:t+2}^{[1**]},
\boldsymbol\Omega_{t+2}^{[1]}-\boldsymbol I_{t+2}
\right).
\end{aligned}
\]
Hence, another application of Lemma~\ref{lem.5} gives
\begin{equation}\label{eq.stop.loss.second.uni}
\begin{aligned}
&\E{
\prod_{j=1}^{t+1}
\Phi\left(
x_j(c_j+\Theta)
\right)
\Phi\left(
\frac{
d_{t+1}+\eta\Theta-\log d
}{\tau}
\right)
\ \bigg|\
S_{\mathcal I_t}=s_{\mathcal I_t}
}\\
&\quad=
\Phi_{t+2}
\left(
\mu_{1:t+2}^{[1**]};
\boldsymbol\Omega_{t+2}^{[1]}
\right).
\end{aligned}
\end{equation}

Consequently, substituting \eqref{eq.stop.loss.first.uni} and
\eqref{eq.stop.loss.second.uni} into
\eqref{eq.stop.loss.numerator.uni}, we obtain
\[
\begin{aligned}
&\E{
(Y_{t+1}-d)_+
1_{\left[Z_{1:t}=z_{1:t}\right]}
\mid
S_{\mathcal I_t}=s_{\mathcal I_t}
}\\
&\quad=
\exp\left(
d_{t+1}
+
\frac{\tau^2}{2}
+
\eta m_t^{[1]}
+
\frac{\eta^2C_t^{[1]}}{2}
\right)
\Phi_{t+2}
\left(
\mu_{1:t+2}^{[1*]};
\boldsymbol\Omega_{t+2}^{[1]}
\right)-
d\,
\Phi_{t+2}
\left(
\mu_{1:t+2}^{[1**]};
\boldsymbol\Omega_{t+2}^{[1]}
\right).
\end{aligned}
\]
Dividing this expression by the denominator obtained from
\eqref{eq.9} proves the result.
\end{proof}







Continuous-discrete mixed nature of $Y_t$, the likelihood function of $Y_{1:t}=y_{1:t}$ is defined by
\begin{equation}\label{eq.likelihood.factorization.uni}
\ell(y_{1:t})
:=
\P{
Z_{1:t}=z_{1:t}
\mid
S_{\mathcal I_t}=s_{\mathcal I_t}
}
f(s_{\mathcal I_t}),
\end{equation}
where \(f(s_{\mathcal I_t})\) denotes the probability mass function of $S_{\mathcal I_t}$ evaluated at $s_{\mathcal I_t}$.
When \(\left\vert\mathcal I_t\right\vert=0\), we set \(f(s_{\mathcal I_t})=1\).

\begin{corollary}[Closed-form likelihood]
\label{cor.likelihood.uni}
Under the same setting as in Theorem~\ref{thm.2.uni},
the likelihood in \eqref{eq.likelihood.factorization.uni}
is given by
\begin{equation}\label{eq.likelihood.uni}
\begin{aligned}
\ell(y_{1:t})
&=
\Phi_t
\left(
m_{1:t}^{[1*]};
\boldsymbol I_t
+
C_t^{[1]}\boldsymbol X_t
\right)
\frac{\sqrt{C_t^{[1]}}}
{\sigma\tau^{n_t}(2\pi)^{n_t/2}
\displaystyle\prod_{j\in\mathcal I_t}s_j}
\exp\left(
-\frac{1}{2\tau^2}
\sum_{j\in\mathcal I_t}
(\log s_j-d_j)^2
+
\frac{\left(m_t^{[1]}\right)^2}{2C_t^{[1]}}
\right).
\end{aligned}
\end{equation}
\end{corollary}



\begin{proof}
From \eqref{eq.likelihood.factorization.uni}, it is enough to
calculate
\[
\P{
Z_{1:t}=z_{1:t}
\mid
S_{\mathcal I_t}=s_{\mathcal I_t}
}
\quad\hbox{and}\quad
f(s_{\mathcal I_t}).
\]
First, the conditional probability can be obtained from \eqref{eq.9}.

It remains to calculate \(f(s_{\mathcal I_t})\).
Suppose that \(n_t=\left\vert\mathcal I_t\right\vert>0\).
Under Model~\ref{m.2.uni}, we have
\[
\Theta\sim N(0,\sigma^2),
\qquad
S_j\mid\Theta
\sim
\operatorname{LN}(d_j+\eta\Theta,\tau^2),
\qquad j\in\mathcal I_t.
\]
Using the conditional independence of the severities given
\(\Theta\), we obtain
\[
\begin{aligned}
f(s_{\mathcal I_t})
&=
\int_{\Real}
\frac{1}{\sigma}
\phi\left(\frac{\theta}{\sigma}\right)
\prod_{j\in\mathcal I_t}
f(s_j\mid\Theta=\theta)
\,{\rm d}\theta\\
&=
\frac{1}
{\sigma\tau^{n_t}(2\pi)^{(n_t+1)/2}
\displaystyle\prod_{j\in\mathcal I_t}s_j}
\int_{\Real}
\exp\left(
-\frac{\theta^2}{2\sigma^2}
-\frac{1}{2\tau^2}
\sum_{j\in\mathcal I_t}
(\log s_j-d_j-\eta\theta)^2
\right)
\,{\rm d}\theta\\
&=
\frac{1}
{\sigma\tau^{n_t}(2\pi)^{(n_t+1)/2}
\displaystyle\prod_{j\in\mathcal I_t}s_j}
\int_{\Real}
\exp\left(
-\frac{(\theta-m_t^{[1]})^2}{2C_t^{[1]}}
-\frac{1}{2\tau^2}
\sum_{j\in\mathcal I_t}
(\log s_j-d_j)^2
+\frac{\left(m_t^{[1]}\right)^2}{2C_t^{[1]}}
\right)
\,{\rm d}\theta\\
&=
\frac{\sqrt{C_t^{[1]}}}
{\sigma\tau^{n_t}(2\pi)^{n_t/2}
\displaystyle\prod_{j\in\mathcal I_t}s_j}
\exp\left(
-\frac{1}{2\tau^2}
\sum_{j\in\mathcal I_t}
(\log s_j-d_j)^2
+\frac{\left(m_t^{[1]}\right)^2}{2C_t^{[1]}}
\right),
\end{aligned}
\]
where the third equality follows from \eqref{eq.222}.
When \(n_t=0\), we have \(C_t^{[1]}=\sigma^2\) and
\(m_t^{[2]}=0\), so the same expression gives
\(f(s_{\mathcal I_t})=1\), in accordance with the definition.


Combining this expression with \eqref{eq.9} in
\eqref{eq.likelihood.factorization.uni} proves the result.
\end{proof}


\begin{corollary}[Posterior distribution of the random effect]
\label{cor.posterior.uni}
Under the setting of Theorem~\ref{thm.2.uni},
the conditional distribution
\[
\Theta\mid Y_{1:t}=y_{1:t}
\]
has density \(\pi(\theta\mid Y_{1:t}=y_{1:t})\) at
\(\theta\in\Real\) given by
\begin{equation}\label{eq.posterior.theta.uni.density}
\begin{aligned}
\pi(\theta\mid Y_{1:t}=y_{1:t})
&=
\frac{
\displaystyle
\frac{1}{\sqrt{C_t^{[1]}}}
\phi\left(
\frac{\theta-m_t^{[1]}}{\sqrt{C_t^{[1]}}}
\right)
\prod_{j=1}^t
\Phi\left(
x_j(c_j+\theta)
\right)
}{
\Phi_t
\left(
m_{1:t}^{[1*]};
\boldsymbol I_t
+
C_t^{[1]}\boldsymbol X_t
\right)
}.
\end{aligned}
\end{equation}
\end{corollary}

\begin{proof}
Using Remark~\ref{rem.1} and the conditional independence in
Model~\ref{m.2.uni}, Bayes' formula gives
\[
\begin{aligned}
\pi(
\theta
\mid
Y_{1:t}=y_{1:t}
)&=
\frac{
\pi(
\theta
\mid
S_{\mathcal I_t}=s_{\mathcal I_t}
)
\P{
Z_{1:t}=z_{1:t}
\mid
\Theta=\theta
}
}{
\P{
Z_{1:t}=z_{1:t}
\mid
S_{\mathcal I_t}=s_{\mathcal I_t}
}
}\\
&=
\frac{
\pi(
\theta
\mid
S_{\mathcal I_t}=s_{\mathcal I_t}
)
\displaystyle\prod_{j=1}^t
\Phi\left(
x_j(c_j+\theta)
\right)
}{
\P{
Z_{1:t}=z_{1:t}
\mid
S_{\mathcal I_t}=s_{\mathcal I_t}
}
}.
\end{aligned}
\]

By part i of Theorem~\ref{thm.2.uni}, we have
\[
\pi(
\theta
\mid
S_{\mathcal I_t}=s_{\mathcal I_t}
)
=
\frac{1}{\sqrt{C_t^{[1]}}}
\phi\left(
\frac{\theta-m_t^{[1]}}{\sqrt{C_t^{[1]}}}
\right).
\]
Substituting this expression and \eqref{eq.9} into Bayes' formula
gives \eqref{eq.posterior.theta.uni.density}, which proves the result.
\end{proof}

\section{Extension to bivariate random effect}




































































\begin{model}[Bernoulli-lognormal probit bivariate random effect model]
\label{model.2}
Let
\[
\rho\in[-1,1], \quad 
\sigma_1, \sigma_2, \tau>0, \quad 
(c_t)_{t\in\mathbb N},(d_t)_{t\in\mathbb N}\subseteq\Real
\]
be predetermined exogenous constants.
We extend Model~\ref{m.2.uni} by replacing the shared random
effect with a bivariate random effect as follows.

\begin{itemize}
  \item[i.] Random effect:
  Assume
  \[
  \Theta_{1:2}
  \sim
  \operatorname{MVN}_2
  \left(
  0_{1:2},
  \begin{pmatrix}
  \sigma_1^2 & \rho\sigma_1\sigma_2\\
  \rho\sigma_1\sigma_2 & \sigma_2^2
  \end{pmatrix}
  \right).
  \]

  \item[ii.] Observation equations:
  Retain the conditional independence assumptions of
  Model~\ref{m.2.uni}, now conditional on \(\Theta_{1:2}\),
  and replace the observation equations by
  \[
  Z_t\mid\Theta_{1:2}
  \sim
  \operatorname{Ber}\left(\Phi(c_t+\Theta_1)\right),
  \qquad
  S_t\mid\Theta_{1:2}
  \sim
  \operatorname{LN}(d_t+\Theta_2,\tau^2).
  \]
  As before, the observed loss is defined by
  \[
  Y_t=Z_tS_t,
  \qquad t\in\mathbb N.
  \]
\end{itemize}
\end{model}



\begin{theorem}[Closed-form one-step mean prediction]
\label{thm.biv}
Consider Model~\ref{model.2}, with \(\rho\in(-1,1)\), and the realized
observations
\[
Y_{1:t}=y_{1:t}\in\Real_{\ge 0}^t,
\]
and let
\[
Z_{1:t}=z_{1:t},
\qquad
S_{\mathcal I_t}=s_{\mathcal I_t},
\qquad
\mathcal I_t,
\qquad
n_t
\]
be defined in Remark~\ref{rem.1}.
For this given
\(Z_{1:t}=z_{1:t}\), define
\[
x_1,\ldots,x_t,x_{t+1}\in\{-1,1\}
\]
as in \eqref{eq.10}, and let the \(k\times k\) matrix
\(\boldsymbol X_k\) be defined in \eqref{eq.8} for
\(k\in\{t,t+1\}\).
Then, we have the following results.

\begin{itemize}
\item[i.] We have
\begin{equation}
\label{eq2.1}
\begin{aligned}
&\Theta_{1:2}
\mid
S_{\mathcal I_t}=s_{\mathcal I_t}\sim
\operatorname{MVN}_2
\left(
\begin{pmatrix}
\dfrac{\rho\sigma_1}{\sigma_2}m_t^{[2]}\\[0.2cm]
m_t^{[2]}
\end{pmatrix},
\begin{pmatrix}
\sigma_1^2(1-\rho^2)
+
\dfrac{\rho^2\sigma_1^2}{\sigma_2^2}C_t^{[2]}
&
\dfrac{\rho\sigma_1}{\sigma_2}C_t^{[2]}\\[0.3cm]
\dfrac{\rho\sigma_1}{\sigma_2}C_t^{[2]}
&
C_t^{[2]}
\end{pmatrix}
\right),
\end{aligned}
\end{equation}
where 
\begin{equation}\label{eq.22}
C_t^{[2]}
:=
\left(
\frac{1}{\sigma_2^2}
+
\frac{n_t}{\tau^2}
\right)^{-1}
\quad\hbox{and}\quad
m_t^{[2]}
:=
\frac{C_t^{[2]}}{\tau^2}
\sum_{j\in\mathcal I_t}
(\log s_j-d_j).
\end{equation}





\item[ii.] Furthermore, we have
\[
\begin{aligned}
&\E{
Y_{t+1}
\mid
Y_{1:t}=y_{1:t}
}\\
&\quad=
\exp\left(
d_{t+1}
+
\frac{\tau^2}{2}
+
m_t^{[2]}
+
\frac{C_t^{[2]}}{2}
\right)
\frac{
\Phi_{t+1}
\left(
m_{1:t+1}^{[2**]};
\boldsymbol I_{t+1}
+
\left\{
\sigma_1^2(1-\rho^2)
+
\frac{\rho^2\sigma_1^2}{\sigma_2^2}C_t^{[2]}
\right\}
\boldsymbol X_{t+1}
\right)
}{
\Phi_t
\left(
m_{1:t}^{[1*]}
;
\boldsymbol I_t
+
\left\{
\sigma_1^2(1-\rho^2)
+
\frac{\rho^2\sigma_1^2}{\sigma_2^2}C_t^{[2]}
\right\}
\boldsymbol X_t
\right)
}.
\end{aligned}
\]
where
\[
m_{1:t}^{[2*]}:=\left(x_1\left(
c_1+
\dfrac{\rho\sigma_1}{\sigma_2}m_t^{[2]}
\right), \cdots,  x_t\left(
c_t+
\dfrac{\rho\sigma_1}{\sigma_2}m_t^{[2]}
\right)\right)^\top
\]
and
\[
m_{1:t+1}^{[2**]}:=\left(
x_1\left(
c_1+
\dfrac{\rho\sigma_1}{\sigma_2}(m_t^{[2]}+C_t^{[2]})
\right),\cdots,
x_{t+1}\left(
c_{t+1}+
\dfrac{\rho\sigma_1}{\sigma_2}(m_t^{[2]}+C_t^{[2]})
\right)
\right)^\top
\]
\end{itemize}
\end{theorem}

\begin{proof}
First, observe that 
the same calculation as in part i of Theorem~\ref{thm.2.uni} gives
\[
\Theta_2
\mid
S_{\mathcal I_t}=s_{\mathcal I_t}
\sim
N(m_t^{[2]},C_t^{[2]}).
\]
Moreover, the conditional
independence of \(\Theta_1\) and \(S_{\mathcal I_t}\) given
\(\Theta_2\) implies
\[
\left[\Theta_1
\mid
\Theta_2,
S_{\mathcal I_t}=s_{\mathcal I_t}\right]
\sim
N\left(
\frac{\rho\sigma_1}{\sigma_2}\Theta_2,
\sigma_1^2(1-\rho^2)
\right).
\]
Combining these two conditional Gaussian distributions gives
\eqref{eq2.1}, which proves part i.

Now, we prove part ii.
As in the proof of Theorem~\ref{thm.2.uni}, we have
\begin{equation}\label{eq.conditional.mean.biv}
\E{
Y_{t+1}
\mid
Y_{1:t}=y_{1:t}
}
=
\frac{
\E{
Y_{t+1}1_{\left[Z_{1:t}=z_{1:t}\right]}
\mid
S_{\mathcal I_t}=s_{\mathcal I_t}
}
}{
\P{
Z_{1:t}=z_{1:t}
\mid
S_{\mathcal I_t}=s_{\mathcal I_t}
}
}.
\end{equation}

We first calculate the denominator.
By part i,
\begin{equation}\label{eq.24}
\left[\Theta_1
\mid
S_{\mathcal I_t}=s_{\mathcal I_t}\right]
\sim
N\left(
\frac{\rho\sigma_1}{\sigma_2}m_t^{[2]},
\sigma_1^2(1-\rho^2)
+
\frac{\rho^2\sigma_1^2}{\sigma_2^2}C_t^{[2]}
\right).
\end{equation}
Hence, the conditional independence in Model~\ref{model.2}
and Lemma~\ref{lem.5} give
\begin{equation}\label{eq.denominator.biv}
\begin{aligned}
\P{
Z_{1:t}=z_{1:t}
\mid
S_{\mathcal I_t}=s_{\mathcal I_t}
}
&=
\E{
\prod_{j=1}^t
\Phi\left(
x_j(c_j+\Theta_1)
\right)
\mid
S_{\mathcal I_t}=s_{\mathcal I_t}
}\\
&=
\Phi_t
\left(
m_{1:t}^{[2*]};
\boldsymbol I_t
+
\left(
\sigma_1^2(1-\rho^2)
+
\frac{\rho^2\sigma_1^2}{\sigma_2^2}C_t^{[2]}
\right)
\boldsymbol X_t
\right)
\end{aligned}
\end{equation}
where the last equality follows from
\[
\begin{aligned}
&\left[
\left(
x_1(c_1+\Theta_1),
\ldots,
x_t(c_t+\Theta_1)
\right)^\top
\ \middle|\
S_{\mathcal I_t}=s_{\mathcal I_t}
\right]\sim
\operatorname{MVN}
\left(
m_{1:t}^{[2*]},
\left(
\sigma_1^2(1-\rho^2)
+
\frac{\rho^2\sigma_1^2}{\sigma_2^2}C_t^{[2]}
\right)
\boldsymbol X_t
\right).
\end{aligned}
\]


Now, we calculate the numerator.
Using the conditional independence in Model~\ref{model.2}
and \(x_{t+1}=1\), the same argument as in
\eqref{eq.numerator.shared} gives
\begin{equation}\label{eq.numerator.biv}
\begin{aligned}
&\E{
Y_{t+1}1_{\left[Z_{1:t}=z_{1:t}\right]}
\mid
S_{\mathcal I_t}=s_{\mathcal I_t}
}\\
&\quad=
\exp\left(
d_{t+1}
+
\frac{\tau^2}{2}
\right)
\E{
\exp(\Theta_2)
\prod_{j=1}^{t+1}
\Phi\left(
x_j(c_j+\Theta_1)
\right)
\mid
S_{\mathcal I_t}=s_{\mathcal I_t}
}.
\end{aligned}
\end{equation}


Since the product of the Gaussian CDFs is bounded by \(1\),
Lemma~\ref{lem.3} with \(h_{1:2}=(0,1)^\top\) to the
conditional expectation in \eqref{eq.numerator.biv}, we have
\[
\begin{aligned}
&\E{
\exp(\Theta_2)
\prod_{j=1}^{t+1}
\Phi\left(
x_j(c_j+\Theta_1)
\right)
\mid
S_{\mathcal I_t}=s_{\mathcal I_t}
}=
\exp\left(
m_t^{[2]}+\frac{C_t^{[2]}}{2}
\right)
\E{
\prod_{j=1}^{t+1}
\Phi\left(
x_j(c_j+\widetilde{\Theta}_1)
\right)
\mid
S_{\mathcal I_t}=s_{\mathcal I_t}
},
\end{aligned}
\]
where
\begin{equation}\label{eq.23}
\begin{aligned}
&\left[\widetilde{\Theta}_{1:2}
\mid
S_{\mathcal I_t}=s_{\mathcal I_t}\right]\sim
\operatorname{MVN}_2
\left(
\begin{pmatrix}
\dfrac{\rho\sigma_1}{\sigma_2}(m_t^{[2]}+C_t^{[2]})\\[0.2cm]
m_t^{[2]}+C_t^{[2]}
\end{pmatrix},
\begin{pmatrix}
\sigma_1^2(1-\rho^2)
+
\dfrac{\rho^2\sigma_1^2}{\sigma_2^2}C_t^{[2]}
&
\dfrac{\rho\sigma_1}{\sigma_2}C_t^{[2]}\\[0.3cm]
\dfrac{\rho\sigma_1}{\sigma_2}C_t^{[2]}
&
C_t^{[2]}
\end{pmatrix}
\right).
\end{aligned}
\end{equation}
Applying Lemma~\ref{lem.5}, we obtain
\begin{equation}\label{eq.tilting.biv}
\begin{aligned}
&\E{
Y_{t+1}1_{\left[Z_{1:t}=z_{1:t}\right]}
\mid
S_{\mathcal I_t}=s_{\mathcal I_t}
}\\
&\quad=
\exp\left(
d_{t+1}
+
\frac{\tau^2}{2}
+
m_t^{[2]}
+
\frac{C_t^{[2]}}{2}
\right)
\Phi_{t+1}
\left(
m_{1:t+1}^{[2**]};
\boldsymbol I_{t+1}
+
\left(
\sigma_1^2(1-\rho^2)
+
\frac{\rho^2\sigma_1^2}{\sigma_2^2}C_t^{[2]}
\right)
\boldsymbol X_{t+1}
\right).
\end{aligned}
\end{equation}
Combining this expression with \eqref{eq.denominator.biv}
in \eqref{eq.conditional.mean.biv} concludes the proof of part ii.
\end{proof}









































































\begin{theorem}[Closed-form one-step stop-loss prediction]
\label{cor.stop.loss.biv}
Under the same setting as in Theorem~\ref{thm.biv}, let \(d>0\),  we have
\[
\begin{aligned}
&\E{
(Y_{t+1}-d)_+
\mid
Y_{1:t}=y_{1:t}
}\\
&\quad=
\frac{
\exp\left(
d_{t+1}
+
\frac{\tau^2}{2}
+
m_t^{[2]}
+
\frac{C_t^{[2]}}{2}
\right)
\Phi_{t+2}
\left(
\mu_{1:t+2}^{[2*]};
\boldsymbol\Omega_{t+2}^{[2]}
\right)-
d\,
\Phi_{t+2}
\left(
\mu_{1:t+2}^{[2**]};
\boldsymbol\Omega_{t+2}^{[2]}
\right)
}{
\Phi_t
\left(
m_{1:t}^{[2*]};
\boldsymbol I_t
+
\left(
\sigma_1^2(1-\rho^2)
+
\frac{\rho^2\sigma_1^2}{\sigma_2^2}C_t^{[2]}
\right)
\boldsymbol X_t
\right)
}
\end{aligned}
\]
where
\[
\mu_{1:t+2}^{[2*]}:=\begin{pmatrix}
m_{1:t+1}^{[2**]}\\[0.15cm]
\dfrac{
d_{t+1}+m_t^{[2]}+C_t^{[2]}+\tau^2-\log d
}{\tau}
\end{pmatrix}
\quad\hbox{and}\quad
\mu_{1:t+2}^{[2**]}:=\begin{pmatrix}
m_{1:t}^{[2*]}\\[0.1cm]
x_{t+1}\left(
c_{t+1}
+
\dfrac{\rho\sigma_1}{\sigma_2}m_t^{[2]}
\right)\\[0.15cm]
\dfrac{
d_{t+1}+m_t^{[2]}-\log d
}{\tau}
\end{pmatrix}
\]
and
\[
\boldsymbol\Omega_{t+2}^{[2]}
:=
\begin{pmatrix}
\boldsymbol I_{t+1}
+
\left(
\sigma_1^2(1-\rho^2)
+
\dfrac{\rho^2\sigma_1^2}{\sigma_2^2}C_t^{[2]}
\right)
\boldsymbol X_{t+1}
&
\dfrac{\rho\sigma_1C_t^{[2]}}{\sigma_2\tau}
x_{1:t+1}
\\[0.4cm]
\dfrac{\rho\sigma_1C_t^{[2]}}{\sigma_2\tau}
x_{1:t+1}^{\top}
&
1+\dfrac{C_t^{[2]}}{\tau^2}
\end{pmatrix}.
\]
\end{theorem}

\begin{proof}
As in the proof of Theorem~\ref{thm.stop.loss.uni}, we have
\[
\E{
(Y_{t+1}-d)_+
\mid
Y_{1:t}=y_{1:t}
}
=
\frac{
\E{
(Y_{t+1}-d)_+
1_{\left[Z_{1:t}=z_{1:t}\right]}
\mid
S_{\mathcal I_t}=s_{\mathcal I_t}
}
}{
\P{
Z_{1:t}=z_{1:t}
\mid
S_{\mathcal I_t}=s_{\mathcal I_t}
}
}.
\]
The denominator is obtained from \eqref{eq.denominator.biv}.

We now calculate the numerator.
The same conditional stop-loss calculation as in
\eqref{eq.stop.loss.numerator.uni}, using the conditional
independence in Model~\ref{model.2} and \(x_{t+1}=1\), gives
\begin{equation}\label{eq.25}
\begin{aligned}
&\E{
(Y_{t+1}-d)_+
1_{\left[Z_{1:t}=z_{1:t}\right]}
\mid
S_{\mathcal I_t}=s_{\mathcal I_t}
}\\
&\quad=
\exp\left(
d_{t+1}+\frac{\tau^2}{2}
\right)
\E{
\exp(\Theta_2)
\prod_{j=1}^{t+1}
\Phi\left(
x_j(c_j+\Theta_1)
\right)
\Phi\left(
\frac{
d_{t+1}+\Theta_2+\tau^2-\log d
}{\tau}
\right)
\ \bigg|\
S_{\mathcal I_t}=s_{\mathcal I_t}
}\\
&\qquad-
d\,
\E{
\prod_{j=1}^{t+1}
\Phi\left(
x_j(c_j+\Theta_1)
\right)
\Phi\left(
\frac{
d_{t+1}+\Theta_2-\log d
}{\tau}
\right)
\ \bigg|\
S_{\mathcal I_t}=s_{\mathcal I_t}
}.
\end{aligned}
\end{equation}

For the expectation in the first term, the product of the
Gaussian CDFs is bounded by \(1\), Lemma~\ref{lem.3} with \(h_{1:2}=(0,1)^\top\), followed by
Lemma~\ref{lem.5}, we obtain
\begin{equation}\label{eq.26}
\begin{aligned}
&\E{
\exp(\Theta_2)
\prod_{j=1}^{t+1}
\Phi\left(
x_j(c_j+\Theta_1)
\right)
\Phi\left(
\frac{
d_{t+1}+\Theta_2+\tau^2-\log d
}{\tau}
\right)
\ \bigg|\
S_{\mathcal I_t}=s_{\mathcal I_t}
}\\
&\quad=
\exp\left(
m_t^{[2]}+\frac{C_t^{[2]}}{2}
\right)
\E{
\prod_{j=1}^{t+1}
\Phi\left(
x_j(c_j+\widetilde{\Theta}_1)
\right)
\Phi\left(
\frac{
d_{t+1}+\widetilde{\Theta}_2+\tau^2-\log d
}{\tau}
\right)
\ \bigg|\
S_{\mathcal I_t}=s_{\mathcal I_t}
}\\
&\quad=
\exp\left(
m_t^{[2]}+\frac{C_t^{[2]}}{2}
\right)
\Phi_{t+2}
\left(
\mu_{1:t+2}^{[2*]};
\boldsymbol\Omega_{t+2}^{[2]}
\right).
\end{aligned}
\end{equation}
Here, the conditional distribution of
\(\widetilde{\Theta}_{1:2}\) is given by \eqref{eq.23},
and the last equality follows from the following observation
\[
\begin{aligned}
&\left[
\begin{pmatrix}
x_1(c_1+\widetilde{\Theta}_1)\\
\vdots\\
x_{t+1}(c_{t+1}+\widetilde{\Theta}_1)\\[0.15cm]
\dfrac{
d_{t+1}+\widetilde{\Theta}_2+\tau^2-\log d
}{\tau}
\end{pmatrix}
\ \middle|\
S_{\mathcal I_t}=s_{\mathcal I_t}
\right]\sim
\operatorname{MVN}_{t+2}
\left(
\mu_{1:t+2}^{[2*]},
\boldsymbol\Omega_{t+2}^{[2]}-\boldsymbol I_{t+2}
\right).
\end{aligned}
\]

For the expectation in the second term in \eqref{eq.25},
part i of Theorem~\ref{thm.biv} gives
\[
\begin{aligned}
&\left[
\begin{pmatrix}
x_1(c_1+\Theta_1)\\
\vdots\\
x_{t+1}(c_{t+1}+\Theta_1)\\[0.15cm]
\dfrac{
d_{t+1}+\Theta_2-\log d
}{\tau}
\end{pmatrix}
\ \middle|\
S_{\mathcal I_t}=s_{\mathcal I_t}
\right]\sim
\operatorname{MVN}_{t+2}
\left(
\mu_{1:t+2}^{[2**]},
\boldsymbol\Omega_{t+2}^{[2]}-\boldsymbol I_{t+2}
\right).
\end{aligned}
\]
Hence, applying Lemma~\ref{lem.5}, we have
\begin{equation}\label{eq.27}
\begin{aligned}
&\E{
\prod_{j=1}^{t+1}
\Phi\left(
x_j(c_j+\Theta_1)
\right)
\Phi\left(
\frac{
d_{t+1}+\Theta_2-\log d
}{\tau}
\right)
\ \bigg|\
S_{\mathcal I_t}=s_{\mathcal I_t}
}=
\Phi_{t+2}
\left(
\mu_{1:t+2}^{[2**]};
\boldsymbol\Omega_{t+2}^{[2]}
\right).
\end{aligned}
\end{equation}

Substituting \eqref{eq.26} and \eqref{eq.27} into \eqref{eq.25}
and dividing by the denominator in \eqref{eq.denominator.biv}
proves the result.
\end{proof}





\begin{corollary}[Closed-form likelihood]
\label{cor2.1}
Under the same setting as in Theorem~\ref{thm.biv}, the likelihood
in \eqref{eq.likelihood.factorization.uni} is given by
\[
\begin{aligned}
\ell(y_{1:t})
&=
\Phi_t
\left(
m_{1:t}^{[2*]};
\boldsymbol I_t
+
\left\{
\sigma_1^2(1-\rho^2)
+
\frac{\rho^2\sigma_1^2}{\sigma_2^2}C_t^{[2]}
\right\}
\boldsymbol X_t
\right)\\
&\quad\times
\frac{\sqrt{C_t^{[2]}}}
{\sigma_2\tau^{n_t}(2\pi)^{n_t/2}
\displaystyle\prod_{j\in\mathcal I_t}s_j}
\exp\left\{
-\frac{1}{2\tau^2}
\sum_{j\in\mathcal I_t}
(\log s_j-d_j)^2
+
\frac{\left(m_t^{[2]}\right)^2}{2C_t^{[2]}}
\right\}.
\end{aligned}
\]
\end{corollary}


\begin{proof}
The conditional probability in \eqref{eq.likelihood.factorization.uni}
is given by \eqref{eq.denominator.biv}.
It remains to calculate \(f(s_{\mathcal I_t})\).

Since the conditional distribution of the severities depends on
the random effect only through \(\Theta_2\sim N(0,\sigma_2^2)\), 
the similar calculation as in the proof of Corollary~\ref{cor.likelihood.uni} for \(n_t>0\) gives 
\[
\begin{aligned}
f(s_{\mathcal I_t})
&=
\frac{1}
{\sigma_2\tau^{n_t}(2\pi)^{(n_t+1)/2}
\displaystyle\prod_{j\in\mathcal I_t}s_j}
\int_{\Real}
\exp\left(
-\frac{\theta_2^2}{2\sigma_2^2}
-\frac{1}{2\tau^2}
\sum_{j\in\mathcal I_t}
(\log s_j-d_j-\theta_2)^2
\right)
\,{\rm d}\theta_2\\
&=
\frac{\sqrt{C_t^{[2]}}}
{\sigma_2\tau^{n_t}(2\pi)^{n_t/2}
\displaystyle\prod_{j\in\mathcal I_t}s_j}
\exp\left(
-\frac{1}{2\tau^2}
\sum_{j\in\mathcal I_t}
(\log s_j-d_j)^2
+
\frac{\left(m_t^{[2]}\right)^2}{2C_t^{[2]}}
\right),
\end{aligned}
\]
where the last equality follows by completing the square using
\eqref{eq.22}. 
When \(n_t=0\), we have \(C_t^{[2]}=\sigma_2^2\) and
\(m_t^{[2]}=0\), so the this expression gives
\(f(s_{\mathcal I_t})=1\), in accordance with the definition.
Combining this expression with \eqref{eq.denominator.biv}
in \eqref{eq.likelihood.factorization.uni} proves the result.
\end{proof}

\begin{corollary}[Posterior distribution of the random effect]
\label{cor.61}
Under the setting of Theorem~\ref{thm.biv}, the conditional distribution
\(
\Theta_{1:2}\mid Y_{1:t}=y_{1:t}
\)
has density
\(\pi(\theta_{1:2}\mid Y_{1:t}=y_{1:t})\) at
\(
\theta_{1:2}=(\theta_1,\theta_2)^\top\in\Real^2
\)
given by
\begin{equation}
\label{eq.posterior.theta.biv.density}
\begin{aligned}
&\pi(\theta_{1:2}\mid Y_{1:t}=y_{1:t})\\
&\quad=
\frac{
\displaystyle
\phi_2
\left(
\begin{pmatrix}
\theta_1-\dfrac{\rho\sigma_1}{\sigma_2}m_t^{[2]}\\[0.2cm]
\theta_2-m_t^{[2]}
\end{pmatrix};
\begin{pmatrix}
\sigma_1^2(1-\rho^2)
+
\dfrac{\rho^2\sigma_1^2}{\sigma_2^2}C_t^{[2]}
&
\dfrac{\rho\sigma_1}{\sigma_2}C_t^{[2]}\\[0.3cm]
\dfrac{\rho\sigma_1}{\sigma_2}C_t^{[2]}
&
C_t^{[2]}
\end{pmatrix}
\right)
\displaystyle\prod_{j=1}^t
\Phi\left(
x_j(c_j+\theta_1)
\right)
}{
\displaystyle
\Phi_t
\left(
m_{1:t}^{[2*]};
\boldsymbol I_t
+
\left(
\sigma_1^2(1-\rho^2)
+
\frac{\rho^2\sigma_1^2}{\sigma_2^2}C_t^{[2]}
\right)
\boldsymbol X_t
\right)
}.
\end{aligned}
\end{equation}
\end{corollary}

\begin{proof}
As in the proof of Corollary~\ref{cor.posterior.uni},
Bayes' formula and the conditional independence in
Model~\ref{model.2} give
\[
\begin{aligned}
&\pi(
\theta_{1:2}
\mid
Y_{1:t}=y_{1:t}
)\\
&\quad=
\frac{
\pi(
\theta_{1:2}
\mid
S_{\mathcal I_t}=s_{\mathcal I_t}
)
\displaystyle\prod_{j=1}^t
\Phi\left(
x_j(c_j+\theta_1)
\right)
}{
\P{
Z_{1:t}=z_{1:t}
\mid
S_{\mathcal I_t}=s_{\mathcal I_t}
}
}.
\end{aligned}
\]
Substituting the Gaussian density from part i of
Theorem~\ref{thm.biv} and the conditional probability in
\eqref{eq.denominator.biv} gives
\eqref{eq.posterior.theta.biv.density}, which proves the result.
\end{proof}

%\begin{corollary}[Posterior distribution of the random effect]
%\label{cor.61}
%Under the setting of Theorem~\ref{thm.biv}, the conditional distribution
%\(
%\Theta_{1:2}\mid Y_{1:t}=y_{1:t}
%\)
%has density
%\(\pi(\theta_{1:2}\mid Y_{1:t}=y_{1:t})\) at
%\(
%\theta_{1:2} = (\theta_1,\theta_2)^\top
%\in\Real^2
%\)
%given by
%\begin{equation}
%\label{eq.posterior.theta.biv.density}
%\begin{aligned}
%&\pi(\theta_{1:2}\mid Y_{1:t}=y_{1:t})\\
%&\quad=
%\frac{
%\displaystyle
%\phi_2
%\left(
%\begin{pmatrix}
%\theta_1-\dfrac{\rho\sigma_1}{\sigma_2}m_t\\[0.2cm]
%\theta_2-m_t
%\end{pmatrix};
%\begin{pmatrix}
%\sigma_1^2(1-\rho^2)
%+
%\dfrac{\rho^2\sigma_1^2}{\sigma_2^2}C_t
%&
%\dfrac{\rho\sigma_1}{\sigma_2}C_t\\[0.3cm]
%\dfrac{\rho\sigma_1}{\sigma_2}C_t
%&
%C_t
%\end{pmatrix}
%\right)
%\displaystyle\prod_{j=1}^t
%\Phi\left(
%x_j(c_j+\theta_1)
%\right)
%}{
%\displaystyle
%\Phi_t
%\left(
%m_{1:t}^{*};
%\boldsymbol I_t
%+
%\left(
%\sigma_1^2(1-\rho^2)
%+
%\frac{\rho^2\sigma_1^2}{\sigma_2^2}C_t
%\right)
%\boldsymbol X_t
%\right)
%}.
%\end{aligned}
%\end{equation}
%\end{corollary}
%
%\begin{proof}
%Bayes' formula gives
%\[
%\begin{aligned}
%&\pi(
%\theta_{1:2}
%\mid
%Y_{1:t}=y_{1:t}
%)\\
%&\quad=
%\frac{
%\pi(
%\theta_{1:2}
%\mid
%S_{\mathcal I_t}=s_{\mathcal I_t}
%)
%\P{Z_{1:t}=z_{1:t}\mid \Theta_{1:2}=\theta_{1:2}}
%}{
%\P{
%Z_{1:t}=z_{1:t}
%\mid
%S_{\mathcal I_t}=s_{\mathcal I_t}
%}
%}\\
%&\quad=
%\frac{
%\pi(
%\theta_{1:2}
%\mid
%S_{\mathcal I_t}=s_{\mathcal I_t}
%)
%\displaystyle\prod_{j=1}^t
%\Phi\left(
%x_j(c_j+\theta_1)
%\right)
%}{
%\P{
%Z_{1:t}=z_{1:t}
%\mid
%S_{\mathcal I_t}=s_{\mathcal I_t}
%}
%}.
%\end{aligned}
%\]
%
%By part i of Theorem~\ref{thm.biv}, we have
%\[
%\begin{aligned}
%&\pi(
%\theta_{1:2}
%\mid
%S_{\mathcal I_t}=s_{\mathcal I_t}
%)\\
%&\quad=
%\phi_2
%\left(
%\begin{pmatrix}
%\theta_1-\dfrac{\rho\sigma_1}{\sigma_2}m_t\\[0.2cm]
%\theta_2-m_t
%\end{pmatrix};
%\begin{pmatrix}
%\sigma_1^2(1-\rho^2)
%+
%\dfrac{\rho^2\sigma_1^2}{\sigma_2^2}C_t
%&
%\dfrac{\rho\sigma_1}{\sigma_2}C_t\\[0.3cm]
%\dfrac{\rho\sigma_1}{\sigma_2}C_t
%&
%C_t
%\end{pmatrix}
%\right).
%\end{aligned}
%\]
%Moreover, substituting this and \eqref{eq.denominator.biv} into Bayes' formula gives
%\eqref{eq.posterior.theta.biv.density}, which proves the result.
%\end{proof}

%
%\section{State-space expansion}\label{sec.5}
%
%Model~\ref{model.2} assumes that the latent risks affecting the
%occurrence and positive parts remain constant over time. We now extend
%this model by allowing the two latent risks to evolve according to a
%stationary bivariate Gaussian VAR(1) process. The two components may
%have different persistence parameters, whereas their contemporaneous
%dependence is represented by their stationary correlation.
%
%\begin{model}[Bernoulli-lognormal probit stationary VAR(1) state-space model]
%\label{model.3}
%Consider the predetermined exogenous constants
%\[
%\eta_1,\eta_2\in(-1,1),
%\qquad
%q_{11},q_{22},\tau\in\Real_{>0},
%\qquad
%q_{12}\in\Real, \qquad (c_t)_{t\in\mathbb N},
%(d_t)_{t\in\mathbb N}
%\subseteq\Real
%\]
%satisfying
%\[
%q_{12}^2<q_{11}q_{22}.
%\]
%We define the sequence of observation
%\((Y_t)_{t\in\mathbb N}\) and the sequence of latent variables $(\Theta_t^{[1]}, \Theta_t^{[2]} )_{t\in\mathbb N}$ through the following evolution and
%observation equations.
%
%\begin{itemize}
%\item[i.] Evolution equation:
%Define the innovations of the latent variables by 
%\[
%\begin{pmatrix}
%  \epsilon_{t}^{[1]}\\
%  \epsilon_{t}^{[2]}
%\end{pmatrix}
%\sim
%\operatorname{MVN}_2
%\left(
%0_{1:2},
%\begin{pmatrix}
%q_{11}&q_{12}\\
%q_{12}&q_{22}
%\end{pmatrix}
%\right),
%\qquad t=1,2,\ldots
%\]
%
%The latent process is a stationary Gaussian
%VAR(1) process initialized by
%\[
%\begin{pmatrix}
%  \Theta_1^{[1]} \\
%  \Theta_1^{[2]}
%\end{pmatrix}
%\sim
%\operatorname{MVN}_2
%\left(
%0_{1:2},
%\begin{pmatrix}
%\dfrac{q_{11}}{1-\eta_1^2}
%&
%\dfrac{q_{12}}{1-\eta_1\eta_2}
%\\[0.4cm]
%\dfrac{q_{12}}{1-\eta_1\eta_2}
%&
%\dfrac{q_{22}}{1-\eta_2^2}
%\end{pmatrix}
%\right),
%\]
%and, for $t\in\mathbb N$, evolve with
%\[
%\begin{pmatrix}
%  \Theta_{t+1}^{[1]} \\
%  \Theta_{t+1}^{[2]}
%\end{pmatrix}
%=
%\begin{pmatrix}
%\eta_1&0\\
%0&\eta_2
%\end{pmatrix}\begin{pmatrix}
%  \Theta_t^{[1]} \\
%  \Theta_t^{[2]}
%\end{pmatrix}+\begin{pmatrix}
%  \epsilon_{t+1}^{[1]}\\
%  \epsilon_{t+1}^{[2]}
%\end{pmatrix}.
%\]
%
%
%\item[ii.] Observation equation:
%Conditional on the latent process
%$(\Theta_t^{[1]}, \Theta_t^{[2]} )_{t\in\mathbb N}$, assume that
%\[
%(Z_t,S_t),
%\qquad t=1,2,\ldots,
%\]
%are mutually independent over \(t\), and that \(Z_t\) and \(S_t\) are
%conditionally independent given the latent process, with
%\[
%Z_t
%\mid
%(\Theta_t^{[1]}, \Theta_t^{[2]} )_{t\in\mathbb N}
%\sim
%\operatorname{Ber}
%\left(
%\Phi(c_t+\Theta_t^{[1]})
%\right),
%\]
%and
%\[
%S_t
%\mid
%(\Theta_t^{[1]}, \Theta_t^{[2]} )_{t\in\mathbb N}
%\sim
%\operatorname{LN}
%\left(
%d_t+\Theta_t^{[2]},
%\tau^2
%\right).
%\]
%Finally, define
%\[
%Y_t=Z_tS_t,
%\qquad t=1,2,\ldots.
%\]
%\end{itemize}
%\end{model}
%
%
%
%\begin{theorem}[Closed-form one-step mean prediction]
%\label{thm.biv.ssm}
%Consider Model~\ref{model.3} with the realized observations
%\[
%Y_{1:t}=y_{1:t}\in\Real_{\ge0}^t,
%\]
%and let
%\[
%Z_{1:t}=z_{1:t},
%\qquad
%S_{\mathcal I_t}=s_{\mathcal I_t},
%\qquad
%\mathcal I_t,
%\qquad
%n_t
%\]
%be defined in Remark~\ref{rem.1}. For the given realization
%\(Z_{1:t}=z_{1:t}\), define
%\[
%x_1,\ldots,x_t,x_{t+1}\in\{-1,1\}
%\]
%as in \eqref{eq.10}, and define
%\[
%\boldsymbol M_t
%:=
%\operatorname{diag}(z_1,\ldots,z_t).
%\]
%Denote the covariance matrix of $$\begin{pmatrix}
%  \Theta_{1:t}^{[1]}\\
%   \Theta_{1:t}^{[2]}
%\end{pmatrix} \equiv \left(\Theta_1^{[1]},\cdots,\Theta_t^{[1]}, \Theta_1^{[2]},\ldots,\Theta_t^{[2]} \right)^\top$$ by
%\begin{equation}\label{eq.33}
%\begin{aligned}
%\boldsymbol\Gamma_t
%:=
%\begin{pmatrix}
%\dfrac{q_{11}}{1-\eta_1^2}
%\left[
%\eta_1^{|i-j|}
%\right]_{i,j=1}^t
%&
%\dfrac{q_{12}}{1-\eta_1\eta_2}
%\left[
%\begin{cases}
%\eta_1^{i-j}, & i\ge j,\\
%\eta_2^{j-i}, & i<j
%\end{cases}
%\right]_{i,j=1}^t
%\\[0.6cm]
%\dfrac{q_{12}}{1-\eta_1\eta_2}
%\left[
%\begin{cases}
%\eta_2^{i-j}, & i\ge j,\\
%\eta_1^{j-i}, & i<j
%\end{cases}
%\right]_{i,j=1}^t
%&
%\dfrac{q_{22}}{1-\eta_2^2}
%\left[
%\eta_2^{|i-j|}
%\right]_{i,j=1}^t
%\end{pmatrix}.
%\end{aligned}
%\end{equation}
%Then, we have the following
%results.
%
%\begin{itemize}
%\item[i.] Define
%\[
%r_j
%:=
%\begin{cases}
%\log s_j-d_j, & z_j=1,\\
%0, & z_j=0,
%\end{cases}
%\qquad j=1,\ldots,t.
%\]
%Then,
%\begin{equation}
%\label{eq.posterior.biv.ssm.severity}
%\begin{pmatrix}
%  \Theta_{1:t}^{[1]}\\
%   \Theta_{1:t}^{[2]}
%\end{pmatrix} 
%\,\bigg\vert\,
%S_{\mathcal I_t}=s_{\mathcal I_t}
%\sim
%\operatorname{MVN}_{2t}
%\left(
%\begin{pmatrix} \mu_{1:t}^{[1]}\\ \mu_{1:t}^{[2]} \end{pmatrix},
%\boldsymbol H_t
%\right),
%\end{equation}
%where
%\[
%\boldsymbol H_t
%:=
%\left[
%\left(
%\boldsymbol\Gamma_t
%\right)^{-1}
%+
%\frac{1}{\tau^2}
%\begin{pmatrix}
%0_{t\times t}&0_{t\times t}\\
%0_{t\times t}&\boldsymbol M_t
%\end{pmatrix}
%\right]^{-1}
%\]
%and
%\[
%\begin{pmatrix} \mu_{1:t}^{[1]}\\ \mu_{1:t}^{[2]} \end{pmatrix}\equiv(\mu_1^{[1]}, \cdots, \mu_t^{[1]}, 
%\mu_1^{[2]}, \cdots, \mu_t^{[2]})^\top
%:=
%\boldsymbol H_t
%\frac{1}{\tau^2}
%\begin{pmatrix}
%0_{1:t}^{\top}\\
%\boldsymbol M_t r_{1:t}
%\end{pmatrix}\in\Real^{2t}.
%\]
%
%\item[ii.] Define
%\[
%\boldsymbol D_t
%:=
%\operatorname{diag}(x_1,\ldots,x_t)\quad\hbox{and}\quad
%\boldsymbol B_t
%:=
%\begin{pmatrix}
%\boldsymbol D_t&0_{t\times t}
%\end{pmatrix}.
%\]
%Furthermore, define
%\[
%a_{1:t}
%:=
%(x_1c_1,\ldots,x_tc_t)^\top
%, \qquad 
%a_{1:t+1}^{\ast}
%:=
%\left(
%x_1c_1,\ldots,x_tc_t,
%x_{t+1}
%\frac{
%c_{t+1}+q_{12}
%}{
%\sqrt{1+q_{11}}
%}
%\right)^\top
%\]
%and
%\[
%\boldsymbol B_{t+1}^{\ast}
%:=
%\begin{pmatrix}
%\boldsymbol D_t&0_{t\times t}\\[0.2cm]
%x_{t+1}
%\dfrac{\eta_1}{\sqrt{1+q_{11}}}
%e_{1:t}^\top
%&
%0_{1\times t}
%\end{pmatrix}.
%\]
%Then,
%\[
%\begin{aligned}
%&\E{
%Y_{t+1}
%\mid
%Y_{1:t}=y_{1:t}
%}\\
%&\quad=
%\exp\left(
%d_{t+1}
%+
%\frac{\tau^2}{2}
%+
%\frac{q_{22}}{2}
%+\eta_2\mu_t^{[2]}
%+
%\frac{\eta_2^2}{2}
%%(0, \cdots, 0, \eta_2)
%%\boldsymbol H_t
%%(0, \cdots, 0, \eta_2)^\top
%\left[\boldsymbol H_t\right]_{2t,2t}
%\right)\\
%&\qquad\times
%\frac{
%\Phi_{t+1}
%\left(
%a_{1:t+1}^{\ast}
%+
%\boldsymbol B_{t+1}^{\ast}
%\left(
%\begin{pmatrix} \mu_{1:t}^{[1]}\\ \mu_{1:t}^{[2]} \end{pmatrix}
%+
%\eta_2\boldsymbol H_t e_{1:2t}
%\right);
%\boldsymbol I_{t+1}
%+
%\boldsymbol B_{t+1}^{\ast}
%\boldsymbol H_t
%(\boldsymbol B_{t+1}^{\ast})^\top
%\right)
%}{
%\Phi_t
%\left(
%a_{1:t}
%+
%\boldsymbol B_t\begin{pmatrix} \mu_{1:t}^{[1]}\\ \mu_{1:t}^{[2]} \end{pmatrix};
%\boldsymbol I_t
%+
%\boldsymbol B_t
%\boldsymbol H_t
%\boldsymbol B_t^\top
%\right)
%}.
%\end{aligned}
%\]
%\end{itemize}
%\end{theorem}
%
%\begin{proof}
%First, we note that the straightforward calculation shows that $\boldsymbol\Gamma_t$ defined in \eqref{eq.33} becomes the covariance matrix of $\begin{pmatrix}
%  \Theta_{1:t}^{[1]}\\
%   \Theta_{1:t}^{[2]}
%\end{pmatrix}$. 
%Now, we prove part i. By the stationary Gaussian VAR(1) assumption in
%Model~\ref{model.3}, one can easily show
%\[
%\begin{pmatrix} \Theta_{1:t}^{[1]}\\ \Theta_{1:t}^{[2]} \end{pmatrix}
%\sim
%\operatorname{MVN}_{2t}
%\left(
%0_{1:2t},
%\boldsymbol\Gamma_t
%\right).
%\]
%Moreover,
%\[
%\left.\log S_j
%\middle\vert
%\begin{pmatrix} \Theta_{1:t}^{[1]}\\ \Theta_{1:t}^{[2]} \end{pmatrix}
%\right.
%\sim
%N\left(
%d_j+\Theta_j^{[2]},
%\tau^2
%\right),
%\qquad j=1,\ldots,t.
%\]
%Therefore, using Bayes' formula, as a function of
%\[
%\begin{pmatrix}
%\theta_{1:t}^{[1]}\\
%\theta_{1:t}^{[2]}
%\end{pmatrix}
%\in\Real^{2t},
%\]
%we have
%\[
%\begin{aligned}
%&\pi\left(
%\theta_{1:t}^{[1]}, 
%\theta_{1:t}^{[2]}
%\mid
%S_{\mathcal I_t}=s_{\mathcal I_t}
%\right)\\
%&\quad\propto
%\exp\left(
%-\frac12
%\begin{pmatrix} \theta_{1:t}^{[1]}\\ \theta_{1:t}^{[2]} \end{pmatrix}^\top
%\left(
%\boldsymbol\Gamma_t
%\right)^{-1}
%\begin{pmatrix}
%\theta_{1:t}^{[1]}\\
%\theta_{1:t}^{[2]}
%\end{pmatrix}
%-
%\frac{1}{2\tau^2}
%\left(
%r_{1:t}-\theta_{1:t}^{[2]}
%\right)^\top
%\boldsymbol M_t
%\left(
%r_{1:t}-\theta_{1:t}^{[2]}
%\right)
%\right)\\
%&\quad\propto
%\exp\left(
%-\frac12\left(
%\begin{pmatrix} \theta_{1:t}^{[1]}\\ \theta_{1:t}^{[2]} \end{pmatrix}^\top
%\left[
%\left(
%\boldsymbol\Gamma_t
%\right)^{-1}
%+
%\frac{1}{\tau^2}
%\begin{pmatrix}
%\boldsymbol 0_{t\times t}
%&
%\boldsymbol 0_{t\times t}\\
%\boldsymbol 0_{t\times t}
%&
%\boldsymbol M_t
%\end{pmatrix}
%\right]
%\begin{pmatrix} \theta_{1:t}^{[1]}\\ \theta_{1:t}^{[2]} \end{pmatrix}
%-
%\frac{2}{\tau^2}
%r_{1:t}^\top
%\boldsymbol M_t
%\theta_{1:t}^{[2]}\right)
%\right)\\
%&\quad=
%\exp\left(
%-\frac12
%\begin{pmatrix} \theta_{1:t}^{[1]}\\ \theta_{1:t}^{[2]} \end{pmatrix}^\top
%\boldsymbol H_t^{-1}
%\begin{pmatrix} \theta_{1:t}^{[1]}\\ \theta_{1:t}^{[2]} \end{pmatrix}
%+
%\begin{pmatrix} \theta_{1:t}^{[1]}\\ \theta_{1:t}^{[2]} \end{pmatrix}^\top
%\boldsymbol H_t^{-1}
%\begin{pmatrix} \mu_{1:t}^{[1]}\\ \mu_{1:t}^{[2]} \end{pmatrix}
%\right)\\
%&\quad\propto
%\exp\left(
%-\frac12
%\left(
%\begin{pmatrix} \theta_{1:t}^{[1]}\\ \theta_{1:t}^{[2]} \end{pmatrix}^\top-\begin{pmatrix} \mu_{1:t}^{[1]}\\ \mu_{1:t}^{[2]} \end{pmatrix}^\top
%\right)
%\boldsymbol H_t^{-1}
%\left(
%\begin{pmatrix} \theta_{1:t}^{[1]}\\ \theta_{1:t}^{[2]} \end{pmatrix}-\begin{pmatrix} \mu_{1:t}^{[1]}\\ \mu_{1:t}^{[2]} \end{pmatrix}
%\right)
%\right)
%\end{aligned}
%\]
%where the equality is from the definition of $\boldsymbol H_t$ and \(\begin{pmatrix} \mu_{1:t}^{[1]}\\ \mu_{1:t}^{[2]} \end{pmatrix}\).
%This proves part i.
%
%Now, we prove part ii. By the definition of conditional expectation,
%we have
%\begin{equation}
%\label{eq.conditional.mean.biv.ssm}
%\begin{aligned}
%&\E{
%Y_{t+1}
%\mid
%Y_{1:t}=y_{1:t}
%}\\
%&\quad=
%\E{
%Y_{t+1}
%\mid
%Z_{1:t}=z_{1:t},
%S_{\mathcal I_t}=s_{\mathcal I_t}
%}\\
%&\quad=
%\frac{
%\E{
%Y_{t+1}
%1_{[Z_{1:t}=z_{1:t}]}
%\mid
%S_{\mathcal I_t}=s_{\mathcal I_t}
%}
%}{
%\P{
%Z_{1:t}=z_{1:t}
%\mid
%S_{\mathcal I_t}=s_{\mathcal I_t}
%}
%}.
%\end{aligned}
%\end{equation}
%
%We first calculate the denominator in
%\eqref{eq.conditional.mean.biv.ssm}. By the conditional independence in
%Model~\ref{model.3}, we have
%\[
%\begin{aligned}
%&\P{
%Z_{1:t}=z_{1:t}
%\mid
%S_{\mathcal I_t}=s_{\mathcal I_t}
%}\\
%&\quad=
%\E{
%\P{
%Z_{1:t}=z_{1:t}
%\,\bigg\vert\,
%\begin{pmatrix} \Theta_{1:t}^{[1]}\\ \Theta_{1:t}^{[2]} \end{pmatrix},
%S_{\mathcal I_t}=s_{\mathcal I_t}
%}
%\middle\vert
%S_{\mathcal I_t}=s_{\mathcal I_t}
%}\\
%&\quad=
%\E{
%\prod_{j=1}^t
%\Phi\left(
%x_j(c_j+\Theta_j^{[1]})
%\right)
%\middle\vert
%S_{\mathcal I_t}=s_{\mathcal I_t}
%}.
%\end{aligned}
%\]
%From the following observation
%\[
%\left( x_1c_1+x_1\Theta_{1}^{[1]}, \cdots, x_tc_t+x_t\Theta_{t}^{[1]}\right)^\top=
%a_{1:t}
%+
%\boldsymbol B_t\begin{pmatrix} \Theta_{1:t}^{[1]}\\ \Theta_{1:t}^{[2]} \end{pmatrix}
%\]
% the result in part i gives
%\[
%\begin{aligned}
%&\left[
%a_{1:t}
%+
%\boldsymbol B_t\begin{pmatrix} \Theta_{1:t}^{[1]}\\ \Theta_{1:t}^{[2]} \end{pmatrix}
%\ \middle|\
%S_{\mathcal I_t}=s_{\mathcal I_t}
%\right]\sim
%\operatorname{MVN}_t
%\left(
%a_{1:t}
%+
%\boldsymbol B_t\begin{pmatrix} \mu_{1:t}^{[1]}\\ \mu_{1:t}^{[2]} \end{pmatrix},
%\boldsymbol B_t
%\boldsymbol H_t
%\boldsymbol B_t^\top
%\right).
%\end{aligned}
%\]
%Therefore, Lemma~\ref{lem.5} gives
%\begin{equation}
%\label{eq.denominator.biv.ssm}
%\begin{aligned}
%&\P{
%Z_{1:t}=z_{1:t}
%\mid
%S_{\mathcal I_t}=s_{\mathcal I_t}
%}=
%\Phi_t
%\left(
%a_{1:t}
%+
%\boldsymbol B_t\begin{pmatrix} \mu_{1:t}^{[1]}\\ \mu_{1:t}^{[2]} \end{pmatrix};
%\boldsymbol I_t
%+
%\boldsymbol B_t
%\boldsymbol H_t
%\boldsymbol B_t^\top
%\right).
%\end{aligned}
%\end{equation}
%
%Next, we calculate the numerator in
%\eqref{eq.conditional.mean.biv.ssm}. 
%We start by showing
%\begin{equation}
%\label{eq.future.mean.biv.ssm}
%\begin{aligned}
%&\E{
%Y_{t+1}
%\middle\vert
%\begin{pmatrix} \Theta_{t}^{[1]}\\ \Theta_{t}^{[2]} \end{pmatrix}
%}=
%\exp\left(
%d_{t+1}
%+
%\frac{\tau^2}{2}
%+
%\frac{q_{22}}{2}
%\right)
%\exp\left(
%\eta_2\Theta_t^{[2]}
%\right)
%\Phi\left(
%\frac{
%c_{t+1}
%+
%q_{12}
%+
%\eta_1\Theta_t^{[1]}
%}{
%\sqrt{1+q_{11}}
%}
%\right).
%\end{aligned}
%\end{equation}
%From the evolution equation in
%Model~\ref{model.3}, we have
%\[
%\left[
%\begin{pmatrix}
%\Theta_{t+1}^{[1]}\\
%\Theta_{t+1}^{[2]}
%\end{pmatrix}
%\ \middle|\
%\begin{pmatrix} \Theta_{t}^{[1]}\\ \Theta_{t}^{[2]} \end{pmatrix}
%\right]
%\sim
%\operatorname{MVN}_2
%\left(
%\begin{pmatrix}
%\eta_1\Theta_t^{[1]}\\
%\eta_2\Theta_t^{[2]}
%\end{pmatrix},
%\begin{pmatrix}
%q_{11}&q_{12}\\
%q_{12}&q_{22}
%\end{pmatrix}
%\right).
%\]
%Using the conditional independence between \(Z_{t+1}\) and \(S_{t+1}\)
%given \((\Theta_{t+1}^{[1]}, \Theta_{t+1}^{[2]})\), we have
%\begin{equation}\label{eq.31}
%\begin{aligned}
%&\E{
%Y_{t+1}
%\middle\vert
%\begin{pmatrix} \Theta_{t}^{[1]}\\ \Theta_{t}^{[2]} \end{pmatrix}
%}\\
%&\quad=
%\E{
%\E{
%Z_{t+1}S_{t+1}
%\middle\vert
%\begin{pmatrix} \Theta_{t+1}^{[1]}\\ \Theta_{t+1}^{[2]} \end{pmatrix}
%}
%\middle \vert
%\begin{pmatrix} \Theta_{t}^{[1]}\\ \Theta_{t}^{[2]} \end{pmatrix}
%}\\
%&\quad=
%\exp\left(
%d_{t+1}+\frac{\tau^2}{2}
%\right)
%\E{
%\exp\left(\Theta_{t+1}^{[2]}\right)
%\Phi\left(c_{t+1}+\Theta_{t+1}^{[1]}\right)
%\middle\vert
%\begin{pmatrix} \Theta_{t}^{[1]}\\ \Theta_{t}^{[2]} \end{pmatrix}
%}.
%\end{aligned}
%\end{equation}
%
%Applying Lemma~\ref{lem.3} conditionally on \((\Theta_{t}^{[1]}, \Theta_{t}^{[2]})\), with the
%tilting vector \((0,1)^\top\), gives
%\[
%\begin{aligned}
%&\E{
%\exp\left(\Theta_{t+1}^{[2]}\right)
%\Phi\left(c_{t+1}+\Theta_{t+1}^{[1]}\right)
%\middle\vert
%\begin{pmatrix} \Theta_{t}^{[1]}\\ \Theta_{t}^{[2]} \end{pmatrix}
%}\\
%&\quad=
%\exp\left(
%\eta_2\Theta_t^{[2]}
%+
%\frac{q_{22}}{2}
%\right)
%\E{
%\Phi\left(
%c_{t+1}+\widetilde{\Theta}_{t+1}^{[1]}
%\right)
%\middle\vert
%\begin{pmatrix} \Theta_{t}^{[1]}\\ \Theta_{t}^{[2]} \end{pmatrix}
%},
%\end{aligned}
%\]
%where
%\[
%\widetilde{\Theta}_{t+1}^{[1]}
%\,\bigg\vert\,
%\begin{pmatrix} \Theta_{t}^{[1]}\\ \Theta_{t}^{[2]} \end{pmatrix}
%\sim
%N\left(
%\eta_1\Theta_t^{[1]}+q_{12},
%q_{11}
%\right).
%\]
%Hence, Lemma~\ref{lem.5} gives
%\[
%\begin{aligned}
%\E{
%\exp\left(\Theta_{t+1}^{[2]}\right)
%\Phi\left(c_{t+1}+\Theta_{t+1}^{[1]}\right)
%\middle\vert
%\begin{pmatrix} \Theta_{t}^{[1]}\\ \Theta_{t}^{[2]} \end{pmatrix}
%} &= 
%\exp\left(\eta_2\Theta_t^{[2]}
%+
%\frac{q_{22}}{2}\right)
%\E{
%\Phi\left(
%c_{t+1}+\widetilde{\Theta}_{t+1}^{[1]}
%\right)
%\middle\vert
%\begin{pmatrix} \Theta_{t}^{[1]}\\ \Theta_{t}^{[2]} \end{pmatrix}
%}\\
%&=
%\exp\left(\eta_2\Theta_t^{[2]}
%+
%\frac{q_{22}}{2}\right)
%\Phi\left(
%\frac{
%c_{t+1}
%+
%q_{12}
%+
%\eta_1\Theta_t^{[1]}
%}{
%\sqrt{1+q_{11}}
%}
%\right)
%\end{aligned}
%\]
%which together with \eqref{eq.31} shows \eqref{eq.future.mean.biv.ssm}.
%
%
%Now, we come back to the numerator
%in \eqref{eq.conditional.mean.biv.ssm} which becomes
%\begin{equation}\label{eq.41}
%\begin{aligned}
%&\E{
%Y_{t+1}
%1_{[Z_{1:t}=z_{1:t}]}
%\mid
%S_{\mathcal I_t}=s_{\mathcal I_t}
%}\\
%&\quad=
%\exp\left(
%d_{t+1}
%+
%\frac{\tau^2}{2}
%+
%\frac{q_{22}}{2}
%\right)\\
%&\qquad\times
%\E{
%\exp\left(
%\eta_2\Theta_t^{[2]}
%\right)
%\prod_{j=1}^t
%\Phi\left(
%x_j(c_j+\Theta_j^{[1]})
%\right)
%\Phi\left(
%x_{t+1}
%\frac{
%c_{t+1}
%+
%q_{12}
%+
%\eta_1\Theta_t^{[1]}
%}{
%\sqrt{1+q_{11}}
%}
%\right)
% \bigg\vert
%S_{\mathcal I_t}=s_{\mathcal I_t}
%}\\
%&\quad=
%\exp\left(
%d_{t+1}
%+
%\frac{\tau^2}{2}
%+
%\frac{q_{22}}{2}
%\right)\\
%&\qquad\times
%\E{
%\exp\left(
%\eta_2 e_{1:2t}^\top 
%\begin{pmatrix} \Theta_{1:t}^{[1]}\\ \Theta_{1:t}^{[2]} \end{pmatrix}
%\right)
%\prod_{j=1}^{t+1}
%\Phi\left(
%\left[
%a_{1:t+1}^{\ast}
%+
%\boldsymbol B_{t+1}^{\ast}
%\begin{pmatrix} \Theta_{1:t}^{[1]}\\ \Theta_{1:t}^{[2]} \end{pmatrix}
%\right]_j
%\right)
%\ \bigg|\
%S_{\mathcal I_t}=s_{\mathcal I_t}
%}
%\end{aligned}
%\end{equation}
%where the first equality is from \eqref{eq.future.mean.biv.ssm}, and 
%the second equality is from the definitions of \(a_{1:t+1}^{\ast}\) and
%\(\boldsymbol B_{t+1}^{\ast}\).
%
%Since the product of the probit terms is bounded by one, part i and
%Lemma~\ref{lem.3} give
%\[
%\begin{aligned}
%&\E{
%\exp\left(
%\eta_2 e_{1:2t}^\top\begin{pmatrix} \Theta_{1:t}^{[1]}\\ \Theta_{1:t}^{[2]} \end{pmatrix}
%\right)
%\prod_{j=1}^{t+1}
%\Phi\left(
%\left[
%a_{1:t+1}^{\ast}
%+
%\boldsymbol B_{t+1}^{\ast}
%\begin{pmatrix} \Theta_{1:t}^{[1]}\\ \Theta_{1:t}^{[2]} \end{pmatrix}
%\right]_j
%\right)
%\ \bigg|\
%S_{\mathcal I_t}=s_{\mathcal I_t}
%}\\
%&\quad=
%\exp\left(
%\eta_2 e_{1:2t}^\top\begin{pmatrix} \mu_{1:t}^{[1]}\\ \mu_{1:t}^{[2]} \end{pmatrix}
%+
%\frac{\eta_2^2}{2}
%e_{1:2t}^\top
%\boldsymbol H_t
%e_{1:2t}
%\right)\\
%&\qquad\times
%\E{
%\prod_{j=1}^{t+1}
%\Phi\left(
%\left[
%a_{1:t+1}^{\ast}
%+
%\boldsymbol B_{t+1}^{\ast}
%\begin{pmatrix} \widetilde\Theta_{1:t}^{[1]}\\ \widetilde\Theta_{1:t}^{[2]} \end{pmatrix}
%\right]_j
%\right)
%\ \bigg|\
%S_{\mathcal I_t}=s_{\mathcal I_t}
%},
%\end{aligned}
%\]
%where
%\[
%\left.
%\begin{pmatrix} \widetilde\Theta_{1:t}^{[1]}\\ \widetilde\Theta_{1:t}^{[2]} \end{pmatrix}
%\right\vert
%S_{\mathcal I_t}=s_{\mathcal I_t}
%\sim
%\operatorname{MVN}_{2t}
%\left(
%\begin{pmatrix} \mu_{1:t}^{[1]}\\ \mu_{1:t}^{[2]} \end{pmatrix}
%+
%\eta_2 \boldsymbol H_t e_{1:2t},
%\boldsymbol H_t
%\right).
%\]
%Finally, applying Lemma~\ref{lem.5} to \eqref{eq.41} and using the following 
%\[
%\begin{aligned}
%&\left[
%a_{1:t+1}^{\ast}
%+
%\boldsymbol B_{t+1}^{\ast}
%\begin{pmatrix} \widetilde\Theta_{1:t}^{[1]}\\ \widetilde\Theta_{1:t}^{[2]} \end{pmatrix}
%\ \middle|\
%S_{\mathcal I_t}=s_{\mathcal I_t}
%\right]\\
%&\quad\sim
%\operatorname{MVN}_{t+1}
%\left(
%a_{1:t+1}^{\ast}
%+
%\boldsymbol B_{t+1}^{\ast}
%\left(
%\begin{pmatrix} \mu_{1:t}^{[1]}\\ \mu_{1:t}^{[2]} \end{pmatrix}
%+
%\eta_2\boldsymbol H_t e_{1:2t}
%\right),
%\boldsymbol B_{t+1}^{\ast}
%\boldsymbol H_t
%(\boldsymbol B_{t+1}^{\ast})^\top
%\right),
%\end{aligned}
%\]
%we have
%\[
%\begin{aligned}
%&\E{
%Y_{t+1}
%1_{[Z_{1:t}=z_{1:t}]}
%\mid
%S_{\mathcal I_t}=s_{\mathcal I_t}
%}\\
%&\quad=
%\exp\left(
%d_{t+1}
%+
%\frac{\tau^2}{2}
%+
%\frac{q_{22}}{2}
%+
%\eta_2e_{1:2t}^\top\begin{pmatrix} \mu_{1:t}^{[1]}\\ \mu_{1:t}^{[2]} \end{pmatrix}
%+
%\frac{\eta_2^2}{2}
%e_{1:2t}^\top
%\boldsymbol H_t e_{1:2t}
%\right)\\
%&\qquad\times
%\Phi_{t+1}
%\left(
%a_{1:t+1}^{\ast}
%+
%\boldsymbol B_{t+1}^{\ast}
%\left(
%\begin{pmatrix} \mu_{1:t}^{[1]}\\ \mu_{1:t}^{[2]} \end{pmatrix}
%+
%\eta_2\boldsymbol H_t e_{1:2t}
%\right);
%\boldsymbol I_{t+1}
%+
%\boldsymbol B_{t+1}^{\ast}
%\boldsymbol H_t
%(\boldsymbol B_{t+1}^{\ast})^\top
%\right).
%\end{aligned}
%\]
%Combining this expression with \eqref{eq.denominator.biv.ssm} in
%\eqref{eq.conditional.mean.biv.ssm} proves part ii.
%\end{proof}
%
%\begin{theorem}[Closed-form one-step stop-loss prediction]
%\label{cor.stop.loss.biv.ssm}
%Under the same setting as in Theorem~\ref{thm.biv.ssm}, let \(d>0\).
%Define
%\[
%\begin{aligned}
%\boldsymbol\Omega_{t+2}^{\rm ssm}
%&:=
%\begin{pmatrix}
%\boldsymbol I_t
%&
%0_{t\times1}
%&
%0_{t\times1}
%\\
%0_{1\times t}
%&
%1+q_{11}
%&
%\displaystyle
%\frac{q_{12}}{\tau}
%\\[0.25cm]
%0_{1\times t}
%&
%\displaystyle
%\frac{q_{12}}{\tau}
%&
%\displaystyle
%1+\frac{q_{22}}{\tau^2}
%\end{pmatrix}
%+
%\begin{pmatrix}
%\boldsymbol D_t
%&
%0_{t\times t}
%\\
%\eta_1e_{1:t}^{\top}
%&
%0_{1\times t}
%\\
%0_{1\times t}
%&
%\displaystyle
%\frac{\eta_2}{\tau}e_{1:t}^{\top}
%\end{pmatrix}
%\boldsymbol H_t
%\begin{pmatrix}
%\boldsymbol D_t
%&
%0_{t\times t}
%\\
%\eta_1e_{1:t}^{\top}
%&
%0_{1\times t}
%\\
%0_{1\times t}
%&
%\displaystyle
%\frac{\eta_2}{\tau}e_{1:t}^{\top}
%\end{pmatrix}^{\top}.
%\end{aligned}
%\]
%Then,
%\begin{equation}
%\label{eq.stop.loss.biv.ssm}
%\begin{aligned}
%&\E{
%(Y_{t+1}-d)_+
%\mid
%Y_{1:t}=y_{1:t}
%}
%\\
%&\quad=
%\frac{
%\exp\left(
%d_{t+1}
%+
%\frac{\tau^2}{2}
%+
%\frac{q_{22}}{2}
%+
%\eta_2\mu_t^{[2]}
%+
%\frac{\eta_2^2}{2}
%[\boldsymbol H_t]_{2t,2t}
%\right)
%\Phi_{t+2}
%\left(
%a_{1:t+2}^{**}
%;
%\boldsymbol\Omega_{t+2}^{\rm ssm}
%\right)
%-
%d\,
%\Phi_{t+2}
%\left(
%a_{1:t+2}^*
%;
%\boldsymbol\Omega_{t+2}^{\rm ssm}
%\right)
%}
%{
%\Phi_t
%\left(
%a_{1:t}
%+
%\boldsymbol B_t
%\begin{pmatrix}
%\mu_{1:t}^{[1]}\\
%\mu_{1:t}^{[2]}
%\end{pmatrix}
%;
%\boldsymbol I_t
%+
%\boldsymbol B_t
%\boldsymbol H_t
%\boldsymbol B_t^\top
%\right)}
%\end{aligned}
%\end{equation}
%where
%\[
%a_{1:t+2}^*:=
%\begin{pmatrix}
%\displaystyle
%a_{1:t}
%+
%\boldsymbol B_t
%\begin{pmatrix}
%\mu_{1:t}^{[1]}\\
%\mu_{1:t}^{[2]}
%\end{pmatrix}
%\\[0.5cm]
%\displaystyle
%x_{t+1}
%\left(
%c_{t+1}
%+
%\eta_1\mu_t^{[1]}
%\right)
%\\[0.5cm]
%\displaystyle
%\frac{
%d_{t+1}
%-
%\log d
%+
%\eta_2\mu_t^{[2]}
%}{
%\tau
%}
%\end{pmatrix}, \quad 
%a_{1:t+2}^{**}:=
%\begin{pmatrix}
%\displaystyle
%a_{1:t}
%+
%\boldsymbol B_t
%\left(
%\begin{pmatrix}
%\mu_{1:t}^{[1]}\\
%\mu_{1:t}^{[2]}
%\end{pmatrix}
%+
%\eta_2\boldsymbol H_t e_{1:2t}
%\right)
%\\[0.5cm]
%\displaystyle
%c_{t+1}
%+
%\eta_1\mu_t^{[1]}
%+
%q_{12}
%+
%\eta_1\eta_2
%[\boldsymbol H_t]_{t,2t}
%\\[0.5cm]
%\displaystyle
%\frac{
%d_{t+1}
%+
%\tau^2
%+
%q_{22}
%-
%\log d
%+
%\eta_2\mu_t^{[2]}
%+
%\eta_2^2
%[\boldsymbol H_t]_{2t,2t}
%}{
%\tau
%}
%\end{pmatrix}.
%\]
%\end{theorem}
%
%\begin{proof}
%%You do not want to change the following part 
%{We first derive the joint conditional distribution of the latent process
%through time \(t+1\). For this proof, define
%\[
%\boldsymbol A_t
%:=
%\begin{pmatrix}
%\eta_1 e_{1:t}^{\top}
%&
%0_{1\times t}
%\\
%0_{1\times t}
%&
%\eta_2 e_{1:t}^{\top}
%\end{pmatrix}
%\]
%so that 
%\[
%\boldsymbol A_t
%\begin{pmatrix}
%\Theta_{1:t}^{[1]}\\
%\Theta_{1:t}^{[2]}
%\end{pmatrix}
%=
%\begin{pmatrix}
%\eta_1\Theta_t^{[1]}\\
%\eta_2\Theta_t^{[2]}
%\end{pmatrix}.
%\]
%Then, combining part i of Theorem~\ref{thm.biv.ssm} with the one-step
%transition distribution in Model~\ref{model.3}, we obtain
%\begin{equation}
%\label{eq.augmented.posterior.biv.ssm}
%\begin{aligned}
%\left[\left.
%\begin{pmatrix}
%\Theta_{1:t}^{[1]}\\
%\Theta_{1:t}^{[2]}\\
%\Theta_{t+1}^{[1]}\\
%\Theta_{t+1}^{[2]}
%\end{pmatrix}
%\right|
%S_{\mathcal I_t}=s_{\mathcal I_t} \right]
%&=\left[\left.
%\begin{pmatrix}
%\Theta_{1:t}^{[1]}\\
%\Theta_{1:t}^{[2]}\\
%\boldsymbol A_t\begin{pmatrix}
%\Theta_{1:t}^{[1]}\\
%\Theta_{1:t}^{[2]}\\
%\end{pmatrix} + 
%\begin{pmatrix}
%\epsilon_{t+1}^{[1]}\\
%\epsilon_{t+1}^{[2]}\\
%\end{pmatrix}
%\end{pmatrix}
%\right|
%S_{\mathcal I_t}=s_{\mathcal I_t} \right]
%\\
%&\sim
%\operatorname{MVN}_{2t+2}
%\left(
%\begin{pmatrix}
%\mu_{1:t}^{[1]}\\
%\mu_{1:t}^{[2]}\\
%\eta_1\mu_t^{[1]}\\
%\eta_2\mu_t^{[2]}
%\end{pmatrix},
%\begin{pmatrix}
%\boldsymbol H_t
%&
%\boldsymbol H_t\boldsymbol A_t^\top
%\\[0.15cm]
%\boldsymbol A_t\boldsymbol H_t
%&
%\boldsymbol A_t\boldsymbol H_t\boldsymbol A_t^\top
%+
%\begin{pmatrix}
%q_{11}&q_{12}\\
%q_{12}&q_{22}
%\end{pmatrix}
%\end{pmatrix}
%\right).
%\end{aligned}
%\end{equation}
%Note that we have changed the order in \eqref{eq.augmented.posterior.biv.ssm}  to
%keep the block matrix compact.
%
%As in \eqref{eq.conditional.mean.biv.ssm}, we have
%\begin{equation}
%\label{eq.conditional.stop.loss.biv.ssm}
%\begin{aligned}
%&\E{
%(Y_{t+1}-d)_+
%\mid
%Y_{1:t}=y_{1:t}
%}
%\\
%&\quad=
%\frac{
%\E{
%(Y_{t+1}-d)_+
%1_{[Z_{1:t}=z_{1:t}]}
%\mid
%S_{\mathcal I_t}=s_{\mathcal I_t}
%}
%}{
%\P{
%Z_{1:t}=z_{1:t}
%\mid
%S_{\mathcal I_t}=s_{\mathcal I_t}
%}
%}.
%\end{aligned}
%\end{equation}
%The denominator in \eqref{eq.conditional.stop.loss.biv.ssm} is given by
%\eqref{eq.denominator.biv.ssm}.
%
%Now, we calculate the numerator in \eqref{eq.conditional.stop.loss.biv.ssm}. 
%Since $S_{t+1}$ follows the lognormal distribution conditional on
%\(
%(\Theta_{t+1}^{[1]},\Theta_{t+1}^{[2]})
%\), combined with the following observation
%\[
%(Y_{t+1}-d)_+
%=
%Z_{t+1}(S_{t+1}-d)_+,
%\]
%we have
%\[
%\begin{aligned}
%&\E{
%(Y_{t+1}-d)_+
%\mid
%\Theta_{t+1}^{[1]},
%\Theta_{t+1}^{[2]}
%}
%\\
%&\quad=
%\Phi\left(
%c_{t+1}+\Theta_{t+1}^{[1]}
%\right)
%\Bigg[
%\exp\left(
%d_{t+1}
%+
%\Theta_{t+1}^{[2]}
%+
%\frac{\tau^2}{2}
%\right)
%\Phi\left(
%\frac{
%d_{t+1}
%+
%\Theta_{t+1}^{[2]}
%+
%\tau^2
%-
%\log d
%}{
%\tau
%}
%\right)
%\\
%&\hspace{5.5cm}
%-
%d\,
%\Phi\left(
%\frac{
%d_{t+1}
%+
%\Theta_{t+1}^{[2]}
%-
%\log d
%}{
%\tau
%}
%\right)
%\Bigg].
%\end{aligned}
%\]
%As a result, we have
%%Using the preceding conditional stop-loss expectation, the tower
%%property, the conditional independence in Model~\ref{model.3}, and
%%\(x_{t+1}=1\), the numerator in
%\eqref{eq.conditional.stop.loss.biv.ssm} becomes
%\begin{equation}
%\label{eq.numerator.stop.loss.biv.ssm}
%\begin{aligned}
%&\E{
%(Y_{t+1}-d)_+
%1_{[Z_{1:t}=z_{1:t}]}
%\mid
%S_{\mathcal I_t}=s_{\mathcal I_t}
%}
%\\
%&\quad=
%\exp\left(
%d_{t+1}+\frac{\tau^2}{2}
%\right)
%\E{
%\exp\left(
%\Theta_{t+1}^{[2]}
%\right)
%\prod_{j=1}^{t+1}
%\Phi\left(
%x_j(c_j+\Theta_j^{[1]})
%\right)
%\times
%\Phi\left(
%\frac{
%d_{t+1}
%+
%\Theta_{t+1}^{[2]}
%+
%\tau^2
%-
%\log d
%}{
%\tau
%}
%\right)
%\ \bigg|\
%S_{\mathcal I_t}=s_{\mathcal I_t}
%}
%\\
%&\qquad-
%d\,
%\E{
%\prod_{j=1}^{t+1}
%\Phi\left(
%x_j(c_j+\Theta_j^{[1]})
%\right)
%\Phi\left(
%\frac{
%d_{t+1}
%+
%\Theta_{t+1}^{[2]}
%-
%\log d
%}{
%\tau
%}
%\right)
%\ \bigg|\
%S_{\mathcal I_t}=s_{\mathcal I_t}
%}
%\end{aligned}
%\end{equation}
%where we recall $x_{t+1}=1$ by definition. 
%
%
%
%
%We first calculate the expectation in the first term of
%\eqref{eq.numerator.stop.loss.biv.ssm}. Since the product of the probit
%terms is bounded by one, Lemma~\ref{lem.3} 
%with the
%tilting vector
%\[
%\begin{pmatrix}
%0_{1:2t+1}^{\top}\\
%1
%\end{pmatrix}
%\]
%can be applied to get
%\begin{equation}\label{eq.34}
%\begin{aligned}
%&\E{
%\exp\left(
%\Theta_{t+1}^{[2]}
%\right)
%\prod_{j=1}^{t+1}
%\Phi\left(
%x_j(c_j+\Theta_j^{[1]})
%\right)
%\times
%\Phi\left(
%\frac{
%d_{t+1}
%+
%\Theta_{t+1}^{[2]}
%+
%\tau^2
%-
%\log d
%}{
%\tau
%}
%\right)
%\ \bigg|\
%S_{\mathcal I_t}=s_{\mathcal I_t}
%}\\
%&=\exp\left(
%\eta_2\mu_t^{[2]}
%+
%\frac12
%\left\{
%q_{22}
%+
%\eta_2^2
%\left[\boldsymbol H_t\right]_{2t,2t}
%\right\}
%\right)
%%%%%%%%%%%
%\E{
%\prod_{j=1}^{t+1}
%\Phi\left(
%x_j(c_j+\widetilde\Theta_j^{[1]})
%\right)
%\times
%\Phi\left(
%\frac{
%d_{t+1}
%+
%\widetilde\Theta_{t+1}^{[2]}
%+
%\tau^2
%-
%\log d
%}{
%\tau
%}
%\right)
%\ \bigg|\
%S_{\mathcal I_t}=s_{\mathcal I_t}
%}\\
%&=\exp\left(
%\eta_2\mu_t^{[2]}
%+
%\frac12
%\left\{
%q_{22}
%+
%\eta_2^2
%\left[\boldsymbol H_t\right]_{2t,2t}
%\right\}
%\right)
%%%%%%%%%%%
%\E{
%\prod_{j=1}^{t+2}
%\Phi(\widetilde U_j)
%\ \bigg|\
%S_{\mathcal I_t}=s_{\mathcal I_t}
%}.
%\end{aligned}
%\end{equation}
%where
%\begin{equation}\label{eq.42}
%\begin{aligned}
%&\left[\left.
%\begin{pmatrix}
%\widetilde\Theta_{1:t}^{[1]}\\
%\widetilde\Theta_{1:t}^{[2]}\\
%\widetilde\Theta_{t+1}^{[1]}\\
%\widetilde\Theta_{t+1}^{[2]}
%\end{pmatrix}
%\right|
%S_{\mathcal I_t}=s_{\mathcal I_t} \right]\\
%%%%%%%%%%%%
%&\qquad\quad\sim 
%%%%%%%%%%%%%
%\operatorname{MVN}_{2t+2}
%\left(
%\begin{pmatrix}
%\displaystyle
%\begin{pmatrix}
%\mu_{1:t}^{[1]}\\
%\mu_{1:t}^{[2]}
%\end{pmatrix}
%+
%\eta_2\boldsymbol H_t e_{1:2t}
%\\
%\eta_1\mu_t^{[1]}
%+
%q_{12}
%+
%\eta_1\eta_2
%\left[\boldsymbol H_t\right]_{t,2t}
%\\
%\eta_2\mu_t^{[2]}
%+
%q_{22}
%+
%\eta_2^2
%\left[\boldsymbol H_t\right]_{2t,2t}
%\end{pmatrix},
%\begin{pmatrix}
%\boldsymbol H_t
%&
%\boldsymbol H_t\boldsymbol A_t^\top
%\\[0.15cm]
%\boldsymbol A_t\boldsymbol H_t
%&
%\boldsymbol A_t\boldsymbol H_t\boldsymbol A_t^\top
%+
%\begin{pmatrix}
%q_{11}&q_{12}\\
%q_{12}&q_{22}
%\end{pmatrix}
%\end{pmatrix}
%\right)
%\end{aligned}
%\end{equation}
%and
%\[
%\widetilde U_{1:t+2}
%:=
%\begin{pmatrix}
%x_1(c_1+\widetilde\Theta_1^{[1]})\\
%\vdots\\
%x_t(c_t+\widetilde\Theta_t^{[1]})\\
%x_{t+1}(c_{t+1}+\widetilde\Theta_{t+1}^{[1]})\\[0.15cm]
%\displaystyle
%\frac{
%d_{t+1}
%+
%\widetilde\Theta_{t+1}^{[2]}
%+
%\tau^2
%-
%\log d
%}{
%\tau
%}
%\end{pmatrix}.
%\]
%Furthermore, from the following observation
%\[
%\begin{aligned}
%\widetilde U_{1:t+2}
%%&:=
%%\begin{pmatrix}
%%x_1(c_1+\widetilde\Theta_1^{[1]})\\
%%\vdots\\
%%x_t(c_t+\widetilde\Theta_t^{[1]})\\
%%x_{t+1}(c_{t+1}+\widetilde\Theta_{t+1}^{[1]})\\[0.15cm]
%%\displaystyle
%%\frac{
%%d_{t+1}
%%+
%%\widetilde\Theta_{t+1}^{[2]}
%%+
%%\tau^2
%%-
%%\log d
%%}{
%%\tau
%%}
%%\end{pmatrix}\\
%&=\begin{pmatrix}
%x_1c_1\\
%\vdots\\
%x_tc_t\\
%x_{t+1}c_{t+1}\\[0.15cm]
%\displaystyle
%\frac{
%d_{t+1}
%+
%\tau^2
%-
%\log d
%}{
%\tau
%}
%\end{pmatrix} + 
%\begin{pmatrix}
%\boldsymbol D_t
%&
%0_{t\times t}
%&
%0_{t\times1}
%&
%0_{t\times1}
%\\
%0_{1\times t}
%&
%0_{1\times t}
%&
%x_{t+1}
%&
%0
%\\
%0_{1\times t}
%&
%0_{1\times t}
%&
%0
%&
%\displaystyle\frac{1}{\tau}
%\end{pmatrix}
%\begin{pmatrix}
%\widetilde\Theta_{1:t}^{[1]}\\
%\widetilde\Theta_{1:t}^{[2]}\\
%\widetilde\Theta_{t+1}^{[1]}\\
%\widetilde\Theta_{t+1}^{[2]}
%\end{pmatrix}
%\end{aligned},
%\]
%we have
%\[
%\widetilde U_{1:t+2} \mid S_{\mathcal I_t}=s_{\mathcal I_t}
%\]
%follows the multivariate normal with mean 
%\[
%\begin{aligned}
%&\E{
%\widetilde U_{1:t+2}
%\mid
%S_{\mathcal I_t}=s_{\mathcal I_t}
%}
%\\
%&\quad=
%\begin{pmatrix}
%\displaystyle
%a_{1:t}
%+
%\boldsymbol B_t
%\left(
%\begin{pmatrix}
%\mu_{1:t}^{[1]}\\
%\mu_{1:t}^{[2]}
%\end{pmatrix}
%+
%\eta_2\boldsymbol H_t e_{1:2t}
%\right)
%\\[0.5cm]
%\displaystyle
%x_{t+1}
%\left(
%c_{t+1}
%+
%\eta_1\mu_t^{[1]}
%+
%q_{12}
%+
%\eta_1\eta_2
%[\boldsymbol H_t]_{t,2t}
%\right)
%\\[0.5cm]
%\displaystyle
%\frac{
%d_{t+1}
%+
%\tau^2
%-
%\log d
%+
%\eta_2\mu_t^{[2]}
%+
%q_{22}
%+
%\eta_2^2
%[\boldsymbol H_t]_{2t,2t}
%}{
%\tau
%}
%\end{pmatrix}=a_{1:t+2}^{**}.
%\end{aligned}
%\]
%}
%and covariance matrix 
%%\[
%%\begin{pmatrix}
%%\boldsymbol H_t
%%&
%%\boldsymbol H_t\boldsymbol A_t^\top
%%\\[0.15cm]
%%\boldsymbol A_t\boldsymbol H_t
%%&
%%\boldsymbol A_t\boldsymbol H_t\boldsymbol A_t^\top
%%+
%%\begin{pmatrix}
%%q_{11}&q_{12}\\
%%q_{12}&q_{22}
%%\end{pmatrix}
%%\end{pmatrix}.
%%\]
%%Since additive constants do not affect covariance, the definition of
%\[
%\begin{aligned}
%&\Var{
%\widetilde U_{1:t+2}
%\mid
%S_{\mathcal I_t}=s_{\mathcal I_t}
%}
%\\
%&\quad=
%\begin{pmatrix}
%\boldsymbol D_t
%&
%0_{t\times t}
%&
%0_{t\times1}
%&
%0_{t\times1}
%\\
%0_{1\times t}
%&
%0_{1\times t}
%&
%x_{t+1}
%&
%0
%\\
%0_{1\times t}
%&
%0_{1\times t}
%&
%0
%&
%\displaystyle\frac{1}{\tau}
%\end{pmatrix}
%\begin{pmatrix}
%\boldsymbol H_t
%&
%\boldsymbol H_t\boldsymbol A_t^\top
%\\[0.15cm]
%\boldsymbol A_t\boldsymbol H_t
%&
%\boldsymbol A_t\boldsymbol H_t\boldsymbol A_t^\top
%+
%\begin{pmatrix}
%q_{11}&q_{12}\\
%q_{12}&q_{22}
%\end{pmatrix}
%\end{pmatrix}
%\begin{pmatrix}
%\boldsymbol D_t
%&
%0_{t\times t}
%&
%0_{t\times1}
%&
%0_{t\times1}
%\\
%0_{1\times t}
%&
%0_{1\times t}
%&
%x_{t+1}
%&
%0
%\\
%0_{1\times t}
%&
%0_{1\times t}
%&
%0
%&
%\displaystyle\frac{1}{\tau}
%\end{pmatrix}^{\top}.
%\end{aligned}
%\]
%Furthermore, pluging-in the following observation into the second term of above, 
%\[
%\begin{aligned}
%&
%\begin{pmatrix}
%\boldsymbol H_t
%&
%\boldsymbol H_t\boldsymbol A_t^\top
%\\
%\boldsymbol A_t\boldsymbol H_t
%&
%\boldsymbol A_t\boldsymbol H_t\boldsymbol A_t^\top
%+
%\begin{pmatrix}
%q_{11}&q_{12}\\
%q_{12}&q_{22}
%\end{pmatrix}
%\end{pmatrix}
%\\
%&\quad=
%\begin{pmatrix}
%\boldsymbol I_{2t}\\
%\boldsymbol A_t
%\end{pmatrix}
%\boldsymbol H_t
%\begin{pmatrix}
%\boldsymbol I_{2t}
%&
%\boldsymbol A_t^\top
%\end{pmatrix}
%+
%\begin{pmatrix}
%0_{2t\times2t}
%&
%0_{2t\times2}
%\\
%0_{2\times2t}
%&
%\begin{pmatrix}
%q_{11}&q_{12}\\
%q_{12}&q_{22}
%\end{pmatrix}
%\end{pmatrix}
%\end{aligned}
%\]
%we have
%\[
%\begin{aligned}
%&\Var{
%\widetilde U_{1:t+2}
%\mid
%S_{\mathcal I_t}=s_{\mathcal I_t}
%}
%\\
%&\quad=
%\begin{pmatrix}
%0_{t\times t}
%&
%0_{t\times1}
%&
%0_{t\times1}
%\\
%0_{1\times t}
%&
%x_{t+1}^2q_{11}
%&
%\displaystyle
%\frac{x_{t+1}q_{12}}{\tau}
%\\[0.25cm]
%0_{1\times t}
%&
%\displaystyle
%\frac{x_{t+1}q_{12}}{\tau}
%&
%\displaystyle
%\frac{q_{22}}{\tau^2}
%\end{pmatrix}
%+
%\begin{pmatrix}
%\boldsymbol D_t
%&
%0_{t\times t}
%\\
%x_{t+1}\eta_1e_{1:t}^{\top}
%&
%0_{1\times t}
%\\
%0_{1\times t}
%&
%\displaystyle
%\frac{\eta_2}{\tau}e_{1:t}^{\top}
%\end{pmatrix}
%\boldsymbol H_t
%\begin{pmatrix}
%\boldsymbol D_t
%&
%0_{t\times t}
%\\
%x_{t+1}\eta_1e_{1:t}^{\top}
%&
%0_{1\times t}
%\\
%0_{1\times t}
%&
%\displaystyle
%\frac{\eta_2}{\tau}e_{1:t}^{\top}
%\end{pmatrix}^{\top}\\
%%%%%%%%%%%
%%%%%%%%%%
%&\quad=
%\begin{pmatrix}
%0_{t\times t}
%&
%0_{t\times1}
%&
%0_{t\times1}
%\\
%0_{1\times t}
%&
%q_{11}
%&
%\displaystyle
%\frac{q_{12}}{\tau}
%\\[0.25cm]
%0_{1\times t}
%&
%\displaystyle
%\frac{q_{12}}{\tau}
%&
%\displaystyle
%\frac{q_{22}}{\tau^2}
%\end{pmatrix}
%+
%\begin{pmatrix}
%\boldsymbol D_t
%&
%0_{t\times t}
%\\
%\eta_1e_{1:t}^{\top}
%&
%0_{1\times t}
%\\
%0_{1\times t}
%&
%\displaystyle
%\frac{\eta_2}{\tau}e_{1:t}^{\top}
%\end{pmatrix}
%\boldsymbol H_t
%\begin{pmatrix}
%\boldsymbol D_t
%&
%0_{t\times t}
%\\
%\eta_1e_{1:t}^{\top}
%&
%0_{1\times t}
%\\
%0_{1\times t}
%&
%\displaystyle
%\frac{\eta_2}{\tau}e_{1:t}^{\top}
%\end{pmatrix}^{\top},
%\end{aligned}
%\]
%where we have used \(x_{t+1}=1\).
%It now follows directly from Lemma~\ref{lem.5} that
%\[
%\begin{aligned}
%\E{
%\prod_{j=1}^{t+2}
%\Phi(\widetilde U_j)
%\ \bigg|\
%S_{\mathcal I_t}=s_{\mathcal I_t}
%}
%&=
%\Phi_{t+2}
%\left(
%a_{1:t+2}^{**}
%;
%\boldsymbol\Omega_{t+2}^{\rm ssm}
%\right)
%\end{aligned}
%\]
%where we have used
%\[
%\begin{aligned}
%\boldsymbol\Omega_{t+2}^{\rm ssm}
%&=\boldsymbol I_{t+2}
%+
%\Var{
%\widetilde U_{1:t+2}
%\mid
%S_{\mathcal I_t}=s_{\mathcal I_t}
%}\\
%&=
%\begin{pmatrix}
%\boldsymbol I_t
%&
%0_{t\times1}
%&
%0_{t\times1}
%\\
%0_{1\times t}
%&
%1+q_{11}
%&
%\displaystyle
%\frac{q_{12}}{\tau}
%\\[0.25cm]
%0_{1\times t}
%&
%\displaystyle
%\frac{q_{12}}{\tau}
%&
%\displaystyle
%1+\frac{q_{22}}{\tau^2}
%\end{pmatrix}
%+
%\begin{pmatrix}
%\boldsymbol D_t
%&
%0_{t\times t}
%\\
%\eta_1e_{1:t}^{\top}
%&
%0_{1\times t}
%\\
%0_{1\times t}
%&
%\displaystyle
%\frac{\eta_2}{\tau}e_{1:t}^{\top}
%\end{pmatrix}
%\boldsymbol H_t
%\begin{pmatrix}
%\boldsymbol D_t
%&
%0_{t\times t}
%\\
%\eta_1e_{1:t}^{\top}
%&
%0_{1\times t}
%\\
%0_{1\times t}
%&
%\displaystyle
%\frac{\eta_2}{\tau}e_{1:t}^{\top}
%\end{pmatrix}^{\top}.
%\\
%\end{aligned}
%\]
%
%
%Combining this equality with \eqref{eq.34}, we have
%\begin{equation}
%\label{eq.first.stop.loss.term.biv.ssm}
%\begin{aligned}
%&\E{
%\exp\left(
%\Theta_{t+1}^{[2]}
%\right)
%\prod_{j=1}^{t+1}
%\Phi\left(
%x_j(c_j+\Theta_j^{[1]})
%\right)
%\Phi\left(
%\frac{
%d_{t+1}
%+
%\Theta_{t+1}^{[2]}
%+
%\tau^2
%-
%\log d
%}{
%\tau
%}
%\right)
%\ \bigg|\
%S_{\mathcal I_t}=s_{\mathcal I_t}
%}
%\\
%&\quad=
%\exp\left(
%\eta_2\mu_t^{[2]}
%+
%\frac{q_{22}}{2}
%+
%\frac{\eta_2^2}{2}
%[\boldsymbol H_t]_{2t,2t}
%\right)
%\Phi_{t+2}
%\left(
%a_{1:t+2}^{**}
%;
%\boldsymbol\Omega_{t+2}^{\rm ssm}
%\right).
%\end{aligned}
%\end{equation}
%
%We next calculate the expectation in the second term of
%\eqref{eq.numerator.stop.loss.biv.ssm}. From \eqref{eq.augmented.posterior.biv.ssm}, conditional on $S_{\mathcal I_t}=s_{\mathcal I_t}$, the following random vector
%\[
%\begin{pmatrix}
%x_1(c_1+\Theta_1^{[1]})\\
%\vdots\\
%x_t(c_t+\Theta_t^{[1]})\\
%x_{t+1}(c_{t+1}+\Theta_{t+1}^{[1]})\\[0.15cm]
%\displaystyle
%\frac{
%d_{t+1}
%+
%\Theta_{t+1}^{[2]}
%-
%\log d
%}{
%\tau
%}
%\end{pmatrix}
%\]
%follows the multivariate normal with mean $a_{1:t+2}^*$
%and the same covariance matrix as that of
%\(\widetilde U_{1:t+2}\). 
%Note that these two covariances are the same 
%because exponential tilting in \eqref{eq.42} changes only the
%mean and do not affect the covariance. Therefore, Lemma~\ref{lem.5} gives
%\begin{equation}
%\label{eq.second.stop.loss.term.biv.ssm}
%\begin{aligned}
%&\E{
%\prod_{j=1}^{t+1}
%\Phi\left(
%x_j(c_j+\Theta_j^{[1]})
%\right)
%\Phi\left(
%\frac{
%d_{t+1}
%+
%\Theta_{t+1}^{[2]}
%-
%\log d
%}{
%\tau
%}
%\right)
%\ \bigg|\
%S_{\mathcal I_t}=s_{\mathcal I_t}
%}
%=
%\Phi_{t+2}
%\left(
%a_{1:t+2}^*
%;
%\boldsymbol\Omega_{t+2}^{\rm ssm}
%\right).
%\end{aligned}
%\end{equation}
%
%Substituting \eqref{eq.first.stop.loss.term.biv.ssm} and
%\eqref{eq.second.stop.loss.term.biv.ssm} into
%\eqref{eq.numerator.stop.loss.biv.ssm}, and then dividing by the
%denominator in \eqref{eq.denominator.biv.ssm}, proves
%\eqref{eq.stop.loss.biv.ssm}.
%
%\end{proof}
%
%
%\begin{corollary}[Closed-form likelihood]
%\label{cor2.1.ssm}
%Under the same setting as in Theorem~\ref{thm.biv.ssm}, the likelihood
%in \eqref{eq.13} under Model~\ref{model.3} is given by
%\begin{equation}
%\label{eq.likelihood.biv.ssm}
%\begin{aligned}
%\ell(y_{1:t})
%&=
%\Phi_t
%\left(
%a_{1:t}
%+
%\boldsymbol B_t
%\begin{pmatrix}
%\mu_{1:t}^{[1]}\\
%\mu_{1:t}^{[2]}
%\end{pmatrix}
%;
%\boldsymbol I_t
%+
%\boldsymbol B_t
%\boldsymbol H_t
%\boldsymbol B_t^\top
%\right)
%\frac{
%\det(\boldsymbol H_t)^{1/2}
%\exp\left(
%-\frac{1}{2\tau^2}
%r_{1:t}^{\top}
%\boldsymbol M_t
%\left(
%r_{1:t}-\mu_{1:t}^{[2]}
%\right)
%\right)
%}{
%\det(\boldsymbol\Gamma_t)^{1/2}
%\tau^{n_t}
%(2\pi)^{n_t/2}
%\displaystyle\prod_{j\in\mathcal I_t}s_j
%}
%\\
%\end{aligned}
%\end{equation}
%\end{corollary}
%
%\begin{proof}
%From \eqref{eq.13}, it is enough to calculate
%\[
%\P{
%Z_{1:t}=z_{1:t}
%\mid
%S_{\mathcal I_t}=s_{\mathcal I_t}
%}
%\quad\hbox{and}\quad
%f(s_{\mathcal I_t}).
%\]
%First, by \eqref{eq.denominator.biv.ssm}, we have
%\[
%\begin{aligned}
%&\P{
%Z_{1:t}=z_{1:t}
%\mid
%S_{\mathcal I_t}=s_{\mathcal I_t}
%}
%\\
%&\quad=
%\Phi_t
%\left(
%a_{1:t}
%+
%\boldsymbol B_t
%\begin{pmatrix}
%\mu_{1:t}^{[1]}\\
%\mu_{1:t}^{[2]}
%\end{pmatrix}
%;
%\boldsymbol I_t
%+
%\boldsymbol B_t
%\boldsymbol H_t
%\boldsymbol B_t^\top
%\right).
%\end{aligned}
%\]
%
%It remains to calculate \(f(s_{\mathcal I_t})\). Conditional on the
%latent process, the observed severities have joint density
%\[
%\begin{aligned}
%&f\left(
%s_{\mathcal I_t}
%\ \middle|\
%\begin{pmatrix}
%\Theta_{1:t}^{[1]}\\
%\Theta_{1:t}^{[2]}
%\end{pmatrix}
%=
%\begin{pmatrix}
%\theta_{1:t}^{[1]}\\
%\theta_{1:t}^{[2]}
%\end{pmatrix}
%\right)
%\\
%&\quad=
%\frac{1}{
%\tau^{n_t}
%(2\pi)^{n_t/2}
%\displaystyle\prod_{j\in\mathcal I_t}s_j
%}
%\exp\left\{
%-\frac{1}{2\tau^2}
%\left(
%r_{1:t}-\theta_{1:t}^{[2]}
%\right)^\top
%\boldsymbol M_t
%\left(
%r_{1:t}-\theta_{1:t}^{[2]}
%\right)
%\right\}.
%\end{aligned}
%\]
%Consequently,
%\[
%\begin{aligned}
%f(s_{\mathcal I_t})
%&=
%\frac{1}{
%\tau^{n_t}
%(2\pi)^{n_t/2}
%\displaystyle\prod_{j\in\mathcal I_t}s_j
%}
%\\
%&\quad\times
%\int_{\Real^{2t}}
%\phi_{2t}
%\left(
%\begin{pmatrix}
%\theta_{1:t}^{[1]}\\
%\theta_{1:t}^{[2]}
%\end{pmatrix}
%;
%\boldsymbol\Gamma_t
%\right)
%\\
%&\qquad\qquad\times
%\exp\left\{
%-\frac{1}{2\tau^2}
%\left(
%r_{1:t}-\theta_{1:t}^{[2]}
%\right)^\top
%\boldsymbol M_t
%\left(
%r_{1:t}-\theta_{1:t}^{[2]}
%\right)
%\right\}
%\,{\rm d}\theta_{1:t}^{[1]}
%\,{\rm d}\theta_{1:t}^{[2]}.
%\end{aligned}
%\]
%
%By the definitions of \(\boldsymbol H_t\) and
%\(
%\left(
%(\mu_{1:t}^{[1]})^\top,
%(\mu_{1:t}^{[2]})^\top
%\right)^\top
%\),
%the exponent in the integrand can be completed as
%\[
%\begin{aligned}
%&-\frac12
%\begin{pmatrix}
%\theta_{1:t}^{[1]}\\
%\theta_{1:t}^{[2]}
%\end{pmatrix}^{\top}
%\boldsymbol\Gamma_t^{-1}
%\begin{pmatrix}
%\theta_{1:t}^{[1]}\\
%\theta_{1:t}^{[2]}
%\end{pmatrix}
%\\
%&\quad
%-\frac{1}{2\tau^2}
%\left(
%r_{1:t}-\theta_{1:t}^{[2]}
%\right)^\top
%\boldsymbol M_t
%\left(
%r_{1:t}-\theta_{1:t}^{[2]}
%\right)
%\\
%&=
%-\frac12
%\left[
%\begin{pmatrix}
%\theta_{1:t}^{[1]}\\
%\theta_{1:t}^{[2]}
%\end{pmatrix}
%-
%\begin{pmatrix}
%\mu_{1:t}^{[1]}\\
%\mu_{1:t}^{[2]}
%\end{pmatrix}
%\right]^\top
%\boldsymbol H_t^{-1}
%\left[
%\begin{pmatrix}
%\theta_{1:t}^{[1]}\\
%\theta_{1:t}^{[2]}
%\end{pmatrix}
%-
%\begin{pmatrix}
%\mu_{1:t}^{[1]}\\
%\mu_{1:t}^{[2]}
%\end{pmatrix}
%\right]
%\\
%&\quad
%-\frac{1}{2\tau^2}
%r_{1:t}^{\top}
%\boldsymbol M_t
%r_{1:t}
%+
%\frac12
%\begin{pmatrix}
%\mu_{1:t}^{[1]}\\
%\mu_{1:t}^{[2]}
%\end{pmatrix}^{\top}
%\boldsymbol H_t^{-1}
%\begin{pmatrix}
%\mu_{1:t}^{[1]}\\
%\mu_{1:t}^{[2]}
%\end{pmatrix}.
%\end{aligned}
%\]
%Moreover, the definition of the posterior mean gives
%\[
%\boldsymbol H_t^{-1}
%\begin{pmatrix}
%\mu_{1:t}^{[1]}\\
%\mu_{1:t}^{[2]}
%\end{pmatrix}
%=
%\frac{1}{\tau^2}
%\begin{pmatrix}
%0_{1:t}^{\top}\\
%\boldsymbol M_t r_{1:t}
%\end{pmatrix},
%\]
%and hence
%\[
%\begin{aligned}
%\begin{pmatrix}
%\mu_{1:t}^{[1]}\\
%\mu_{1:t}^{[2]}
%\end{pmatrix}^{\top}
%\boldsymbol H_t^{-1}
%\begin{pmatrix}
%\mu_{1:t}^{[1]}\\
%\mu_{1:t}^{[2]}
%\end{pmatrix}
%&=
%\frac{1}{\tau^2}
%(\mu_{1:t}^{[2]})^\top
%\boldsymbol M_t r_{1:t}.
%\end{aligned}
%\]
%Integrating the resulting \(2t\)-dimensional Gaussian kernel therefore
%gives
%\[
%\begin{aligned}
%f(s_{\mathcal I_t})
%&=
%\frac{
%\det(\boldsymbol H_t)^{1/2}
%}{
%\det(\boldsymbol\Gamma_t)^{1/2}
%\tau^{n_t}
%(2\pi)^{n_t/2}
%\displaystyle\prod_{j\in\mathcal I_t}s_j
%}
%\\
%&\quad\times
%\exp\left\{
%-\frac{1}{2\tau^2}
%r_{1:t}^{\top}
%\boldsymbol M_t
%\left(
%r_{1:t}-\mu_{1:t}^{[2]}
%\right)
%\right\}.
%\end{aligned}
%\]
%Multiplying this expression by the conditional probability obtained
%above proves \eqref{eq.likelihood.biv.ssm}.
%
%When \(\mathcal I_t=\varnothing\), we have
%\(\boldsymbol M_t=0_{t\times t}\),
%\(\boldsymbol H_t=\boldsymbol\Gamma_t\), and
%\[
%\begin{pmatrix}
%\mu_{1:t}^{[1]}\\
%\mu_{1:t}^{[2]}
%\end{pmatrix}
%=
%0_{1:2t}^{\top},
%\]
%so that the marginal-severity factor is equal to \(1\), as required.
%\end{proof}
%
%\begin{corollary}[Filtering and one-step predictive distributions of the latent state]
%\label{cor.posterior.biv.ssm}
%Under the setting of Theorem~\ref{thm.biv.ssm}, the following results
%hold.
%
%\begin{itemize}
%
%\item[i.] The filtering distribution
%\[
%\begin{pmatrix}
%\Theta_t^{[1]}\\
%\Theta_t^{[2]}
%\end{pmatrix}
%\ \bigg|\ 
%Y_{1:t}=y_{1:t}
%\]
%has density
%\[
%\pi\left(
%\theta_t^{[1]},\theta_t^{[2]}
%\mid
%Y_{1:t}=y_{1:t}
%\right),
%\qquad
%\begin{pmatrix}
%\theta_t^{[1]}\\
%\theta_t^{[2]}
%\end{pmatrix}
%\in\Real^2,
%\]
%given by
%\begin{equation}
%\label{eq.filtering.state.biv.ssm}
%\begin{aligned}
%&\pi\left(
%\theta_t^{[1]},\theta_t^{[2]}
%\mid
%Y_{1:t}=y_{1:t}
%\right)
%\\
%&\quad=
%\frac{
%\displaystyle
%\phi_2
%\left(
%\begin{pmatrix}
%\theta_t^{[1]}-\mu_t^{[1]}\\[0.1cm]
%\theta_t^{[2]}-\mu_t^{[2]}
%\end{pmatrix};
%\begin{pmatrix}
%[\boldsymbol H_t]_{t,t}
%&
%[\boldsymbol H_t]_{t,2t}
%\\
%[\boldsymbol H_t]_{t,2t}
%&
%[\boldsymbol H_t]_{2t,2t}
%\end{pmatrix}
%\right)
%}{
%\displaystyle
%\Phi_t
%\left(
%a_{1:t}
%+
%\boldsymbol B_t
%\begin{pmatrix}
%\mu_{1:t}^{[1]}\\
%\mu_{1:t}^{[2]}
%\end{pmatrix};
%\boldsymbol I_t
%+
%\boldsymbol B_t
%\boldsymbol H_t
%\boldsymbol B_t^\top
%\right)
%}
%\\
%&\qquad\times
%\Phi_t
%\Bigg(
%a_{1:t}
%+
%\boldsymbol B_t
%\Bigg[
%\begin{pmatrix}
%\mu_{1:t}^{[1]}\\
%\mu_{1:t}^{[2]}
%\end{pmatrix}
%+
%\boldsymbol H_t
%\begin{pmatrix}
%e_{1:t}&0_{t\times1}\\
%0_{t\times1}&e_{1:t}
%\end{pmatrix}
%\begin{pmatrix}
%[\boldsymbol H_t]_{t,t}
%&
%[\boldsymbol H_t]_{t,2t}
%\\
%[\boldsymbol H_t]_{t,2t}
%&
%[\boldsymbol H_t]_{2t,2t}
%\end{pmatrix}^{-1}
%\begin{pmatrix}
%\theta_t^{[1]}-\mu_t^{[1]}\\[0.1cm]
%\theta_t^{[2]}-\mu_t^{[2]}
%\end{pmatrix}
%\Bigg]
%;
%\\
%&\hspace{1.5cm}
%\boldsymbol I_t
%+
%\boldsymbol B_t
%\Bigg[
%\boldsymbol H_t
%-
%\boldsymbol H_t
%\begin{pmatrix}
%e_{1:t}&0_{t\times1}\\
%0_{t\times1}&e_{1:t}
%\end{pmatrix}
%\begin{pmatrix}
%[\boldsymbol H_t]_{t,t}
%&
%[\boldsymbol H_t]_{t,2t}
%\\
%[\boldsymbol H_t]_{t,2t}
%&
%[\boldsymbol H_t]_{2t,2t}
%\end{pmatrix}^{-1}
%\begin{pmatrix}
%e_{1:t}&0_{t\times1}\\
%0_{t\times1}&e_{1:t}
%\end{pmatrix}^{\top}
%\boldsymbol H_t
%\Bigg]
%\boldsymbol B_t^\top
%\Bigg).
%\end{aligned}
%\end{equation}
%
%\item[ii.] The one-step predictive distribution
%\[
%\begin{pmatrix}
%\Theta_{t+1}^{[1]}\\
%\Theta_{t+1}^{[2]}
%\end{pmatrix}
%\ \bigg|\ 
%Y_{1:t}=y_{1:t}
%\]
%has density
%\[
%\pi\left(
%\theta_{t+1}^{[1]},\theta_{t+1}^{[2]}
%\mid
%Y_{1:t}=y_{1:t}
%\right),
%\qquad
%\begin{pmatrix}
%\theta_{t+1}^{[1]}\\
%\theta_{t+1}^{[2]}
%\end{pmatrix}
%\in\Real^2,
%\]
%given by
%\begin{equation}
%\label{eq.predictive.state.biv.ssm}
%\begin{aligned}
%&\pi\left(
%\theta_{t+1}^{[1]},\theta_{t+1}^{[2]}
%\mid
%Y_{1:t}=y_{1:t}
%\right)
%\\
%&\quad=
%\frac{
%\displaystyle
%\phi_2
%\left(
%\begin{pmatrix}
%\theta_{t+1}^{[1]}-\eta_1\mu_t^{[1]}\\[0.1cm]
%\theta_{t+1}^{[2]}-\eta_2\mu_t^{[2]}
%\end{pmatrix};
%\begin{pmatrix}
%q_{11}
%+
%\eta_1^2[\boldsymbol H_t]_{t,t}
%&
%q_{12}
%+
%\eta_1\eta_2[\boldsymbol H_t]_{t,2t}
%\\
%q_{12}
%+
%\eta_1\eta_2[\boldsymbol H_t]_{t,2t}
%&
%q_{22}
%+
%\eta_2^2[\boldsymbol H_t]_{2t,2t}
%\end{pmatrix}
%\right)
%}{
%\displaystyle
%\Phi_t
%\left(
%a_{1:t}
%+
%\boldsymbol B_t
%\begin{pmatrix}
%\mu_{1:t}^{[1]}\\
%\mu_{1:t}^{[2]}
%\end{pmatrix};
%\boldsymbol I_t
%+
%\boldsymbol B_t
%\boldsymbol H_t
%\boldsymbol B_t^\top
%\right)
%}
%\\
%&\qquad\times
%\Phi_t
%\Bigg(
%a_{1:t}
%+
%\boldsymbol B_t
%\Bigg[
%\begin{pmatrix}
%\mu_{1:t}^{[1]}\\
%\mu_{1:t}^{[2]}
%\end{pmatrix}
%+
%\boldsymbol H_t
%\begin{pmatrix}
%\eta_1e_{1:t}&0_{t\times1}\\
%0_{t\times1}&\eta_2e_{1:t}
%\end{pmatrix}
%\\
%&\hspace{4.5cm}\times
%\begin{pmatrix}
%q_{11}
%+
%\eta_1^2[\boldsymbol H_t]_{t,t}
%&
%q_{12}
%+
%\eta_1\eta_2[\boldsymbol H_t]_{t,2t}
%\\
%q_{12}
%+
%\eta_1\eta_2[\boldsymbol H_t]_{t,2t}
%&
%q_{22}
%+
%\eta_2^2[\boldsymbol H_t]_{2t,2t}
%\end{pmatrix}^{-1}
%\begin{pmatrix}
%\theta_{t+1}^{[1]}-\eta_1\mu_t^{[1]}\\[0.1cm]
%\theta_{t+1}^{[2]}-\eta_2\mu_t^{[2]}
%\end{pmatrix}
%\Bigg]
%;
%\\
%&\hspace{1.5cm}
%\boldsymbol I_t
%+
%\boldsymbol B_t
%\Bigg[
%\boldsymbol H_t
%-
%\boldsymbol H_t
%\begin{pmatrix}
%\eta_1e_{1:t}&0_{t\times1}\\
%0_{t\times1}&\eta_2e_{1:t}
%\end{pmatrix}
%\\
%&\hspace{3.0cm}\times
%\begin{pmatrix}
%q_{11}
%+
%\eta_1^2[\boldsymbol H_t]_{t,t}
%&
%q_{12}
%+
%\eta_1\eta_2[\boldsymbol H_t]_{t,2t}
%\\
%q_{12}
%+
%\eta_1\eta_2[\boldsymbol H_t]_{t,2t}
%&
%q_{22}
%+
%\eta_2^2[\boldsymbol H_t]_{2t,2t}
%\end{pmatrix}^{-1}
%\begin{pmatrix}
%\eta_1e_{1:t}&0_{t\times1}\\
%0_{t\times1}&\eta_2e_{1:t}
%\end{pmatrix}^{\top}
%\boldsymbol H_t
%\Bigg]
%\boldsymbol B_t^\top
%\Bigg).
%\end{aligned}
%\end{equation}
%
%\end{itemize}
%\end{corollary}
%
%\begin{proof}
%By part i of Theorem~\ref{thm.biv.ssm},
%\[
%\begin{pmatrix}
%\Theta_{1:t}^{[1]}\\
%\Theta_{1:t}^{[2]}
%\end{pmatrix}
%\ \bigg|\ 
%S_{\mathcal I_t}=s_{\mathcal I_t}
%\sim
%\operatorname{MVN}_{2t}
%\left(
%\begin{pmatrix}
%\mu_{1:t}^{[1]}\\
%\mu_{1:t}^{[2]}
%\end{pmatrix},
%\boldsymbol H_t
%\right).
%\]
%
%The conditional distribution of
%\(
%(\Theta_t^{[1]},\Theta_t^{[2]})^\top
%\)
%given \(S_{\mathcal I_t}=s_{\mathcal I_t}\) is therefore Gaussian with
%mean
%\[
%\begin{pmatrix}
%\mu_t^{[1]}\\
%\mu_t^{[2]}
%\end{pmatrix}
%\]
%and covariance matrix
%\[
%\begin{pmatrix}
%[\boldsymbol H_t]_{t,t}
%&
%[\boldsymbol H_t]_{t,2t}
%\\
%[\boldsymbol H_t]_{t,2t}
%&
%[\boldsymbol H_t]_{2t,2t}
%\end{pmatrix}.
%\]
%The Gaussian conditioning formula gives the conditional mean and
%covariance of
%\[
%\begin{pmatrix}
%\Theta_{1:t}^{[1]}\\
%\Theta_{1:t}^{[2]}
%\end{pmatrix}
%\]
%given
%\(
%(\Theta_t^{[1]},\Theta_t^{[2]})^\top
%\)
%that appear inside the final Gaussian cumulative distribution function
%in \eqref{eq.filtering.state.biv.ssm}. Consequently, conditional
%independence in Model~\ref{model.3} and Lemma~\ref{lem.5} give the
%corresponding conditional probability of
%\(Z_{1:t}=z_{1:t}\). Bayes' formula, together with
%\eqref{eq.denominator.biv.ssm}, proves part i.
%
%For part ii, the evolution equation in Model~\ref{model.3} and part i of
%Theorem~\ref{thm.biv.ssm} give
%\[
%\left[
%\begin{pmatrix}
%\Theta_{t+1}^{[1]}\\
%\Theta_{t+1}^{[2]}
%\end{pmatrix}
%\ \middle|\
%S_{\mathcal I_t}=s_{\mathcal I_t}
%\right]
%\sim
%\operatorname{MVN}_2
%\left(
%\begin{pmatrix}
%\eta_1\mu_t^{[1]}\\
%\eta_2\mu_t^{[2]}
%\end{pmatrix},
%\begin{pmatrix}
%q_{11}
%+
%\eta_1^2[\boldsymbol H_t]_{t,t}
%&
%q_{12}
%+
%\eta_1\eta_2[\boldsymbol H_t]_{t,2t}
%\\
%q_{12}
%+
%\eta_1\eta_2[\boldsymbol H_t]_{t,2t}
%&
%q_{22}
%+
%\eta_2^2[\boldsymbol H_t]_{2t,2t}
%\end{pmatrix}
%\right).
%\]
%Moreover, the conditional cross-covariance between
%\[
%\begin{pmatrix}
%\Theta_{1:t}^{[1]}\\
%\Theta_{1:t}^{[2]}
%\end{pmatrix}
%\quad\hbox{and}\quad
%\begin{pmatrix}
%\Theta_{t+1}^{[1]}\\
%\Theta_{t+1}^{[2]}
%\end{pmatrix}
%\]
%given \(S_{\mathcal I_t}=s_{\mathcal I_t}\) is
%\[
%\boldsymbol H_t
%\begin{pmatrix}
%\eta_1e_{1:t}&0_{t\times1}\\
%0_{t\times1}&\eta_2e_{1:t}
%\end{pmatrix}.
%\]
%The Gaussian conditioning formula therefore gives the conditional mean
%and covariance appearing in the final Gaussian cumulative distribution
%function in \eqref{eq.predictive.state.biv.ssm}. Applying
%Lemma~\ref{lem.5} and Bayes' formula, and again using
%\eqref{eq.denominator.biv.ssm}, proves part ii.
%\end{proof}
%
%
%


%==========================================================================
% Section: Poisson-gamma / binomial-probit thinning model
%          for bivariate claim frequency
%==========================================================================
\section{Binomial thinning model for bivariate claim frequency}
\label{sec.thin}

In this section, we develop a model for two nested claim counts.
For each region \(i\) and period \(t\), let \(N_{it}\) denote the total
number of claims and let \(M_{it}\) denote the number of extreme claims
among them, for example, the claims whose amounts exceed a given
threshold. By construction,
\[
0\le M_{it}\le N_{it}.
\]
A natural way to respect this constraint is to regard \(M_{it}\) as a
binomial thinning of \(N_{it}\): each of the \(N_{it}\) claims is
classified as extreme or not. Repeated observations of the same region
over time are linked through two region-specific random effects, one
for the claim frequency and one for the probability that a claim is
extreme. Throughout this section, the two random effects are assumed to
be independent and static over time, and predictions are obtained as
exact Bayesian posterior means.

\begin{model}[Poisson-gamma binomial-probit thinning model]
\label{model.thin}
Let
\[
r,\sigma\in\Real_{>0},
\qquad
\alpha,\beta\in\Real^p,
\]
and let \(x_{it}\in\Real^p\), \(i=1,\ldots,I\), \(t\in\mathbb N\), be
predetermined exogenous covariates, which are allowed to vary over time.
Define
\[
\lambda_{it}:=\exp\left(x_{it}^\top\alpha\right)\in\Real_{>0},
\qquad
\eta_{it}:=x_{it}^\top\beta\in\Real.
\]
We define the bivariate sequences
\(\left((N_{it},M_{it})\right)_{t\in\mathbb N}\), \(i=1,\ldots,I\),
through the following random effect assumption and observation
equations.
\begin{itemize}
  \item[i.] Random effects: For each region \(i\), assume
  \[
  \Theta_i^{(1)}\sim\operatorname{Gamma}(r,r),
  \qquad
  \Theta_i^{(2)}\sim\operatorname{N}(0,\sigma^2),
  \]
  where \(\Theta_i^{(1)}\) and \(\Theta_i^{(2)}\) are independent.
  Furthermore, the pairs
  \(\left(\Theta_i^{(1)},\Theta_i^{(2)}\right)\), \(i=1,\ldots,I\),
  are mutually independent.

  \item[ii.] Observation equations: Conditional on
  \(\left(\Theta_i^{(1)},\Theta_i^{(2)}\right)\), assume that the pairs
  \((N_{it},M_{it})\), \(t\in\mathbb N\), are mutually independent with
  \[
  N_{it}\mid\Theta_i^{(1)},\Theta_i^{(2)}
  \sim
  \operatorname{Pois}\left(\lambda_{it}\Theta_i^{(1)}\right)
  \]
  and
  \[
  M_{it}\mid N_{it},\Theta_i^{(1)},\Theta_i^{(2)}
  \sim
  \operatorname{Bin}\left(N_{it},\Phi\left(\eta_{it}+\Theta_i^{(2)}\right)\right),
  \]
  Furthermore, the regions are independent: the sequences
  \(\left((N_{it},M_{it})\right)_{t\in\mathbb N}\), \(i=1,\ldots,I\),
  are mutually independent.
\end{itemize}
\end{model}
Here, \(\operatorname{Bin}(0,p)\) is the point mass at \(0\), so that
\(M_{it}=0\) whenever \(N_{it}=0\).

Since \(\E{\Theta_i^{(1)}}=1\) and \(\E{\Theta_i^{(2)}}=0\), the
intercepts in \(\alpha\) and \(\beta\) are identifiable, and
\(\lambda_{it}\) is the a priori expected total claim count. The random
effect \(\Theta_i^{(1)}\) captures the unobserved heterogeneity of the
claim frequency across regions, whereas \(\Theta_i^{(2)}\) captures the
unobserved heterogeneity of the propensity of a claim to be extreme.
For notational convenience, write
\[
p_{it}(\theta):=\Phi\left(\eta_{it}+\theta\right),
\qquad \theta\in\Real.
\]

\begin{remark}\label{rem.thin.exposure}
All results in this section use \(\lambda_{it}\) and \(\eta_{it}\) only
as known constants. Hence they remain valid for any
\(\lambda_{it}\in\Real_{>0}\) and \(\eta_{it}\in\Real\); in particular,
an exposure offset \(e_{it}>0\) can be included through
\(\lambda_{it}=e_{it}\exp\left(x_{it}^\top\alpha\right)\).
\end{remark}

The following lemma is the well-known thinning property of the Poisson
distribution. It shows that, conditional on the random effects, the
extreme and non-extreme claim counts are independent Poisson random
variables.

\begin{lemma}[Poisson thinning]
\label{lem.thin}
Under Model~\ref{model.thin}, conditional on
\(\left(\Theta_i^{(1)},\Theta_i^{(2)}\right)\), the random variables
\(M_{it}\) and \(N_{it}-M_{it}\) are independent with
\[
M_{it}\mid\Theta_i^{(1)},\Theta_i^{(2)}
\sim
\operatorname{Pois}\left(
\lambda_{it}\Theta_i^{(1)}p_{it}\left(\Theta_i^{(2)}\right)
\right)
\]
and
\[
N_{it}-M_{it}\mid\Theta_i^{(1)},\Theta_i^{(2)}
\sim
\operatorname{Pois}\left(
\lambda_{it}\Theta_i^{(1)}
\left\{1-p_{it}\left(\Theta_i^{(2)}\right)\right\}
\right).
\]
\end{lemma}

\begin{proof}
Fix \(\theta_1>0\) and \(\theta_2\in\Real\), and write
\(\Lambda:=\lambda_{it}\theta_1\) and \(p:=p_{it}(\theta_2)\). For
\(m,k\in\mathbb N_0\), Model~\ref{model.thin} gives
\[
\begin{aligned}
&\P{M_{it}=m,\,N_{it}-M_{it}=k\mid
\Theta_i^{(1)}=\theta_1,\Theta_i^{(2)}=\theta_2}\\
&\quad=
\P{N_{it}=m+k\mid\Theta_i^{(1)}=\theta_1}
\binom{m+k}{m}p^m(1-p)^k\\
&\quad=
\frac{e^{-\Lambda}\Lambda^{m+k}}{(m+k)!}
\frac{(m+k)!}{m!\,k!}p^m(1-p)^k\\
&\quad=
\frac{e^{-\Lambda p}(\Lambda p)^m}{m!}
\cdot
\frac{e^{-\Lambda(1-p)}\left\{\Lambda(1-p)\right\}^k}{k!},
\end{aligned}
\]
which proves the result.
\end{proof}

Since the regions are independent, it suffices to study a single
region. The following remark collects the notation for its observed
history.

\begin{remark}\label{rem.thin.notation}
Fix a region \(i\) and \(t\in\mathbb N\), and consider a realized
history
\[
\mathcal H_{it}:=
\left\{
(N_{ij},M_{ij})=(n_{ij},m_{ij}):\ j=1,\ldots,t
\right\},
\qquad
n_{ij},m_{ij}\in\mathbb N_0,
\quad
m_{ij}\le n_{ij}.
\]
Define the cumulative claim count and the cumulative a priori
frequency by
\[
n_{i\bullet t}:=\sum_{j=1}^t n_{ij},
\qquad
\lambda_{i\bullet t}:=\sum_{j=1}^t \lambda_{ij}.
\]

Each of the \(n_{i\bullet t}\) observed claims is regarded as a
Bernoulli trial, which is a success if the claim is extreme. To this
end, for \(j=1,\ldots,t\), define the sign vector
\[
u^{(j)}:=
\big(
\underbrace{1,\ldots,1}_{m_{ij}},
\underbrace{-1,\ldots,-1}_{n_{ij}-m_{ij}}
\big)^\top
\in\{-1,1\}^{n_{ij}},
\]
which is empty when \(n_{ij}=0\), and stack these vectors as
\[
u_{1:n_{i\bullet t}}:=
\left(
u^{(1)\top},\ldots,u^{(t)\top}
\right)^\top,
\qquad
h_{1:n_{i\bullet t}}:=
\big(
\underbrace{\eta_{i1},\ldots,\eta_{i1}}_{n_{i1}},
\ldots,
\underbrace{\eta_{it},\ldots,\eta_{it}}_{n_{it}}
\big)^\top .
\]
That is, \(u_k\) is the sign of the \(k\)th claim and \(h_k\) is the
value of \(\eta_{ij}\) in the period \(j\) of the \(k\)th claim. We then
define
\[
\mu_{it}:=
\left(u_1h_1,\ldots,u_{n_{i\bullet t}}h_{n_{i\bullet t}}\right)^\top
\in\Real^{n_{i\bullet t}},
\qquad
\boldsymbol U_{it}:=
u_{1:n_{i\bullet t}}u_{1:n_{i\bullet t}}^\top
\in\Real^{n_{i\bullet t}\times n_{i\bullet t}},
\]
which are the mean vector and the rank-one matrix appearing in
Lemma~\ref{lem.bp} for the trials formed by the observed claims.

For \(n,m\in\mathbb N_0\) with \(m\le n\), let
\(\mu_{it}^{(n,m)}\in\Real^{n_{i\bullet t}+n}\) and
\(\boldsymbol U_{it}^{(n,m)}\in\Real^{(n_{i\bullet t}+n)\times(n_{i\bullet t}+n)}\)
denote the corresponding quantities for the augmented history
\(\mathcal H_{it}\cup\{(N_{i,t+1},M_{i,t+1})=(n,m)\}\). That is,
\(n\) additional trials with \(\eta\)-value \(\eta_{i,t+1}\) are appended,
\(m\) of them with sign \(1\) and \(n-m\) of them with sign \(-1\).

Finally, define
\[
D_{it}(\theta):=
\prod_{j=1}^t
p_{ij}(\theta)^{m_{ij}}
\left\{1-p_{ij}(\theta)\right\}^{n_{ij}-m_{ij}},
\qquad\theta\in\Real.
\]
Since \(1-\Phi(a)=\Phi(-a)\) for \(a\in\Real\), we have
\begin{equation}\label{eq.thin.D}
D_{it}(\theta)
=
\prod_{k=1}^{n_{i\bullet t}}
\Phi\left(u_k(h_k+\theta)\right).
\end{equation}
For \(n_{i\bullet t}=0\), we adopt the conventions
\(D_{it}(\theta)=1\) and
\(\Phi_0\left(\mu_{it};\boldsymbol I_0+\sigma^2\boldsymbol U_{it}\right)=1\).
\end{remark}

The following theorem shows that the joint probability of the observed
history factorizes into a closed-form Poisson-gamma part, which depends
only on \((\alpha,r)\), and a Bernoulli-probit part of the form in
Lemma~\ref{lem.bp}, which depends only on \((\beta,\sigma)\).

\begin{theorem}[Likelihood factorization]
\label{thm.thin.lik}
Consider Model~\ref{model.thin} with the realized history
\(\mathcal H_{it}\), and let the notation be as in
Remark~\ref{rem.thin.notation}. Then
\begin{equation}\label{eq.thin.factorization}
\P{\mathcal H_{it}}
=
L^{(1)}_{it}\,L^{(2)}_{it},
\end{equation}
where
\begin{equation}\label{eq.thin.L1}
L^{(1)}_{it}
:=
\P{N_{i,1:t}=n_{i,1:t}}
=
\frac{\Gamma\left(r+n_{i\bullet t}\right)}
{\Gamma(r)\prod_{j=1}^t n_{ij}!}
\,
\frac{r^r\prod_{j=1}^t\lambda_{ij}^{n_{ij}}}
{\left(r+\lambda_{i\bullet t}\right)^{r+n_{i\bullet t}}}
\end{equation}
and
\begin{equation}\label{eq.thin.L2}
\begin{aligned}
L^{(2)}_{it}
&:=
\P{M_{i,1:t}=m_{i,1:t}\mid N_{i,1:t}=n_{i,1:t}}\\
&=
\prod_{j=1}^t\binom{n_{ij}}{m_{ij}}
\int_{\Real}
\frac{1}{\sigma}\phi\left(\frac{\theta}{\sigma}\right)
D_{it}(\theta)
\,{\rm d}\theta\\
&=
\prod_{j=1}^t\binom{n_{ij}}{m_{ij}}
\,
\Phi_{n_{i\bullet t}}
\left(
\mu_{it};
\boldsymbol I_{n_{i\bullet t}}
+
\sigma^2\boldsymbol U_{it}
\right).
\end{aligned}
\end{equation}
\end{theorem}

\begin{proof}
By the conditional independence in Model~\ref{model.thin},
\[
\P{\mathcal H_{it}}
=
\E{
\prod_{j=1}^t
\P{N_{ij}=n_{ij}\mid\Theta_i^{(1)}}
\P{M_{ij}=m_{ij}\mid N_{ij}=n_{ij},\Theta_i^{(2)}}
}.
\]
The first factor in the product depends only on \(\Theta_i^{(1)}\) and
the second only on \(\Theta_i^{(2)}\). Since \(\Theta_i^{(1)}\) and
\(\Theta_i^{(2)}\) are independent, we obtain
\[
\P{\mathcal H_{it}}
=
\E{
\prod_{j=1}^t
\P{N_{ij}=n_{ij}\mid\Theta_i^{(1)}}
}
\,
\E{
\prod_{j=1}^t
\P{M_{ij}=m_{ij}\mid N_{ij}=n_{ij},\Theta_i^{(2)}}
}.
\]

For the first expectation, the gamma integral gives
\[
\begin{aligned}
\E{
\prod_{j=1}^t
\P{N_{ij}=n_{ij}\mid\Theta_i^{(1)}}
}
&=
\int_0^\infty
\prod_{j=1}^t
\frac{e^{-\lambda_{ij}\theta}(\lambda_{ij}\theta)^{n_{ij}}}{n_{ij}!}
\,
\frac{r^r}{\Gamma(r)}\theta^{r-1}e^{-r\theta}
\,{\rm d}\theta\\
&=
\frac{r^r\prod_{j=1}^t\lambda_{ij}^{n_{ij}}}
{\Gamma(r)\prod_{j=1}^t n_{ij}!}
\int_0^\infty
\theta^{r+n_{i\bullet t}-1}
e^{-(r+\lambda_{i\bullet t})\theta}
\,{\rm d}\theta,
\end{aligned}
\]
which equals the right-hand side of \eqref{eq.thin.L1}. Summing the
first display over all \(m_{i,1:t}\) with \(0\le m_{ij}\le n_{ij}\)
shows that this expectation is \(\P{N_{i,1:t}=n_{i,1:t}}\), because
the second expectation sums to one.

For the second expectation, we have
\[
\E{
\prod_{j=1}^t
\P{M_{ij}=m_{ij}\mid N_{ij}=n_{ij},\Theta_i^{(2)}}
}
=
\prod_{j=1}^t\binom{n_{ij}}{m_{ij}}
\E{D_{it}\left(\Theta_i^{(2)}\right)},
\]
which proves the second line of \eqref{eq.thin.L2}. It then follows
from \eqref{eq.thin.factorization} and \eqref{eq.thin.L1} that this
expression is \(\P{M_{i,1:t}=m_{i,1:t}\mid N_{i,1:t}=n_{i,1:t}}\).
Finally, by \eqref{eq.thin.D},
\[
\E{D_{it}\left(\Theta_i^{(2)}\right)}
=
\E{
\prod_{k=1}^{n_{i\bullet t}}
\Phi\left(u_k\left(h_k+\Theta_i^{(2)}\right)\right)
},
\]
and applying Lemma~\ref{lem.bp} with \(k=n_{i\bullet t}\),
\(u_{1:k}\), and \(h_{1:k}\) proves the last line of
\eqref{eq.thin.L2}.
\end{proof}

\begin{remark}[Separable estimation]
\label{rem.thin.estimation}
By Theorem~\ref{thm.thin.lik} and the independence of the regions,
the log-likelihood of the observations
\(\mathcal H_{1T},\ldots,\mathcal H_{IT}\) decomposes as
\[
\ell(\alpha,r,\beta,\sigma)
=
\ell^{(1)}(\alpha,r)+\ell^{(2)}(\beta,\sigma),
\qquad
\ell^{(k)}:=\sum_{i=1}^I\log L^{(k)}_{iT},
\quad k=1,2.
\]
Hence the maximum likelihood estimators of \((\alpha,r)\) and
\((\beta,\sigma)\) are obtained from two separate optimizations, and
the Fisher information matrix is block-diagonal. Here,
\(\ell^{(1)}\) is the log-likelihood of a negative binomial
random effect panel model for \(N_{it}\), and \(\ell^{(2)}\) is the
log-likelihood of a binomial probit model for \(M_{it}\) given
\(N_{it}\) with a Gaussian random intercept. Periods with \(n_{it}=0\)
do not contribute to \(\ell^{(2)}\).
\end{remark}

\begin{remark}[Computation]
\label{rem.thin.computation}
The dimension \(n_{i\bullet t}\) of the multivariate normal
distribution function in \eqref{eq.thin.L2} equals the cumulative
number of claims of region \(i\), which can be large. However, since
\(\boldsymbol U_{it}\) is of rank one, the second line of
\eqref{eq.thin.L2} expresses the same quantity as the one-dimensional
integral
\[
\Phi_{n_{i\bullet t}}
\left(
\mu_{it};
\boldsymbol I_{n_{i\bullet t}}+\sigma^2\boldsymbol U_{it}
\right)
=
\int_{\Real}\phi(z)\,D_{it}(\sigma z)\,{\rm d}z,
\]
which depends on the data only through the \(t\) pairs
\((n_{ij},m_{ij})\) and can be evaluated accurately by Gauss--Hermite
quadrature.
\end{remark}

The following corollary shows that the two random effects remain
independent a posteriori, and that the frequency random effect has a
conjugate gamma posterior.

\begin{corollary}[Posterior distribution of the random effects]
\label{cor.thin.posterior}
Under the setting of Theorem~\ref{thm.thin.lik}, the following
results hold.
\begin{itemize}
  \item[i.] Conditional on \(\mathcal H_{it}\), the random effects
  \(\Theta_i^{(1)}\) and \(\Theta_i^{(2)}\) are independent.

  \item[ii.] We have
  \[
  \Theta_i^{(1)}\mid\mathcal H_{it}
  \sim
  \operatorname{Gamma}\left(r+n_{i\bullet t},\,r+\lambda_{i\bullet t}\right).
  \]

  \item[iii.] The conditional distribution
  \(\Theta_i^{(2)}\mid\mathcal H_{it}\) has density
  \(\pi(\theta\mid\mathcal H_{it})\) at \(\theta\in\Real\) given by
  \begin{equation}\label{eq.thin.posterior2}
  \pi(\theta\mid\mathcal H_{it})
  =
  \frac{
  \displaystyle
  \frac{1}{\sigma}\phi\left(\frac{\theta}{\sigma}\right)
  \prod_{k=1}^{n_{i\bullet t}}
  \Phi\left(u_k(h_k+\theta)\right)
  }{
  \Phi_{n_{i\bullet t}}
  \left(
  \mu_{it};
  \boldsymbol I_{n_{i\bullet t}}+\sigma^2\boldsymbol U_{it}
  \right)
  }.
  \end{equation}
\end{itemize}
\end{corollary}

\begin{proof}
Let \(\pi_1\) and \(\pi_2\) denote the prior densities of
\(\Theta_i^{(1)}\) and \(\Theta_i^{(2)}\), respectively. By Bayes'
formula and the proof of Theorem~\ref{thm.thin.lik}, the joint
posterior density of
\(\left(\Theta_i^{(1)},\Theta_i^{(2)}\right)\) at
\((\theta_1,\theta_2)\in\Real_{>0}\times\Real\) is
\[
\frac{
\pi_1(\theta_1)
\prod_{j=1}^t\P{N_{ij}=n_{ij}\mid\Theta_i^{(1)}=\theta_1}
}{L^{(1)}_{it}}
\cdot
\frac{
\pi_2(\theta_2)
\prod_{j=1}^t\binom{n_{ij}}{m_{ij}}D_{it}(\theta_2)
}{L^{(2)}_{it}}.
\]
Each factor integrates to one, which proves part i. The first factor
is proportional to
\[
\theta_1^{r+n_{i\bullet t}-1}e^{-(r+\lambda_{i\bullet t})\theta_1},
\]
which proves part ii. Substituting \eqref{eq.thin.D} and
\eqref{eq.thin.L2} into the second factor proves part iii.
\end{proof}

By part ii of Corollary~\ref{cor.thin.posterior}, the posterior of the
frequency random effect depends on the history only through
\(n_{i\bullet t}\); in particular, the extreme claim counts
\(m_{i,1:t}\) carry no information on \(\Theta_i^{(1)}\). Conversely,
by part iii, the total claim counts affect the posterior of
\(\Theta_i^{(2)}\) only by determining the number of Bernoulli trials.
The density \eqref{eq.thin.posterior2} is a Gaussian density
multiplied by a product of probit functions, normalized by the
multivariate normal distribution function in Lemma~\ref{lem.bp}.

We now present the exact Bayesian one-step-ahead predictions.

\begin{theorem}[Closed-form one-step prediction]
\label{thm.thin.pred}
Consider Model~\ref{model.thin} with the realized history
\(\mathcal H_{it}\), and let the notation be as in
Remark~\ref{rem.thin.notation}. Then, we have the following results.
\begin{itemize}
  \item[i.] For \(n,m\in\mathbb N_0\) with \(m\le n\), the predictive
  probability is
  \[
  \begin{aligned}
  &\P{N_{i,t+1}=n,\,M_{i,t+1}=m\mid\mathcal H_{it}}\\
  &\quad=
  \frac{\Gamma\left(r+n_{i\bullet t}+n\right)}
  {\Gamma\left(r+n_{i\bullet t}\right)n!}
  \left(
  \frac{r+\lambda_{i\bullet t}}{r+\lambda_{i\bullet t}+\lambda_{i,t+1}}
  \right)^{r+n_{i\bullet t}}
  \left(
  \frac{\lambda_{i,t+1}}{r+\lambda_{i\bullet t}+\lambda_{i,t+1}}
  \right)^{n}\\
  &\qquad\times
  \binom{n}{m}
  \frac{
  \Phi_{n_{i\bullet t}+n}
  \left(
  \mu_{it}^{(n,m)};
  \boldsymbol I_{n_{i\bullet t}+n}+\sigma^2\boldsymbol U_{it}^{(n,m)}
  \right)
  }{
  \Phi_{n_{i\bullet t}}
  \left(
  \mu_{it};
  \boldsymbol I_{n_{i\bullet t}}+\sigma^2\boldsymbol U_{it}
  \right)
  }.
  \end{aligned}
  \]

  \item[ii.] The one-step-ahead prediction of the total claim count is
  \[
  \E{N_{i,t+1}\mid\mathcal H_{it}}
  =
  \lambda_{i,t+1}\,
  \frac{r+n_{i\bullet t}}{r+\lambda_{i\bullet t}}
  =
  \lambda_{i,t+1}
  \left\{
  \zeta_{it}\frac{n_{i\bullet t}}{\lambda_{i\bullet t}}
  +
  \left(1-\zeta_{it}\right)
  \right\},
  \qquad
  \zeta_{it}:=\frac{\lambda_{i\bullet t}}{r+\lambda_{i\bullet t}}.
  \]

  \item[iii.] The predictive probability that a claim in period
  \(t+1\) is extreme is
  \[
  \widehat p_{i,t+1}
  :=
  \E{p_{i,t+1}\left(\Theta_i^{(2)}\right)\mid\mathcal H_{it}}
  =
  \frac{
  \Phi_{n_{i\bullet t}+1}
  \left(
  \mu_{it}^{(1,1)};
  \boldsymbol I_{n_{i\bullet t}+1}+\sigma^2\boldsymbol U_{it}^{(1,1)}
  \right)
  }{
  \Phi_{n_{i\bullet t}}
  \left(
  \mu_{it};
  \boldsymbol I_{n_{i\bullet t}}+\sigma^2\boldsymbol U_{it}
  \right)
  }.
  \]

  \item[iv.] For \(n\in\mathbb N_0\), the one-step-ahead prediction of
  the extreme claim count given the total claim count is
  \[
  \E{M_{i,t+1}\mid\mathcal H_{it},N_{i,t+1}=n}
  =
  n\,\widehat p_{i,t+1}.
  \]

  \item[v.] The one-step-ahead prediction of the extreme claim count is
  \[
  \E{M_{i,t+1}\mid\mathcal H_{it}}
  =
  \lambda_{i,t+1}\,
  \frac{r+n_{i\bullet t}}{r+\lambda_{i\bullet t}}\,
  \widehat p_{i,t+1}.
  \]
\end{itemize}
\end{theorem}

\begin{proof}
We first show part i. Applying Theorem~\ref{thm.thin.lik} to the
augmented history
\(\mathcal H_{it}\cup\{(N_{i,t+1},M_{i,t+1})=(n,m)\}\) and to
\(\mathcal H_{it}\), we obtain
\[
\P{N_{i,t+1}=n,\,M_{i,t+1}=m\mid\mathcal H_{it}}
=
\frac{L^{(1)}_{i,t+1}}{L^{(1)}_{it}}
\cdot
\frac{L^{(2)}_{i,t+1}}{L^{(2)}_{it}},
\]
where \(L^{(1)}_{i,t+1}\) and \(L^{(2)}_{i,t+1}\) are evaluated at the
augmented history. By \eqref{eq.thin.L1},
\[
\frac{L^{(1)}_{i,t+1}}{L^{(1)}_{it}}
=
\frac{\Gamma\left(r+n_{i\bullet t}+n\right)}
{\Gamma\left(r+n_{i\bullet t}\right)n!}
\,
\frac{\lambda_{i,t+1}^{n}\left(r+\lambda_{i\bullet t}\right)^{r+n_{i\bullet t}}}
{\left(r+\lambda_{i\bullet t}+\lambda_{i,t+1}\right)^{r+n_{i\bullet t}+n}},
\]
and, by the last line of \eqref{eq.thin.L2},
\[
\frac{L^{(2)}_{i,t+1}}{L^{(2)}_{it}}
=
\binom{n}{m}
\frac{
\Phi_{n_{i\bullet t}+n}
\left(
\mu_{it}^{(n,m)};
\boldsymbol I_{n_{i\bullet t}+n}+\sigma^2\boldsymbol U_{it}^{(n,m)}
\right)
}{
\Phi_{n_{i\bullet t}}
\left(
\mu_{it};
\boldsymbol I_{n_{i\bullet t}}+\sigma^2\boldsymbol U_{it}
\right)
},
\]
which proves part i.

To show part ii, note that, conditional on
\(\left(\Theta_i^{(1)},\Theta_i^{(2)}\right)\), \(N_{i,t+1}\) is
independent of \(\mathcal H_{it}\). Hence, by part ii of
Corollary~\ref{cor.thin.posterior},
\[
\E{N_{i,t+1}\mid\mathcal H_{it}}
=
\lambda_{i,t+1}\E{\Theta_i^{(1)}\mid\mathcal H_{it}}
=
\lambda_{i,t+1}\frac{r+n_{i\bullet t}}{r+\lambda_{i\bullet t}},
\]
and the second expression follows from
\(
(r+n_{i\bullet t})/(r+\lambda_{i\bullet t})
=
\zeta_{it}\,n_{i\bullet t}/\lambda_{i\bullet t}+1-\zeta_{it}.
\)

To show part iii, part iii of Corollary~\ref{cor.thin.posterior} gives
\[
\widehat p_{i,t+1}
=
\frac{
\E{
p_{i,t+1}\left(\Theta_i^{(2)}\right)
\prod_{k=1}^{n_{i\bullet t}}
\Phi\left(u_k\left(h_k+\Theta_i^{(2)}\right)\right)
}
}{
\Phi_{n_{i\bullet t}}
\left(
\mu_{it};
\boldsymbol I_{n_{i\bullet t}}+\sigma^2\boldsymbol U_{it}
\right)
}.
\]
Since
\(p_{i,t+1}(\theta)=\Phi\left(1\cdot(\eta_{i,t+1}+\theta)\right)\), the
product in the numerator is the product in \eqref{eq.thin.D} for the
history augmented by \((n,m)=(1,1)\). Applying Lemma~\ref{lem.bp} to the
augmented history proves part iii.

To show part iv, we first observe that, conditional on
\(\mathcal H_{it}\), the joint density of
\(\left(\Theta_i^{(1)},\Theta_i^{(2)}\right)\) and the probability mass
function of \(N_{i,t+1}\) satisfy
\[
\pi(\theta_1,\theta_2,n\mid\mathcal H_{it})
=
\pi(\theta_1\mid\mathcal H_{it})
\P{N_{i,t+1}=n\mid\Theta_i^{(1)}=\theta_1}
\,
\pi(\theta_2\mid\mathcal H_{it}),
\]
by part i of Corollary~\ref{cor.thin.posterior} and
Model~\ref{model.thin}. Hence \(\Theta_i^{(2)}\) is independent of
\(N_{i,t+1}\) conditional on \(\mathcal H_{it}\). Together with the
binomial observation equation, this gives
\[
\E{M_{i,t+1}\mid\mathcal H_{it},N_{i,t+1}=n}
=
n\,\E{p_{i,t+1}\left(\Theta_i^{(2)}\right)\mid\mathcal H_{it},N_{i,t+1}=n}
=
n\,\widehat p_{i,t+1}.
\]

Finally, part v follows from part iv and part ii, since
\[
\E{M_{i,t+1}\mid\mathcal H_{it}}
=
\E{N_{i,t+1}\mid\mathcal H_{it}}\,\widehat p_{i,t+1}.
\]
\end{proof}

\begin{remark}[Interpretation of the predictions]
\label{rem.thin.interpretation}
Part ii of Theorem~\ref{thm.thin.pred} is the exact credibility formula
of the Poisson-gamma model: the a priori frequency \(\lambda_{i,t+1}\)
is multiplied by a credibility-weighted average of the observed ratio
\(n_{i\bullet t}/\lambda_{i\bullet t}\) and the prior mean \(1\), with
credibility factor \(\zeta_{it}\), which increases to one as
\(\lambda_{i\bullet t}\to\infty\). Part iii is the posterior predictive
probability that an additional Bernoulli-probit trial, namely a new
claim, is a success.
In particular, if \(n_{i\bullet t}=0\), then
\(
\widehat p_{i,t+1}
=
\Phi\left(\eta_{i,t+1}/\sqrt{1+\sigma^2}\right),
\)
which is the a priori probability that a claim is extreme. Part v shows
that the prediction of the extreme claim count is the product of the
two exact Bayesian predictions, one for the frequency and one for the
proportion of extreme claims, as a consequence of the a posteriori
independence of the random effects.
\end{remark}

We conclude this section with the unconditional moments of
Model~\ref{model.thin}, which describe the dependence induced by the
random effects.

\begin{proposition}[Moments]
\label{prop.thin.moments}
Under Model~\ref{model.thin}, fix a region \(i\), and define
\[
\pi_{it}:=\Phi\left(\frac{\eta_{it}}{\sqrt{1+\sigma^2}}\right),
\qquad
\pi_{i,st}:=
\Phi_2\left(
(\eta_{is},\eta_{it})^\top;
\boldsymbol I_2+\sigma^2\boldsymbol J_2
\right),
\qquad s,t\in\mathbb N,
\]
where \(\boldsymbol J_2\) is the \(2\times2\) matrix of ones. Then, for
\(s,t\in\mathbb N\) with \(s\neq t\), we have
\[
\begin{aligned}
\E{N_{it}}&=\lambda_{it},
&
\E{M_{it}}&=\lambda_{it}\pi_{it},\\
\Var{N_{it}}&=\lambda_{it}+\frac{\lambda_{it}^2}{r},
&
\cov{N_{is},N_{it}}&=\frac{\lambda_{is}\lambda_{it}}{r},\\
\cov{N_{it},M_{it}}&=\pi_{it}\left(\lambda_{it}+\frac{\lambda_{it}^2}{r}\right),
&
\cov{N_{is},M_{it}}&=\frac{\lambda_{is}\lambda_{it}\pi_{it}}{r},
\end{aligned}
\]
and
\[
\begin{aligned}
\Var{M_{it}}
&=
\lambda_{it}\pi_{it}
+
\lambda_{it}^2
\left\{
\left(1+\frac1r\right)\pi_{i,tt}-\pi_{it}^2
\right\},\\
\cov{M_{is},M_{it}}
&=
\lambda_{is}\lambda_{it}
\left\{
\left(1+\frac1r\right)\pi_{i,st}-\pi_{is}\pi_{it}
\right\}.
\end{aligned}
\]
\end{proposition}

\begin{proof}
Write \(\Theta_1:=\Theta_i^{(1)}\) and \(\Theta_2:=\Theta_i^{(2)}\), and
recall \(\E{\Theta_1}=1\) and \(\E{\Theta_1^2}=1+1/r\). By
Lemma~\ref{lem.bp} with \(k=1\) and \(k=2\), and \(u_{1:k}\) the
vector of ones,
\[
\E{p_{it}(\Theta_2)}=\pi_{it},
\qquad
\E{p_{is}(\Theta_2)p_{it}(\Theta_2)}=\pi_{i,st},
\qquad s,t\in\mathbb N.
\]
Conditional on \((\Theta_1,\Theta_2)\), Model~\ref{model.thin} and
Lemma~\ref{lem.thin} give, for \(s\neq t\),
\[
\begin{aligned}
\E{N_{it}\mid\Theta_1,\Theta_2}&=\lambda_{it}\Theta_1,
&
\E{M_{it}\mid\Theta_1,\Theta_2}&=\lambda_{it}\Theta_1p_{it}(\Theta_2),\\
\E{N_{it}^2\mid\Theta_1,\Theta_2}&=\lambda_{it}\Theta_1+\lambda_{it}^2\Theta_1^2,
&
\E{M_{it}^2\mid\Theta_1,\Theta_2}&=\lambda_{it}\Theta_1p_{it}(\Theta_2)
+\lambda_{it}^2\Theta_1^2p_{it}(\Theta_2)^2,
\end{aligned}
\]
and
\[
\begin{aligned}
\E{N_{it}M_{it}\mid\Theta_1,\Theta_2}
&=
p_{it}(\Theta_2)\E{N_{it}^2\mid\Theta_1,\Theta_2}
=
p_{it}(\Theta_2)\left(\lambda_{it}\Theta_1+\lambda_{it}^2\Theta_1^2\right),\\
\E{N_{is}M_{it}\mid\Theta_1,\Theta_2}
&=
\lambda_{is}\lambda_{it}\Theta_1^2p_{it}(\Theta_2),\\
\E{M_{is}M_{it}\mid\Theta_1,\Theta_2}
&=
\lambda_{is}\lambda_{it}\Theta_1^2p_{is}(\Theta_2)p_{it}(\Theta_2).
\end{aligned}
\]
Taking expectations, using the independence of \(\Theta_1\) and
\(\Theta_2\), and subtracting the products of the means proves the
result.
\end{proof}


\section*{Acknowledgements}

Jae Youn Ahn was supported by the National Research Foundation of Korea (NRF) through a grant funded by the Korea government (MSIT) (RS-2025-23524530) and the Basic Science Research Program (Priority Research Institute) funded by the Ministry of Education (2021R1A6A1A10039823).
Parts of this manuscript were written while Mario V.~W\"uthrich was hosted at Ewha Womans University, Seoul.


\bibliographystyle{apalike}
\bibliography{earthquake_bib}





\appendix

\section{Miscelleaneous results}\label{app.1}


\end{document}

